\documentclass[11pt,oneside]{book}

% ---- Trim size: 6in x 9in trade paperback ----
\usepackage[paperwidth=6in,paperheight=9in,margin=0.75in,bottom=0.85in,top=0.85in]{geometry}

% ---- Font handling.
%
%      NOTE: this used to take the classic NFSS route --- \usepackage[T1]
%      {fontenc} plus \usepackage{tgtermes} --- on the theory that it kept
%      the body text off luaotfload. It did not work. Loading fontspec
%      afterwards (which this document must, for EB Garamond and the runic
%      staves) switches the default encoding to TU, at which point the T1
%      family "qtm" is no longer resolvable and LaTeX silently substitutes
%      Latin Modern for the entire body of the book. The build throws no
%      error; it just quietly sets 60,000 words in the wrong face.
%
%      Fixed by declaring the body font through fontspec as well, which is
%      the only consistent option once fontspec is in the document at all.
%      Ligatures and old-style figures are off by design --- this is a
%      Times-class text face and should behave like one. ----
\usepackage{fontspec}
\setmainfont{TeX Gyre Termes}[
  Ligatures = TeX,
  Numbers   = Lining
]

% ---- EB Garamond, reserved for the title page and chapter-opening
%      ornament only (this matches the font the original manuscript is set
%      in throughout — see its own colophon, "This book was set in EB
%      Garamond" — while the revised edition's running body keeps TeX Gyre
%      Termes for series consistency with the rest of the Formaetrix
%      line) ----
\newfontfamily\ebgaramond{EBGaramond-Regular}[
  Path = decl-fonts/,
  Extension = .ttf,
  UprightFont = EBGaramond-Regular,
  BoldFont = EBGaramond-Bold,
  ItalicFont = EBGaramond-Italic,
  BoldItalicFont = EBGaramond-BoldItalic
]

\usepackage{microtype}
\usepackage[normalem]{ulem}
\usepackage{titlesec}
\usepackage{titletoc}
\usepackage{fancyhdr}
\usepackage{setspace}
\usepackage[english]{babel}
\usepackage{hyperref}
\hypersetup{
  hidelinks,
  pdftitle={Declensions of Dark Water},
  pdfauthor={Elian Voigt}
}

% ---- THE NINETEEN STAVES OF VIKAMÁL
%
%      Chapters carry no titles and no Latin numerals. Each is marked
%      with a single stave, and the staves are not ordinals. Each one
%      names a case, and the book declines itself through them in order:
%
%        I    ᚠ  Nefn         nominative     that which names reality
%        II   ᚢ  Eign         genitive       ownership, inheritance
%        III  ᚦ  Gjǫf         dative         what is given or received
%        IV   ᚨ  Leið         accusative     directed action
%        V    ᚱ  Staðr        locative       where memory remains
%        VI   ᚲ  Brott        ablative       separation
%        VII  ᚷ  Kall         vocative       calling across worlds
%        VIII ᚹ  Verk         instrumental   means rather than intent
%        IX   ᚺ  Fylgd        comitative     those who travel together
%        X    ᚾ  Orsök        causal         why something exists
%        XI   ᛁ  Tíð          temporal       time as grammar
%        XII  ᛃ  Skipti       transformative becoming something else
%        XIII ᛇ  Spegill      reflexive      self referring to self
%        XIV  ᛈ  Þögn         privative      meaning through absence
%        XV   ᛉ  Vottur       evidential     witness rather than truth
%        XVI  ᛊ  Skuggi       conditional    what almost happened
%        XVII ᛏ  Þröskuldur   terminative    crossing a boundary
%        XVIII ᛒ Lifandi      living case    grammar still changing
%        XIX  ᛖ  Ónefnt       null case      that which cannot be spoken
%
%      The book runs to seventeen chapters and a coda, which takes
%      XVIII, Lifandi — the case that is still changing. XIX, Ónefnt,
%      the null case, is never printed anywhere in this volume. That is
%      deliberate. A reader who knows the sequence will find the last
%      stave missing, which is the only way a book can set a case that
%      cannot be spoken.
%
%      These are Elder Fuþark forms, not the Younger Fuþark of the
%      Icelandic calendar staves — a deliberate anachronism. Vikamál is
%      older than the settlement and older than the language it now
%      hides inside; the alphabet says so before the prose does.
%
%      \thechapter stays Roman so hyperref, bookmarks and any future
%      cross-references keep working; only the printed display is
%      runic. ----
\renewcommand{\thechapter}{\Roman{chapter}}

% ---- Runic face. Junicode is the recommended choice: it is free, it is
%      cut for medievalists, its proportions sit comfortably beside EB
%      Garamond, and it covers the full Runic block including the three
%      golden-number staves. Drop Junicode-Regular.ttf into decl-fonts/.
%      If it is absent the book still compiles and falls back to Roman
%      numerals, so a missing font degrades the design without breaking
%      the build. ----
\newif\ifrunicavailable
\IfFontExistsTF{decl-fonts/Junicode-Regular.ttf}{
  \newfontfamily\runicface{Junicode-Regular}[
    Path = decl-fonts/,
    Extension = .ttf
  ]
  \runicavailabletrue
}{
  \newcommand{\runicface}{}
  \runicavailablefalse
}

% The staves, in case order. Position 19 (ᛖ Ónefnt, the null case) is
% defined here for completeness and is never called by the document.
\newcommand{\runicnumeral}[1]{%
  \ifrunicavailable
    {\runicface\ifcase#1\or ᚠ\or ᚢ\or ᚦ\or ᚨ\or ᚱ\or ᚲ\or ᚷ\or ᚹ\or ᚺ\or
     ᚾ\or ᛁ\or ᛃ\or ᛇ\or ᛈ\or ᛉ\or ᛊ\or ᛏ\or ᛒ\or ᛖ\fi}%
  \else
    % Fallback keys off the argument, not the chapter counter, so the
    % coda's stave does not silently print the previous chapter's number.
    \MakeUppercase{\romannumeral#1}%
  \fi}

% The coda is not a \chapter, so its stave is set explicitly. It takes
% XVIII, Lifandi — the living case, grammar still changing.
\newcommand{\codastave}{\runicnumeral{18}}

% ---- Paragraph style: novel-standard, first-line indent, no extra vertical gap ----
\setlength{\parindent}{1.15em}
\setlength{\parskip}{0pt}
\linespread{1.06}

% ---- Scene break ornament ----
\newcommand{\sceneline}{%
  \par\vspace{1.2em}
  \begin{center}
    \fontsize{11}{11}\selectfont \textbullet\quad\textbullet\quad\textbullet
  \end{center}
  \vspace{0.6em}\noindent\ignorespaces
}

% ---- Field apparatus: Rowan's notebook. Set in a narrower measure, one
%      point down, to read as a hand separate from the narrating voice.
%      Restrained by design — the novel is prose first and a documentary
%      object only where the documentation is itself dramatic. ----
\newenvironment{fieldnote}
  {\par\vspace{0.9em}\begingroup
   \leftskip=1.6em \rightskip=1.6em
   \fontsize{10}{12.4}\selectfont\itshape
   \setlength{\parindent}{0pt}\noindent\ignorespaces}
  {\par\endgroup\vspace{0.9em}\noindent\ignorespaces}

% ---- A glossed form: the Vikamál, then its morphological reading. Used
%      sparingly, where Rowan is actually working rather than merely
%      noticing. ----
\newcommand{\gloss}[2]{\emph{#1}\,\textsuperscript{\textup{\tiny#2}}}

% ---- A paradigm: the small declension tables Rowan draws when a
%      hypothesis is still standing. Deliberately plain — he rules them
%      by hand in a notebook, not in a typesetting system. ----
\newenvironment{paradigm}
  {\par\vspace{0.9em}\begingroup\centering
   \fontsize{9.5}{12}\selectfont
   \renewcommand{\arraystretch}{1.15}}
  {\par\endgroup\vspace{0.9em}\noindent\ignorespaces}

% ---- A struck hypothesis: Rowan crossing out a theory he has just lost
%      confidence in, without erasing the evidence that he held it. ----
\newcommand{\struck}[1]{\sout{#1}}

% ---- Chapter heading style: "Chapter N", the Aldine-leaf ornament (as
%      called out in the original's own colophon), then the subtitle
%      where one exists. Untitled chapters simply omit the subtitle line. ----
% The chapter opening is now a single cut mark and the Aldine leaf. No
% word "Chapter", no numeral in Latin script, no title. The stave is set
% large because it is meant to read as something carved rather than
% something typeset.
\titleformat{\chapter}[display]
  {\normalfont\centering}
  {\fontsize{26}{30}\selectfont\runicnumeral{\value{chapter}}\\[0.7em]%
   {\fontsize{11}{13}\selectfont\ebgaramond ❦}}
  {10pt}
  {}
\titlespacing*{\chapter}{0pt}{2.2\baselineskip}{2.2\baselineskip}

\pagestyle{fancy}
\fancyhf{}
\fancyhead[LE]{\small\scshape Declensions of Dark Water}
\fancyhead[RO]{\small\scshape Elian Voigt}
\fancyfoot[C]{\thepage}
\renewcommand{\headrulewidth}{0.3pt}

\fancypagestyle{plain}{%
  \fancyhf{}
  \fancyfoot[C]{\thepage}
  \renewcommand{\headrulewidth}{0pt}
}

\begin{document}

\frontmatter
\pagestyle{empty}

% ---------------- TITLE PAGE ----------------
\begin{titlepage}
\centering
\ebgaramond
\vspace*{2.3in}
{\fontsize{24}{28}\selectfont\addfontfeature{Letters=SmallCaps}Declensions of Dark Water\par}
\vspace{0.5in}
{\Large ❦\par}
\vspace{0.5in}
{\large\itshape a novel\par}
\vspace{1.1in}
{\Large Elian Voigt\par}
\vfill
{\small\addfontfeature{Letters=SmallCaps}Form\ae trix\par}
\vspace*{0.6in}
\end{titlepage}

% ---------------- COPYRIGHT PAGE ----------------
\newpage
\thispagestyle{empty}
\vspace*{1in}
\small
\noindent Copyright \textcopyright\ 2026 by Elian Voigt

\vspace{1em}
\noindent All rights reserved. No part of this book may be reproduced, stored in a retrieval system, or transmitted in any form or by any means, electronic, mechanical, photocopying, recording, or otherwise, without prior written permission of the publisher, except brief quotations used in reviews and critical articles.

\vspace{1em}
\noindent This is a work of fiction. Names, characters, places, and incidents either are products of the author's imagination or are used fictitiously. Any resemblance to actual persons, living or dead, events, or locales is entirely coincidental.

\vspace{1em}
\noindent Cover design and interior typography by Form\ae trix.\\
Set in EB Garamond.\\
Second Edition, Revised, 2026.

\vspace{1em}
\noindent ISBN 978-1-XXXXXX-XX-X

\vspace{1em}
\noindent formaetrix.com

% ---------------- DEDICATION ----------------
\newpage
\thispagestyle{empty}
\vspace*{2.6in}
\begin{center}
\itshape
For those who keep the grammar.
\end{center}

\newpage
\thispagestyle{empty}
\mbox{}

\mainmatter
\pagestyle{fancy}

\chapter{}

He empties his mother's flat the way he would open a site.

This is not a decision. He arrives on the first Saturday in February with
a roll of bin bags, a marker pen, and eighteen flat-packed archive boxes
bought in bulk from a supplier the department uses, and by the time he
has the kettle on he has already established a working order --- back
bedroom first, kitchen last, nothing to leave the building unlabelled ---
and it is only when he catches himself writing \emph{BEDROOM 2 / TEXTILE
/ NO PROVENANCE} on the side of a box of his mother's cardigans that he
stops with the pen in his hand and stands very still in the middle of
the room for about a minute.

Then he finishes writing it, because the alternative is a box with
nothing on it.

The flat is the ground floor of a converted house on the Grantchester
road, and it is small and warm and profoundly, aggressively tidy, and it
takes him four weekends to establish that it contains nothing. Not
nothing of value --- there are the cardigans, and a good set of pans, and
eleven hundred pounds in premium bonds, and a chest of drawers that a
dealer in Cherry Hinton will eventually give him three hundred for. He
means nothing of \emph{her}. No letters. No address book older than the
current one. No diaries, no keepsakes, no drawer of the unfileable that
every human being on earth maintains. He finds a shoebox at the top of
the wardrobe and his heart does something complicated and it is full of
receipts for a boiler, sorted by year.

The photographs begin at roughly twenty-five. He knows this already;
he has known it since he was a boy. But it is different to know it and
to lay it out on a bed. There is a whole life, orderly and well
documented: Elizabeth at thirty at a colleague's wedding, Elizabeth at
forty holding him at a swimming gala with the pinched proud expression
of a woman being photographed against her preference, Elizabeth at
sixty in a garden in Suffolk. And then the sequence simply runs out at
the near end, the way a road runs out, and behind it is the white of a
page nobody wrote on.

He is fifty-five years old. He is one of the four people in Britain who
could reasonably be described as an authority on how communities stop
saying things. And he stands in his mother's bedroom with her life laid
out chronologically on the duvet, and thinks, with the clean shock of
something one has been avoiding at eye level for decades: \emph{I have
never once made her tell me.}

\sceneline

Her coat is on the back of the door.

It is navy, wool, from a shop on Green Street that closed years ago, and
it is the last thing in the flat and he has left it deliberately until
last, and when he takes it off the hook the door feels different in the
room. He stands with it over his arm. It is a coat. It is eleven hundred
grammes of wool and a lining that has gone at one shoulder. It does not
hold anything.

He hangs it back on the hook.

He tells himself he will take it next time, and there is a next time,
and he does not take it then either, and in the end he hands the keys to
the agent with the coat still hanging on the back of the door, and does
not mention it, and drives home, and this remains the single most
uncharacteristic act of his adult life.

\sceneline

The woman in the upstairs flat is called Mrs.~Ogilvie and she comes down
on the third weekend with a tin of shortbread and stays fifty minutes.

She is kind and slightly deaf and clearly lonely, and she wants to tell
him that his mother was a lovely neighbour, very private, wonderful in
the snow. Rowan makes tea and lets her talk, and does what he always
does with an informant, which is to say almost nothing and leave the
gaps open.

``She never had people,'' Mrs.~Ogilvie says. ``Not that she was
unfriendly. She'd stop and talk on the path for half an hour. But she
never had people in.''

``No.''

``Nineteen years I lived over her.'' She takes a shortbread. ``Do you
know, I couldn't have told you where she was from.''

Rowan holds his cup. ``What would you have guessed?''

``Oh --- '' Mrs.~Ogilvie waves vaguely, ``--- north somewhere? There was
a thing in the vowels sometimes, on the phone. Not often.'' She thinks
about it. ``I did ask her once. Years ago. And she said --- '' and here
Mrs.~Ogilvie's face does something interesting, a small backwards
groping, ``--- she said something perfectly nice, and I remember thinking
what a nice answer it was, and I couldn't tell you now what it was.''

She laughs at herself.

Rowan does not laugh. He is sitting very still, because he has just
watched a technique he has spent his career describing, performed by a
master, on a woman who did not notice it happening and has not noticed
in the forty years since. A response with the correct affective shape
and no informational content. Warmth as a sealant. He has published on
this. He has a diagram.

``You're a professor, aren't you,'' Mrs.~Ogilvie says comfortably.
``What is it you do?''

``Languages,'' Rowan says. ``More or less.''

``Lovely,'' says Mrs.~Ogilvie, satisfied, and Rowan hears himself do it
--- the identical manoeuvre, the same afternoon, in the same room ---
and does not correct it.

\sceneline

The tin is under the bed, at the back, behind the vacuum cleaner.

It is a biscuit tin, decorated with a Highland scene that has no
plausible connection to anything, and it does not rattle. He sits down
on the floor with his back to the bed and takes the lid off, and the lid
is stiff, and comes free with a small sucking noise.

Three things.

A sealed envelope, cream, unstamped. On the face, in his mother's hand,
four characters and a full stop. \emph{for E.B. only}.

A photograph, unlabelled on the front. A young woman on a shore of black
rock, wind in her hair, the surf behind her the colour of pewter, and
she is squinting, and half-laughing at whoever holds the camera. She is
perhaps seventeen. Rowan looks at it for a long moment as a photograph
of a stranger --- notes the coat, the period, the rock type, dates it
loosely to the mid-sixties --- and then the eyes come up out of it and
find him, and the room tilts, and he puts the photograph face down on
the carpet and sits with his hand flat on the back of it.

On the reverse, in pencil, one word. \emph{Brynjavík.}

And forty or so scraps of paper, all sizes, all in her handwriting.

\sceneline

He takes the scraps home and it is three weeks before he understands
them, and the three weeks are not idleness. He looks at them almost
every night.

They are Icelandic. That much is instant. They are not, as he assumes
for the first fortnight, phrases --- not copied proverbs, not vocabulary,
not the notes of an adult learner. He assumes an adult learner because
that is the ordinary explanation and because it is bearable: a woman in
her forties in a small flat in Grantchester, teaching herself the
language of a country she came from, privately, the way people take up
the piano. He builds this picture in some detail. He is fond of it.

In the third week he lays them all out on the floor of his own front
room, in no order, because he has run out of better ideas, and he
crouches over them with a mug going cold on the mantelpiece, and he sees
it in about four seconds.

They are drills.

Each scrap carries one stem, run through a fixed sequence of endings.
The hand is even. The spacing is even. There are no crossings-out. There
is no English anywhere on any of the forty scraps --- not a gloss, not a
translation, not a note-to-self, nothing. A learner glosses. A learner
cannot help it; the whole point of a learner's notebook is the bridge
back to the language they own. These are not bridges. These are
repetitions, produced at speed, by a hand that did not have to stop and
think about what came next.

And then, because he cannot not, he reads the sequence.

Nominative. Accusative. Dative. Genitive.

And then three more columns.

He kneels on the carpet with his glasses pushed up on his forehead and
looks at the fifth, sixth and seventh columns of a paradigm that has
four, and the headings are abbreviated to two or three characters in a
convention he does not recognise, and on every single scrap the seventh
column runs off the right-hand edge of the paper, because every single
scrap has been torn.

Not cut. Torn, by hand, along the same margin, forty times.

Rowan sits back on his heels.

He is very good at this. That is the thing he will think about later, in
a rented cottage nine hundred miles north, with the sea writing on the
glass: that he was very good at this, and that it took him four seconds,
and that the whole apparatus of a career was available to him on the
floor of his own front room. He has, sitting on the shelf above him,
a monograph with his name on it about the material traces of suppressed
grammatical categories. He has given the plenary. He knows precisely
what it means when a paradigm is drilled past the point where the living
language has cases to fill, and he knows exactly what it means when the
part past the fourth column is missing, forty times, by hand.

He gets up. He puts the kettle on. He comes back and gathers the scraps
carefully into a stack, squares them, and puts them in the tin with the
envelope and the photograph, and puts the lid on.

He does not write it down.

That is the whole of it, and it takes four seconds and lasts thirty
years, and no one is in the room. He has a notebook on the arm of the
chair, because there is always a notebook on the arm of the chair, and
he does not pick it up. He carries the tin upstairs and puts it in the
bottom of the wardrobe, and goes to bed, and sleeps badly, and in the
morning he goes to a supervision and is, by common consent, on
excellent form.

\sceneline

The grant proposal takes him eleven days and it is the best thing he has
written in a decade.

\emph{Morphological attrition in a terminal speech community: case
syncretism and lexical avoidance on the north Icelandic coast.} He does
not lie in it. This is important to him afterwards, and it is true: there
is not a false sentence in the document. The site is genuinely
under-described. The dialect is genuinely conservative. He is genuinely
the correct person. Every claim is supportable and every citation is
real and the methodology is sound and he has costed it to the pound.

It is simply that the document does not contain the reason.

He submits it in April. The committee meets in June. He is asked two
questions, one about equipment insurance and one about whether four
months is long enough, and no one asks him the third question, which is
the only question, and which he has rehearsed an answer to eleven times
while shaving.

The answer he prepared is: \emph{My mother was born there.}

He never has to use it. He is not sure, afterwards, whether he would
have.

\sceneline

Marianne catches him in the corridor outside the departmental office in
the last week of term, with a coffee in each hand and her bag falling
off her shoulder.

``I hear you got Iceland.''

``I got Iceland.''

``God, well done. Nobody gets fieldwork any more.'' She wedges the bag
back up with her elbow. ``Whereabouts?''

``North coast. Tiny place. Hundred people.''

``Lovely.'' She is already half past him, and then she turns back,
pleasantly, in the doorway, holding both coffees, and says: ``Isn't your
mother Icelandic?''

And Rowan hears himself say, in a tone of warm and complete
friendliness:

``She left a long time ago.''

Marianne says, ``Well, quite,'' and goes.

He stands in the corridor for a moment.

It is a flawless sentence. He admires it, briefly and coldly, the way one
admires a joint in wood. It is true. It answers. It closes. It does not
mention that his mother is dead, or that she died eight months ago, or
that he has her photograph in his desk drawer twenty feet away, or that
he is going to the village she was born in, or that he has a sealed
envelope in a biscuit tin at the bottom of a wardrobe with somebody's
initials on it. Every one of those facts is relevant. Marianne would
have been kind about all of them. She is a kind woman.

\emph{She left a long time ago.}

He walks back to his room, and sits down, and does not work for
twenty minutes.

\sceneline

He books the flights in the last week of August.

The confirmation comes through while he is standing in his kitchen with
the door open on a warm evening --- Keflavík, one way, hold luggage,
seat 14A, a car reserved from a firm in Reykjavík with a scraped white
Dacia in it that he does not know about yet. He reads the email on his
phone with a glass of wine in his other hand and the smell of somebody
else's barbecue coming over the fence.

Nothing has happened.

That is worth setting down plainly, because of everything that comes
after: on this evening, in this kitchen, nothing whatsoever has happened.
A woman died of a stroke at seventy-seven in a hospital in Cambridge.
Her son emptied her flat, methodically, and found a tin under the bed
containing a photograph, an envelope, and some old papers, which is
approximately what one finds under beds. He is a linguist and the
papers were in Icelandic and the place in the photograph has an
under-documented dialect and he applied for a grant and got it. There
is no mystery here. There is a bereaved man of fifty-five with a good
research question and an airline confirmation, standing in his kitchen
in the last of the light.

He finishes the wine and washes the glass and puts it on the rack.

Then he goes upstairs, and opens the wardrobe, and takes out the tin,
and does not open it, and puts it into the top of the duffel bag that is
already half packed on the spare bed --- and then takes it out again, and
finds a wool sock, and wraps it, and puts it at the bottom, underneath
everything, where a bag would have to be fully emptied to reach it.

He stands looking into the bag for a while.

\emph{For the padding,} he thinks, in English, in the nominative, and
hears the sentence do its work, and goes to bed.

\chapter{}

The Flybus takes forty-five minutes from Keflavík and Rowan spends all
of them looking at moss.

It grows over the lava in a grey-green pelt, unbroken, kilometre after
kilometre, and it is the first thing he has seen in eight months that
does not require him to have an opinion about it. He rests his temple
against the cold glass. Somewhere in the hold, in a duffel he packed
four times, is a tin box wrapped in a wool sock. He has not thought
about the tin since Heathrow, which is a lie he permits himself the way
he permits himself the second whisky on a night flight, knowingly, and
with a certain administrative respect for his own dishonesty.

The woman across the aisle is explaining to her husband that the moss
takes a century to recover from a boot. She says it in Icelandic. Rowan
understands every word, including \emph{að jafna sig,} to right itself,
to come back to level, and he feels the small warm click of comprehension
that has been the reliable pleasure of his working life. Twenty-two
years. He can read a saga at speed, follow a radio panel show, argue
about a bill. He has never lived here for longer than a fieldwork season
and he has an accent that will not come out, a flatness in the vowels
that marks him at the first sentence, and he has made his peace with
that. A foreign accent is only a signature. It says where the hand came
from. It does not say the hand is unsteady.

He does not know yet that this is the last week of his life in which he
will regard his own fluency as an asset.

\sceneline

The hotel is on Laugavegur above a shop selling puffin keyrings and
volcanic-ash soap. The room is small and very clean and smells of
lemon and dry heat. He unpacks by his method --- laptop, recorder,
notebooks, the empty tin --- and aligns each object parallel to the
edge of the desk, and stands looking at the result for a moment before
he goes down.

The clerk is perhaps twenty-five, with a septum ring and an English
paperback face-down on the counter. She takes his passport and types
without looking at the keys.

``Business or holiday?''

``Work.'' He says it in Icelandic and she answers in Icelandic, which
does not always happen, and he notes that it has happened, and is
pleased, and notes the pleasure. ``Fieldwork. I'm a linguist.''

``Where?''

``North. A village called Brynjavík. Small place, up past---''

``Brynjavík.'' She stops typing. It is not a dramatic stop. It is the
pause of a person retrieving a word from a shelf they have not visited
in a while. ``My grandmother grew up there.''

``Really.'' He leans on the counter. This is the part of the work he is
good at, the part that cannot be taught: the flat friendly opening of a
door. ``Does she still---''

``She died. Two thousand and nine.'' The clerk resumes typing. ``She
talked about it. The harbour and the boats and the wall they built when
she was a girl. She had a whole thing about the wall.''

``Have you been?''

``No.'' She slides his passport back.

``It's a long way,'' Rowan offers.

``It's four hours,'' the clerk says, and then frowns very slightly, at
the screen, the way a person frowns at a sum that has come out wrong
twice. ``I go to Akureyri all the time. My boyfriend's from Dalvík.''
She hands him the key card. ``I don't know. We never went. I never
thought about it, actually.'' She looks up, and her expression is
perfectly open and mildly interested, a woman making conversation at the
end of a shift. ``That's strange, isn't it. My own grandmother.''

``People are busy,'' Rowan says.

``People are busy,'' she agrees, and picks up her paperback, and that
is the end of it.

He goes up to the room and writes it down, because he writes everything
down.

\begin{fieldnote}
Hotel clerk, Laugavegur. Grandmother b.\ Brynjavík, d.\ 2009. Clerk has
never visited despite regular travel to the same region. No stated
reason. No apparent avoidance --- she raised it herself, unprompted, and
was mildly surprised by her own answer.

Cf.\ standard pattern: descendants of emigrants from declining rural
districts routinely do not visit. Nothing remarkable. Filed for
completeness.
\end{fieldnote}

He looks at what he has written. \emph{Filed for completeness} is what
he writes when he has recorded something for a reason he cannot yet
give. There are four hundred such entries in twenty-two years of
notebooks, and perhaps nine of them ever came to anything, and the ratio
is the discipline.

\sceneline

The Árni Magnússon Institute occupies a building of pale concrete and
enormous windows, and Dr.~Ragnheiður Ólafsdóttir occupies an office in
it that is warm as a greenhouse and stacked to the ceiling with offprints
in three languages. She is in her sixties, broad-shouldered, in a
cardigan the colour of oatmeal, and she is so immediately, unguardedly
generous that Rowan feels the small shame of a man who came prepared to
be charming.

``Sit, sit. Move those.'' She waves at a chair holding a decade of
\emph{Íslenskt mál}. ``You are the Cambridge one. Morphological
attrition.'' She says the phrase in English, with relish, as though it
were the name of a cocktail. ``I read the Faroese paper. The one about
the dative going soft in the fishing vocabulary.''

``That was ten years ago.''

``It was good. Most of them are not good.'' She pours coffee without
asking. ``So. What are you doing to us?''

He tells her. The north coast, the peninsula, a small parish with an
unusually conservative dialect. He does not say \emph{my mother was born
there.} He has rehearsed not saying it, and the rehearsal works, and he
is aware of the machinery working, and hates it slightly.

``Brynjavík,'' Ragnheiður says. She swivels to a shelf, runs a finger
along it, does not find what she is looking for, swivels back. ``Yes.
Interesting little dialect area.''

``Interesting how?''

``Oh --- conservative, as you say. Retentions. Some case usage that
looks archaic and is probably just isolation, you know how it goes, one
family with a strong grandmother and suddenly the whole fjord has a
feature.'' She sips. ``There is a lot of it up there. The whole
north-west is a museum of things that should have died in 1890.''

``Have you worked on it?''

``No.''

He waits. She does not elaborate, and there is nothing in her face to
elaborate --- no shadow, no evasion, only a woman who has answered a
question completely.

``Any particular reason?'' he asks. ``It sounds like exactly your
thing.''

``It is exactly my thing.'' She says it cheerfully. Then she pauses, and
looks past him at the window, at the flat white sky over Vatnsmýri, and
for the first time in the conversation she takes slightly too long.
``I nearly did. Twice. In --- ninety-six, and again about eight years
ago, when we had the mapping money. Both times it came off the list.''
She turns her hands over. ``The first time I think there was a funding
problem. The second time I do not remember. Something else was more
urgent. There is always something more urgent.'' She smiles. ``That is
what a career is, Dr.~Hale. It is the list of things that came off the
list.''

It is a good line. She has clearly used it before. He writes it down
because it is charming, and it is only in the car, four days later, that
he will take it out and look at it again.

She is helpful for another hour. She gives him three names, two of them
dead, and the third a retired schools inspector in Sauðárkrókur who will
not answer his letters. She tells him where the parish records should
be and where they will actually be, which are different buildings. She
warns him about the road.

``It is not a bad road. It is a road that was built for a bigger place
than the one at the end of it.'' She writes something on a card. ``Read
these two before you go. Halldórsdóttir on the coastal case system, and
there is a thing by Björnsson in \emph{Íslenskt mál} --- volume
nineteen, I think, or twenty --- comparing four of the northern parishes.
Old, but it is the only comparative work anyone has done.''

``Four parishes?''

``Mm. Brynjavík and three others.'' She hands him the card. ``Hvalsker
is one of them, I remember, because I stayed there once as a student and
the guesthouse had a cat with no tail.''

\sceneline

The Halldórsdóttir is exactly where she said it would be.

The Björnsson is not in volume nineteen.

Rowan spends forty minutes on this, which is thirty-five minutes longer
than the problem deserves, because the problem is trivial and the
solutions are obvious and he works through them all. It is not in
nineteen. It is not in twenty. He checks eighteen and twenty-one on the
principle that a scholar misremembering a volume number is the most
ordinary event in the world, and it is not in those either. He searches
the index of the journal for \emph{Björnsson} across the whole run and
finds eleven articles by four different Björnssons, none of them
comparative, none of them about the north coast. He searches for the
four parish names. He searches for \emph{samanburður,} comparison. He
finds nothing.

He sits back. The reading room is quiet in the particular way of
national libraries, a hush composed of ventilation and the small dry
sound of other people's pages.

The explanation is not mysterious. Ragnheiður is sixty-something and
reads three hundred articles a year and was recalling, without notes, a
paper she last looked at two decades ago. She has confused a journal, or
an author, or --- most likely of all --- an article she was told about
and never actually read, which is a thing every scholar alive has done
and none of them admits. Rowan himself once cited, in print, a
monograph that turned out to be a conference abstract. He is not
troubled. He is mildly, professionally amused.

He notes it anyway.

\begin{fieldnote}
Björnsson, comparative study of 4 northern parishes, \emph{Íslenskt mál}
vol.~19 or 20 --- not located. Journal run complete on shelf. Author
index searched. Not present under any volume.

Almost certainly misremembered attribution (cf.\ my own Hafnarfjörður
citation, 2011, which did not exist either). Will ask R.Ó.\ by email.

N.B.\ she named Hvalsker as one of the four and produced a specific
sensory memory attached to it (guesthouse, cat, no tail). People do not
usually invent cats.
\end{fieldnote}

He requests the parish holdings before he leaves, and goes to lunch, and
does not think about it again until the following morning.

\sceneline

The catalogue lists seventeen.

He works through the slips at a desk by the window with the winter light
coming flat across the wood. Parish registers, three volumes.
Confirmation rolls. A visitation report from 1874. A bound run of
district survey annals. A schoolbook. Two manuscript items described
only as \emph{blandað efni,} mixed material, which in his experience
means either treasure or a churchwarden's laundry accounts. A hymnal.

Six arrive.

The young man who brings them is apologetic in the way of someone who
has delivered this news many times and has stopped finding it
interesting. ``Sorry. These four are showing as not-on-shelf --- that
usually means they went out for conservation and the record didn't get
updated. These three are showing a location that doesn't exist any
more.'' He turns the slip around to demonstrate. ``See, that's a
basement stack we haven't used since the refurbishment. And this one
just says \emph{deposited}, which is not a location, that's someone in
1978 not doing their job.''

``And the last two?''

``Those ones are the annoying ones.'' He almost smiles. ``They're in the
catalogue and there's no accession record. So somebody typed them in off
a list at some point --- probably off the old card index in the
nineties, when everything got keyed in --- and either the books never
came to us at all, or they came and were never accessioned, or the card
was wrong in the first place.'' He shrugs. ``Retroconversion. You would
not believe what is in our catalogue. We have a chair.''

``A chair?''

``A chair. Catalogued, with a shelfmark. Nobody knows.''

Rowan laughs, and means it.

``The district ones are worse, though,'' the young man adds, gathering
the rejected slips. ``Anything that came in through the county offices.
There was a burst pipe in eighty-one and a fire in a records store in
Blönduós before that, and a lot of it was never listed properly to begin
with. Half the north is like this. It's not just your parish.''

\emph{It's not just your parish.} Rowan turns that over for a moment and
files it, because it is genuinely reassuring, and because he notices
that he needed reassuring, and that is the part worth recording.

Six of seventeen. He works through the six until the light goes. The
register is a register. The visitation report of 1874 complains about
damp in the vestry and the poor attendance of the fishing families in
season and the state of the pastor's fence, and is so completely,
gloriously ordinary that Rowan reads all nine pages of it with something
close to affection. The hymnal in the catalogue is not the hymnal;
it is a Reykjavík printing of 1902 with nothing in it at all, and he
checks page nine out of a habit he does not yet have a reason for, and
finds \emph{verði þinn vilji,} standard, correct, unremarkable, and
closes it and thinks nothing whatever about it.

\sceneline

The map shop is on a side street off Bankastræti and smells of paper and
electric heat, and Rowan is happy in it in the uncomplicated way he is
happy in very few places.

He buys three sheets of the same coast.

The current 1:50,000, which is the one he will drive with. A 1:100,000
tourist sheet, printed the year before last, of the sort sold in petrol
stations. And, from a box at the back, a 1987 survey sheet, foxed at the
folds, which he buys for eight hundred krónur because he has never in
his life left a shop containing an old map of a place he is going to.

He spreads all three on the hotel bed that evening.

The 1:50,000 gives the population as 113. The tourist sheet gives 96.
The 1987 sheet gives 41.

He sits down on the edge of the bed.

Then he gets up and makes tea in the little plastic kettle and comes
back and looks at them properly, because the first reading of anything
is worthless.

The 1987 figure is thirty-eight years old and is presumably a census
figure from a census taken before he finished school. The tourist sheet
is derivative and lazy and probably lifts its numbers from whatever
database was cheapest. The 1:50,000 is the survey and is presumably
right. Three sources, three vintages, three numbers --- and the numbers
do not decline. That is the only genuinely odd thing, and he sits with
it, and works it out in under a minute: 41 in 1987 and 113 now would
mean the village nearly tripled, which does not happen to Icelandic
fishing villages, but 41 in 1987 could easily be a misprint for 141,
or a figure for the settlement proper against 113 for the whole
parish including the outlying farms, or --- most probably --- the
1987 sheet is simply using a different unit than he assumes. Rural
population figures are a swamp. He has published on a related swamp.

He photographs the three sheets side by side anyway.

The roads take longer.

On the 1:50,000 the connector road from the pass to the village is drawn
as a solid single-lane track, class F, with a spot height and a bridge
symbol. On the tourist sheet it is dashed --- unsurfaced, seasonal, the
convention for a road that may not be passable. On the 1987 sheet it is
not drawn at all. The village sits at the end of nothing, reachable in
the cartography only by the coastal track from the east, which the
1:50,000 shows as long abandoned.

He looks at this for a while.

Roads are built. Roads are downgraded. A road that did not exist in 1987
and exists now is a road that was built between 1987 and now, which is
not a mystery, it is a public works project. He is being asked to be
unsettled by infrastructure.

\begin{fieldnote}
Three sheets, same coast.

Pop.: 113 (1:50k, current) / 96 (tourist, 2024) / 41 (survey, 1987).
Non-monotonic. Probable unit mismatch (settlement vs.\ parish) or
transcription error in the 1987 figure.

Connector rd.: solid class F (current) / dashed seasonal (tourist) /
absent (1987). Consistent with construction post-1987.

\struck{Check date of road construction --- district records, Sauðárkrókur.}

Have checked. There is no road construction recorded for that district
between 1970 and 2003. The budget lines are complete and itemised for
every year and there is nothing in them but resurfacing.

Which means either the road existed in 1987 and the surveyors left it
off the sheet, or it was built without an entry.

The first of these is ordinary. Sheets have errors. Every sheet has
errors. I have found errors on sheets myself, twice, and reported them,
and received a polite letter.
\end{fieldnote}

He reads that last paragraph back and recognises, dimly, the sound of a
man reassuring himself in writing, and does not know what to do about
it, so he underlines nothing and closes the book.

\sceneline

The parish office of the national church is on the fourth floor of a
building with a lift that announces each floor in a woman's recorded
voice, and the woman who receives him is called Gerður and has been
doing this job, she tells him, since before the computers.

She is delighted to be asked. Nobody asks.

``Brynjavík, Brynjavík.'' She types. ``Yes. Active parish. Resident
incumbent --- séra Hjálmar Pétursson, installed two thousand and six.
Before him séra Guðrún Ingimarsdóttir, and before her --- '' she
scrolls, ``--- oh, this goes back a long way. Do you want it all?''

``How far back does it go?''

``Seventeen thirty-two.'' She says it without any particular reverence.
``That is when the returns start. Some parishes are older but the
paperwork is not.''

``And it's continuous?''

``It is continuous.'' She turns the monitor slightly so he can see, which
she is probably not supposed to do. Column after column of years and
names, unbroken, three centuries of a small parish sending its numbers
to the capital. ``That is unusual, actually. Most of them have gaps ---
the wars, the eruptions, the influenza, somebody's grandmother dying and
the box going in the attic. This one has no gaps.''

``None at all?''

``None at all.'' She scrolls to the bottom. ``And they are always on
time. Look --- the returns are due the fifteenth of January. Every year
they arrive in the first week. Every year.'' She laughs, a short fond
laugh. ``I have parishes in this city that I chase for four months. And
then there is Brynjavík, one hundred and thirteen people at the end of a
road, and every January the first week, like a clock.''

She says \emph{one hundred and thirteen} and Rowan feels the number land
against the three maps on his hotel bed, and says nothing, and thanks
her.

At the lift he stops.

``Sorry --- one more. Are there any records of services not being held?
Cancellations?''

``Oh, they log those.'' Gerður brightens; this is clearly a good
question. ``The parishes are meant to record any Sunday without a
service, with the reason. Most of them do not bother any more.'' She
crosses to the terminal again and taps. ``Brynjavík. Three.''

``Three this decade?''

``Three since seventeen thirty-two.'' She turns and looks at him over
her glasses with the pleasure of a specialist producing her best
exhibit. ``With the reasons written out --- a blizzard, a funeral in
another parish, and one in 1918 that just says \emph{illness}, which
will have been the influenza and which I find rather awful in three
letters.'' She sits back. ``Do you know how strange that is? I have been
here thirty-one years and I have never seen a parish that logged
three.''

He rides down four floors in a lift that names each one, and stands on
the street in the blue afternoon dark, and finds that he is smiling.

Because it is a good detail. It is the kind of detail that makes a
funding body sit up: an isolated community with an extraordinarily
strong institutional continuity, a parish that has not missed a return in
two hundred and ninety-three years. It is culture. It is
\emph{conservatism}, in the technical sense, exactly what he flew here to
find, and the fact that a village of a hundred and thirteen people
maintains its paperwork with more rigour than a parish in central
Reykjavík is not sinister, it is the whole thesis, it is a community
that has decided, over centuries, that certain records must be kept.

He walks back to Laugavegur through the wet light and buys a coffee he
does not want and thinks: they are careful people. It will be a good
site.

That is the last afternoon of his life in which \emph{careful} means
what he currently thinks it means.

\sceneline

The rental is a white Dacia with eleven thousand kilometres on it and a
scraped wing the desk agent photographs from four angles.

He loads it in the dark of eight in the morning: the duffel, the
equipment case, the box of photocopies, and, on the passenger seat where
he can reach it, the 1:50,000 sheet with the route pencilled in and the
warnings marked --- false turnouts, a slide zone, a bridge replaced by a
culvert. The tin goes at the bottom of the duffel, wrapped in a wool
sock, because he cannot make himself carry it any other way.

He sits for a moment with his hands on the wheel and the engine running
and the heater roaring at the windscreen.

Six documents of seventeen. A road that appears on one sheet and not
another. An article that does not exist in the volume where a good
scholar remembered it. Three population figures that do not resolve. A
woman who has never visited her grandmother's village and cannot say
why. A parish that has not missed a Sunday in three hundred years.

Six things. He has been on eleven fieldwork sites and there is not one
of them that would not yield six such things in four days --- the
archives of small places are made of this, of pipes bursting and clerks
retiring and boxes going into attics, and a researcher who cannot look
at institutional noise without hearing a signal in it is not a
researcher, he is a man with a hobby. Every one of the six has an
explanation. He has produced the explanations himself, unprompted,
without strain. That is the job. That is precisely the job.

He pulls out into the Reykjavík traffic and joins Route One heading
north, and the city gives way to suburb and the suburb to lava, and by
the time he clears Borgarnes the light has come up grey over water on
his left, and he is thinking about nothing much at all.

It is only somewhere past Blönduós, four hours later, on an empty
straight with the mountains going blue on his right, that he hears
himself, in the privacy of the car, say aloud, in Icelandic, without
having decided to:

\emph{``Sex hlutir.''}

Six things.

And then, because it is a habit older than his career, because he has
been declining nouns in his head since he was nine years old in a house
where a woman would not tell him where she was from, he declines it. Out
loud. To the windscreen. Nominative, accusative, dative, genitive.
\emph{Sex hlutir, sex hluti, sex hlutum, sex hluta.}

The last one sits oddly in his mouth. \emph{Of six things.} Possessive.
As though the six things owned something.

He turns the radio on, and it plays a pop song about the summer, and he
drives north into the shortening afternoon.

\chapter{}

He feels the Icelandic wind before it strikes. On the final hairpin
above Brynjavík a gale broadsides the rental, lifting the left wheels,
dropping them back with a jolt that travels up through his spine. Rowan
grips the wheel and counts the gusts: one, two, three, four. Here they
never come in pairs. He read that once, in a footnote, and filed it
under local lore, and now, on the road, it feels less like folklore than
like a rule someone forgot to explain.

The single lane curls along a razor's edge, salt-pitted guardrail to one
side, black-rock void to the other. He hugs the rail, though instinct
tells him the rail is not the danger. The danger is the place itself,
too exposed, too willing.

On the passenger seat lies his 1:50,000 survey map, his thumb tracing
the penciled warnings: false turnouts, slides, a collapsed bridge
patched by a culvert. At kilometer forty-seven the village should
appear. The odometer ticks over, 46.7, 47.0, and a fluorescent stake
flashes past. He brakes onto gravel, reverses, inches forward, studies
it. Too sharp. Too new. Planted eight hundred meters early. Not drift,
not a surveyor's rounding. He photographs the stake, the boot-prints in
the mud beside it, the flat horizon, and drives on at a crawl.

Fog rolls in like something with intent. Thirty meters of visibility,
then twelve, then nothing. Another stake surfaces, then another, spaced
without agreement, as if the men who set them had been quarrelling about
where the road belonged. Rowan thinks: the road is only here when the
ground consents to it. The thought arrives whole and unwelcome, and he
sets it down where he sets down all such thoughts, in the margin,
unsigned.

His GPS flickers and pins the village half a kilometer out to sea. The
blue dot hops inland, hunts, gives up. He closes the app and proceeds by
memory, which is the only instrument he has ever fully trusted.

A thinning of fog. A huddle of low roofs, ink-black against a grey that
is not quite sky. A weather-worn sign: \emph{Brynjavík}, and beneath it
a column of figures rather than a figure, the earlier numbers struck
through and the newer ones painted smaller each time, as though the sign
itself were being economical with what remained. \emph{604.} \emph{412.}
\emph{237.} \emph{158.} And at the bottom, in a hand that has given up
on the pretence of permanence, \emph{113.} A hundred and fifty years of
subtraction, kept on a board at the roadside the way a farmer keeps a
tally of lambs. He photographs it, and thinks that no other village he
has worked in has bothered to show its arithmetic. Most repaint. This
one strikes through and lets you count.

He parks beneath the single lit window of the community hall and kills the
engine. The silence is immediate. Only the wind humming in the power
lines, and under it, further down, the sound the sea makes when it is
not doing anything in particular.

Outside, the houses press close, eaves almost touching. The alleys
barely admit a body. The basalt cliffs stand at the back of the village
like a spine, and though the geology is young, the rock feels
inevitable, a thing the place has decided to lean against. He
photographs the sign, the buildings, the cliffs, then compares them
against the Street View he downloaded a month ago in a bright room in
Cambridge. A blue house that stood thirty meters back now abuts the
pavement. A playground has become gravel and a sodden bench. He paces
thirteen strides from curb to stoop and writes \emph{spatial
compression} in his field book, his thumb smudging the wet ink. He
shivers, and does not decide whether it is cold or arithmetic.

He checks his duffel out of habit, cataloguing: tape recorder, notepads,
and, at the bottom, wrapped in a wool sock because he could not make
himself carry it any other way, a small tin box. Inside the box, a
sealed envelope. On its face, in his mother's hand, four characters and
a full stop: \emph{for E.B. only}.

He does not take the box out. He knows the weight of it without lifting
it.

Elizabeth Hale had been reticent to the point of architecture; a woman
built of what she declined to say. When he was a boy he had asked once,
twice, about \emph{before Cambridge}, about her own parents, and each
time she had answered him in a way that was not an answer, a smooth grey
surface that gave the question nothing to grip. He had learned, without
deciding to, to read the shapes of her silences the way other children
learned to read faces. He had built, over thirty years, a set of tools
precise enough to translate a woman who would not be translated. He
calls this his vocation. He has never once called it what it is.

He still does not know who E.B. is. He knows that his mother wrote to
her once and did not send it. He knows there are forty scraps of paper
under the envelope carrying a paradigm with three more columns than the
language has, and a torn edge where the last of them should be. He knows
he has driven eleven hundred kilometres without once permitting himself
to finish the sentence that begins \emph{a person only drills a paradigm
that far if.}

He puts the duffel down. He reads the email printout: \emph{Arrival any
time after 17:00. Please come directly to office. --- Thorsteinn.} He
climbs the hall steps, photographs the door, and knocks, once with his
knuckles, once with his palm.

\sceneline

Inside, a man in navy, close-shaven, watching him. ``You're early. Good.
We've already lost an hour to the wind.''

The office smells of cleaning solution and wet wool. Rowan signs the
visitor log --- \emph{Dr.~R. Hale, Univ. of Cambridge, 17:34} --- while
Thorsteinn watches the pen. The wall behind the desk carries maps in
overlay, layer on layer of coastline, shorelines shifting inward year by
year, boundaries redrawn in a hand that grows less confident with each
revision.

``Subsidence,'' Thorsteinn says, when Rowan asks. ``Erosion. Other
things.'' He slides a heavy key across the desk. The tag reads a
six-digit code. ``Guest house first. No television. The wireless
reaches, some days. Breakfast at eight. The café is simple.''

Rowan takes the key. ``I'd like to speak with Edda Brynjarsdóttir.''

Thorsteinn's fingers pause on the blotter. It is a small stall, the kind
Rowan has spent his life learning to see. ``She knows the stories,''
Thorsteinn says. ``She will not always answer. If you go, bring a gift.
Salt. Matches. Good tea.'' He aligns a stack of forms that are already
aligned. ``And you will not ask her about the water by its name. Not the
sea. Not the bay. Not any word that points straight at it. Here we do
not do that.''

``How do you refer to it, then?''

``We don't refer to it.'' Thorsteinn says it flatly, as a man states an
occupational hazard. ``We work around it. You will learn the shape of
the hole by the things placed at its edges.'' He stands, which is the
end of the interview. At the threshold he adds, without turning: ``The
north edge is harder ground. Do not walk it alone.'' Outside, the wind
takes Rowan by the collar like a hand checking a coat for a name tag.

\sceneline

On the quay, four figures bend over green netting, oilskin and wool
layered like strata in a cliff face. The oldest has hands that look
carved by water. Beside him a woman works the twine without looking at
it. And a young one, no more than twenty, smoking, his fair hair
whipping loose from under a cap. Rowan waits for the rhythm of the
mending, net-shuttle and pull, before he steps into it.

``Góðan dag,'' he tries.

The woman nods. The old man grunts. Rowan gives his name, his purpose.
``Dr.~Hale. Cambridge. Here to document the dialect.''

``Second time you've passed,'' the old man says. ``Nothing much to
see.''

``You can take the photographs,'' the woman adds. ``The colours
flatten.''

The young one, the smoker, watches Rowan with more openness than the
others, the way the young in a closed place watch anyone who has been
elsewhere. ``Everything moves here,'' he says. ``Boats, quotas. The
village breathes. Some years the space opens. Some years it squeezes.''

``That's Bjarni,'' the old man says, in the tone of a man apologizing
for weather. ``He talks.''

Rowan writes \emph{subjective spatial variation} and, beside it, the
name. Bjarni.

``How do you refer to the water?'' Rowan asks.

The twine twitches and stops. Finally the old man says, ``We call it the
west.''

``Or the fog,'' the woman says.

``Or the deep,'' Bjarni ventures, and grins, testing it, a boy leaning
over a rail.

``The bay?'' Rowan offers.

They laugh, but it is not quite laughter. ``You're not one of us,''
Bjarni says. ``So you can say it. It slides off you.'' He says it almost
with envy.

``You name what wants naming,'' the old man says. ``You leave the
rest.'' He spits into the gap between the stones. ``It's not law. It's
how it's done.''

Rowan raises his recorder, and they give him the safe words the way one
hands over ballast: \emph{the edge, the salt, the slip. That way brings
the weather, not the fish.} He notes \emph{avoidant referents, formal
weight,} and, when he plays the file back that evening, he will find a
smooth blank where a certain word should sit. Not a glitch. A held
breath. Later he will recognize the quality of that silence, though he
does not know it yet. It is the exact texture of every answer his mother
ever gave him about her own childhood: a place in the sentence where the
floor had been carefully removed.

Bjarni nods toward the northmost house on the ridge. ``Talk to Edda. She
keeps the old grammar.'' He flicks his cigarette into the water and
watches where it lands, longer than a man watches a cigarette.

\sceneline

Her house is the last before the stone wall. The door is newer than the
walls around it, the paint unweathered. He knocks once, then again,
shifting the emphasis, and the door opens to a measured width.

Edda Brynjarsdóttir is smaller than the village made her sound and
denser, as if she has been compressed by the same pressure that
shortened the road. Her sweater is undyed wool the grey of granite. Her
eyes take his measure, the package, the stance, and settle.

``Dr.~Hale.'' Not a question. An entry in a ledger she has been keeping
a long time.

He states his name, his purpose, the gift. She accepts the tea and the
salt without a tremor. She looks past him toward the water, then back.
``Best to say little this early,'' she says. ``The weather's
undecided.'' Not a warning. A note on procedure.

Inside: wool and ash and dried herbs, shelves loaded in careful balance,
no mirrors. She sits and gestures him to the chair opposite. The window
behind her is set at an angle that shows nothing continuous, a fragment
of basalt, a slice of bay, the view deliberately broken.

He sets the recorder between them and does not switch it on. Instead he
takes the tin from his bag, opens it, and slides the envelope across the
table. \emph{for E.B. only}, in his mother's hand.

Edda reads the name. She does not react. Not a flicker, not a breath out
of rhythm. She lets it lie there on the wood between the salt and the
tea, and does not touch it, and does not open it. He watches her
not-react with a scholar's attention and finds, at the bottom of his
attention, something older and colder, the child at the closed door,
listening. He puts the child back in the margin.

He asks his prepared questions. Weather. Childhood. The shifting edge of
the village. She answers in phrases that compress more than they
release. ``Years tighten. Some widen. My mother spoke of a winter when
the distance doubled, and one when the houses pressed so close you could
not lose a child in the street.''

He asks about language. ``Words get thin at the edge,'' she says. ``The
endings don't always hold. Sometimes there's enough. Sometimes it is
nothing.'' She gives him three examples, the same stem, only the ending
altered, her vowels lengthening on the final syllable as if to make the
sound carry a weight.

He asks, obliquely, about the word no one uses. She does not answer. Her
hand closes around the tea, the knuckles whitening, and that is the
answer, and he marks it.

He thanks her. She does not acknowledge the thanks; she rises to tend
the stove. The interview is over. Procedure, not dismissal. As he
gathers his things his eyes go, once, to the envelope, still sealed,
still lying where he left it. She has not moved it toward herself. She
has not moved it away.

``You'll leave that,'' she says. Not a question either.

He leaves it.

\sceneline

The guest house smells of salt and the dry tick of the baseboard heater.
He unpacks, aligning each object parallel to the desk edge --- laptop,
recorder, notebooks, the empty tin. The window overlooks water that
merges with sky. He transcribes: \emph{Village compressed. Road
inconsistent. Dialect avoids forbidden referent.} He types, then
backspaces, every instance of \emph{ocean} and \emph{waterway}, settling
on \emph{the edge, the outer margin.} He does not decide to do this. He
notices he has done it.

Fog gathers against the glass and draws thin lines that bisect the view,
dew beading and running in ways he does not quite watch. It is nothing
yet. It is a page not yet written on. He lies down with the recorder
running, and the last thing he registers before sleep is that the tide
has a cadence, long troughs and clustered peaks, and that it sounds, if
he lets it, like a room in which someone is choosing not to speak. He
has slept in that room his whole life. He does not think to write the
forbidden word. He would not have known which one it was.

\chapter{}

Thorsteinn provides the key at 08:00. The prefab was a holding pen, he
explains, suitable for transients; for a scholar meaning to stay the
month, the Annex is required. It stands further up the slope, grounded
in the volcanic scree, its roof a heavy brow of sod turning brown in the
salt air.

Rowan carries his equipment in two trips. He counts one hundred and
twelve steps from the road to the threshold. The interior smells of
linseed oil and dry timber, a single room partitioned by use rather than
by walls: a sleeping alcove to the east, a workspace facing west. The
furniture is heavy, built of driftwood and bolted to the floor. He tests
the desk. It does not wobble.

He begins the setup, because method is a form of defence. He aligns the
recorder parallel to the window frame, three centimeters from the edge.
He stacks the notebooks by date of anticipated use. He builds a small
grid of order inside the cottage's stubborn geography.

The light changes, bars of cold illumination falling across the
floorboards, and he sees the wear pattern in the pine, a trough worn
from desk to window and another from stove to door. Not general traffic.
A loop. A patrol. He runs his hand along the desk edge and finds the
wood scarred, and leans closer, and the scars resolve into cuts. Not
scratches. Inscriptions.

He fetches his loupe. The markings are shallow, made with a blade that
was probably dull, and he recognizes their shapes at once, though their
presence here makes no sense he can hold.

\emph{-um.} The dative plural ending.

He moves to the window frame, where a hand would rest, and finds another
cluster carved into the sill. \emph{-ar.} Genitive singular. The
bedframe: \emph{-in}, the suffixed article. There are no stems anywhere.
The room is littered with grammatical tails, suffixes cut loose from
their nouns, waiting for a subject to fasten to. He photographs each
instance and labels the files by coordinate. He circles the room,
cataloguing. The carvings are not vandalism; they are corrections. They
cluster over knots in the wood, at joints where the carpentry failed, as
if a former occupant had tried to conjugate the room itself, to force
the crooked timber into a declension it could not hold.

He writes in his journal: \emph{Structure contains extensive marginalia.
Morphology appears reactive. Prior occupants strictly codified space.}
Then he sits, and the window frames the grey and the boundary where the
village meets the tide, and he does not name the water. He notes only
that the horizon is active.

\sceneline

A knock, single and dull. A woman on the slate step, in a coat of boiled
wool dyed the colour of wet peat, her hair pulled back from a face that
looks etched by the same blunt blade that marked the table.

``Dr.~Hale. I am Sigrún Ásgeirsdóttir.''

She wipes her boots on the mat, three strokes, left, right, left, before
she will step onto the floor. She sets a canvas satchel on the table,
avoiding the carved edge. ``Thorsteinn said you required histories. The
parish annals. Nineteen hundred to nineteen forty-five. After that the
records become sparse.''

She unbuckles the satchel and lifts out the black-bound volumes one at a
time, squaring each to the one beneath it. Last, she sets down a small
book bound in cracked calf, and hesitates over it half a second longer
than she hesitated over the others.

``Séra Hjálmar sends this. He says you will want the hymnal as well as
the register, since you are here for the speech and not the sheep.'' A
pause. ``He also says he is at the parsonage most mornings and that you
are not to knock before nine.''

``Hjálmar Pétursson,'' Rowan says. ``Installed two thousand and six.''

Sigrún's hands go still on the satchel buckle for about a second.

``You have been to the parish office.''

``I went to a great many offices.'' He hears how it sounds and does not
soften it, because he has learned in three days here that softening is
read as concealment. ``They were very helpful. They were pleased that
somebody asked.''

``They would be.'' She resumes with the buckle. ``There has been a
resident pastor in Brynjavík since 1732. There has been a service every
Sunday of that time except three, and the three are recorded, with the
reasons.'' She sets her hand
briefly on the calf-bound book, and takes it away again. ``He is a good
man. He will talk to you longer than anyone else here will. Do not
mistake that for the same thing as telling you more.''

``Thank you. I can make tea. I have the brand Thorsteinn suggested.''

``No tea.'' She keeps her coat on. She stands near the center of the
room, hands clasped, reading the space, or the density of the air in it.

``I've been reviewing the samples from the quay,'' Rowan says,
organizing her stack. ``There's a variance in how the locals address the
geography.'' He turns to the window and gestures at the horizon, at the
iron-grey line where sky meets surface. In fluent Icelandic he says,
``It's fascinating. The way the ocean dominates the perspective, yet the
vocabulary minimizes it.''

Sigrún's head snaps up. ``We do not use that construction.''

He pauses. ``The vocabulary?''

``The subject.'' Her voice goes flat, administrative. ``You placed the
agent in the water. We do not say \emph{the ocean dominates.} We say the
view is filled.''

He files it: \emph{passive required; agent removal.} ``A grammatical
avoidance of agency. So if I were to describe the state of the
waves---''

``Currents.''

``If I were to say the waves are rising---''

The light in the room dies.

It is not a cloud. It is an erasure. Fog hits the window with the force
of a solid body, the glass white and blind in an instant, the wind that
had been a constant thrum cut clean out. The silence is a seal applied
to the cottage. Rowan looks at the window, then the lamp, which flickers
once and steadies. His watch's second hand still sweeps. The room feels
held.

Sigrún has not moved. She does not look at the window. She is looking at
his hands, and adjusting her cuff.

``The forms must be precise,'' she says, her tone exactly what it was,
as though the blindness of the room were a small clerical thing he has
caused and must correct. ``This close to the glass, Dr.~Hale, you do not
give action to the outside. The outside is a location. It does not
act.''

The white thins to grey wisps. The wind bleeds back, treble first, then
the low thud of the turf roof.

``I understand,'' Rowan says, level, though his pulse has spiked and he
notes the spike the way he notes everything, at a distance, as data.
\emph{Grammatical breach followed immediately by environmental
obscuration.}

Sigrún nods. The lesson is not moral; it is structural. ``The annals are
in order. Do not read them past sunset in poor light; it strains the
eyes. And the carvings.'' She is at the door now, not looking back. ``Do
not let your fingers trace them too long. They are unfinished for a
reason.'' The latch, and she is gone, and the window is clear again,
framing the thing it is forbidden to name.

\sceneline

The annals are bound in black cloth, spines blank. \emph{District
Survey, 1912.} The paper is brittle and smells of acid and time. Crop
yields, sheep counts, an influenza outbreak in November, and then the
fishing reports, in which the catch is meticulously recorded --- cod,
haddock, weights, prices --- but the source of the fish is never named.
The sentences bend into the passive or stagger under heavy
circumlocutions. \emph{The nets were drawn from the grey expanse. The
boats returned from the lower meridian. The yield from the westward open
was substantial.}

A map folded into the back cover renders Brynjavík in fanatical ink
detail, every fence line, every pen. But to the west, where the shore
ends, the mapmaker simply stopped. No waves, no soundings, no depth. A
blank cream void, labelled with a single coordinate: \emph{64°N ---
Undesignated.}

He logs it. \emph{Cartographic erasure. Linguistic containment.} And,
because the phrase belongs to a colder discipline than the one he means
to practice here, he underlines it and does not know why.

He sets the annals aside and takes up the hymnal, expecting nothing from
it. Icelandic Lutheran hymnals are among the most standardized documents
in Europe; the text has been fixed by synod and reprinted without
meaningful variation since the nineteenth century, and he opens it the
way a man opens a phone directory in a hotel room, to be doing
something.

He finds the Lord's Prayer on page nine and reads it without attention,
and stops in the third line.

\emph{Faðir vor, þú sem ert á himnum. Helgist þitt nafn, til komi þitt
ríki, verði að þínum vilja, svo á jörðu sem á himni.}

He reads it again. Then he reads it aloud, quietly, to hear whether his
eye is lying to him.

\emph{Verði að þínum vilja.} Dative. \emph{Let it be done according to
thy will.}

Every hymnal in Iceland reads \emph{verði þinn vilji.} Nominative. The
will is the subject. The will is the thing that does. It is one of the
oldest and most stable lines in the language --- it survived the
Reformation, it survived Guðbrandur's Bible, it survived four centuries
of orthographic reform --- and here, in a parish hymnal printed, he
checks the plate, in Akureyri in 1928, the divine will has been quietly
moved out of the nominative and set into an oblique case, where it can
be measured against but cannot act.

\begin{fieldnote}
Hymnal, Brynjavík parish, Akureyri impression 1928. Pater Noster, l.~3.
Standard text \emph{verði þinn vilji} (nom.) replaced by \emph{verði að
þínum vilja} (dat., governed by \emph{að}). Not a misprint --- the
setting is clean, the line is centred, the spacing is correct for the
longer reading. A misprint costs a letter. This cost a plate.

\struck{Provincial variant? Check Akureyri press for house readings.}

\struck{Hyper-corrective? Rural registers overuse \emph{að} + dat.}

Neither. Both those hypotheses predict variation elsewhere in the text.
There is none. The rest of the prayer is standard to the comma. One
case, changed once, deliberately, in the only line of the prayer in
which something other than the speaker is permitted to want.
\end{fieldnote}

He closes the hymnal and does not move for a while. He is a
morphologist; he has spent twenty years being pleased by things exactly
this size. He waits for the pleasure and it does not come. What comes
instead is the first cold suspicion that he has been handed a
document by a woman who wiped her boots three times before she would
step on the floor, and that she knew what page he would stop on.

The hours compress. The light fails, first to the long northern twilight
and then to dark, and he switches on the lamp, and his eyes drift to the
beam above the desk where \emph{-um} catches the light. Back to the
book. Page forty-two, a sentence that ends without ending: \emph{The
storm took the boat and its crew to the---} No object. No full stop. He
looks at the beam. \emph{-um.} The dative plural. \emph{To the {[}dative
plural{]}.} The occupants did not carve these marks to damage the wood.
They were storing the endings they were forbidden to speak. Offloading
the grammar into the timber, the way a man buries what he cannot keep in
the house and cannot throw away.

He stops at the window.

The exterior dark is total. The glass returns the room to him --- desk,
lamp, his own pale face --- and then something else begins to form on
the surface. Condensation. The room is dry; the heater cycles evenly;
yet the moisture gathers, and it does not run in streaks. It clusters.
He steps closer and adjusts his glasses.

The water is writing.

A curve. A vertical stroke. A crossbar. The droplets hold against
gravity and shape a character. \emph{Þ.} Thorn. Beside it, another
cluster. \emph{u.} Then \emph{ng.}

\emph{Þung.} Heavy. Or pregnant. Or difficult.

He does not gasp. He observes the surface tension holding the beads in
place where surface tension should not, and files the improbability, and
lifts his hand. The heat of his palm should break the formation, should
merge and run the droplets. He brings his index finger to hover over the
character.

The pane vibrates. A low frequency, barely audible, a cello string bowed
too slowly, travelling the air gap into his skin. Patterned. Short,
long, short. He presses his fingertip to the \emph{Þ.}

The droplets do not smear. They set. The vibration spikes, a sharp
closed-circuit buzz, and beneath his finger the character dissolves and
reforms. \emph{S.} Then \emph{j.} Then \emph{ó.} He pulls his hand back.
The vibration stops. The condensation loses cohesion and weeps down the
glass in ordinary chaotic rivulets, and the word it was becoming is gone
before he can decide whether he wanted it to finish.

He stands in the center of the room. The lamp hums. The books lie open
to their amputated sentences. His finger is wet. He is not afraid, he
tells himself, and hears, distinctly, how the sentence is built to hold
something at arm's length, and recognizes the construction, because it
is his mother's, because it is his. \emph{He is not afraid.} Agent
removed from the water. The fear placed nowhere, so that it cannot act.

He revises his earlier hypothesis. The cottage is not a blank page. It
is a prompt. He sits and picks up his pen and does not write a
conclusion. He writes a question. \emph{System active. Interaction
confirms syntax is physical. What is the required case for entry?}

Outside, the wind taps against the wall. One. Two. Three. Four. Rowan
waits. The lesson is beginning, and some part of him he will not look at
directly is not surprised to find that he has been enrolled a long time.

\chapter{}

Rowan takes the path to the north ridge, his boots finding purchase on
slate that feels brittle, as if the salt air has over-cured the very
geology. The wind here does not buffet; it presses, a laminar flow that
lies the grass flat against the soil. He carries his field bag ---
canvas, waterproofed, holding the recorder, a stack of legal pads, three
pens tested for ink flow.

On the quay below, small in the fog, Bjarni is coiling line. He lifts a
hand when Rowan passes --- the only villager who has yet greeted him
unprompted --- and Rowan lifts one back, and files the fact that the
young man watches the water while he coils, the way a person watches a
dog he half trusts.

Edda's cottage sits low in the scree, a huddle of turf and basalt that
seems less built than extruded. The door is painted a red so deep it
reads as black in the twilight. He checks his watch, 14:00 exactly, and
knocks, and the door opens before his knuckles can graze the wood a
second time.

She is in the same undyed wool she wore on the first day, and Rowan
suspects it is not the same garment but the only kind she owns. Her
hands, resting on the frame, are roped with veins that suggest high
pressure held a long time.
``Dr.~Hale,'' she says. It is a verification of inventory.

``Ms.~Brynjarsdóttir. Thank you for agreeing to speak again.'' He offers
the smile he perfected in Cambridge common rooms --- disarming,
deferential, slightly worn.

She steps back, letting a wedge of light spill onto the stones. ``The
weather holds. Come in before the pressure drops.''

Inside, the air is dry and smells of burnt peat and something sharper,
ozone or metal filings. The room is a masterclass in compression.
Shelves on every vertical surface, jars, folded linen, tools whose
purposes have eroded from memory. A single paraffin lamp burns on the
central table, its flame motionless in its glass chimney. The envelope
is nowhere in sight. He does not ask. She gestures him to a chair of
heavy timber, back straight.

He begins his ritual --- unzip the bag, remove the recorder, check the
battery, align the device equidistant between them, notebook to the
right, pen uncapped. He feels the familiar hum of academic extraction.
He is here to document, to classify, to save. It is a story he tells
himself about why men like him go to the ends of coasts, and he has told
it so often it has the smoothness of truth.

``I appreciate your time,'' he says, and his voice sounds too round in
the damped acoustics. ``Thorsteinn says you're the custodian of the
older forms.''

``Thorsteinn manages the maps.'' Her vowels are drawn out, stretched
like wire. ``I manage the speech.''

His eye catches the ceiling beams above her, old drift-timber dark with
smoke. Carvings, like the Annex, but denser. The wood is not scarred; it
is annotated. He recognizes the declension markers --- \emph{-um, -ar,
-i} --- and, crowded into the grain like refugees seeking high ground,
stems he does not know.

``You've marked the beams,'' he says, pointing with his pen. ``Mnemonic
aids?''

``Structural reinforcement.'' She does not look up.

He notes it --- \emph{informant views grammar as architectural} --- and
presses record. ``Let's begin with the fishing terminology. The
specialized lexicon for the tools, the nets, the hooks. How has the
vocabulary shifted since the herring collapse?''

She folds her hands and waits two beats. ``The collapse was not of
herring. It was of counting.'' Her voice lacks the melodic rise and fall
of standard Icelandic; it is flat, rhythmic, engineered, built in blocks
of sound. ``When the nets went down, we pulled up weight. But the weight
had no name. To name a thing is to fasten it. We had no fasteners for
what came up.''

He writes quickly. \emph{Lexical gaps.} ``So you borrowed terms? Coined
new ones?''

``We widened the old ones. We stretched the vowels to hold more water.''
She demonstrates --- the word for net, but the vowel stays in her throat
and resonates, turning the monosyllable into a four-second hum that he
feels behind his sternum.

He looks past her to the window on his left. The glass is clear. He half
expects it to cloud, the way the Annex glass clouded, and is aware of
expecting it, and disciplines the expectation away. Instead his eye goes
to the floor near the hearth, where a low draught has worked the
peat-ash into fine spirals, tight and repeating, the same loop over and
over, as if the room exhales in a hand he is beginning to be able to
read.

``Fascinating,'' he says, confident again on the ground of description.
``Phonological lengthening as a marker for mass. Applied to all nouns,
or only those tied to extraction?''

``Only those that touch the boundary.'' She watches the digital counter
turn. ``You are recording the sound. The sound is only the casing.''

``We can keep to morphology, if you prefer.'' A slight irritation. He
wants data, not metaphysics. But the data is coming, wrapped in the
metaphysics, and he is not yet willing to admit it comes no other way.
``How do you distinguish the shallow catch from the deep?''

``Depth is not a distance.'' She leans forward, and the lamp flame
stretches upward though there is no draught. ``It is a case. We use the
dative for the bottom. We use the genitive for the surface.'' She
pauses. ``Because the surface belongs to the wind. And what belongs to a
thing is genitive. But the bottom is not owned. The bottom is
\emph{given to.} So it takes the dative. You see --- it is not
arbitrary. It is only consistent.''

He writes it, and against his training the consistency snags at him,
because it is consistent, because the rule reaches down and holds its
own logic the way a keystone holds a span. \emph{Spatial relations
mapped to case. Vertical axis: genitive surface, dative floor. System
internally coherent.}

He turns to a clean page and rules it, because a thing this shaped wants
a grid, and because ruling a grid is what he does with his hands when
his mind has begun to move faster than he is comfortable with.

\begin{paradigm}
\begin{tabular}{@{}llll@{}}
& \textsc{surface} & \textsc{column} & \textsc{floor}\\[2pt]
\hline\\[-6pt]
case & genitive & --- & dative\\
sense & \emph{of} & --- & \emph{given to}\\
holder & the wind & --- & unowned\\
attested & yes & --- & yes\\
\end{tabular}
\end{paradigm}

He looks at the middle column for a while.

An empty cell in a paradigm is not a mystery; it is the most ordinary
thing in comparative morphology. Paradigms have gaps. English has no
past participle for \emph{stride} that anybody will commit to. What
makes him sit back, tapping the pen against his thumb, is that this gap
is not a defect. The other two columns are built around it. The system
does not fail to describe the middle. It declines to.

\begin{fieldnote}
Vertical axis, prelim. Two attested terms, one prohibited.

\struck{Gap = attrition. Middle term lost, remaining terms reanalysed.}
--- Against: attrition leaves debris. Lost forms leave doublets,
hypercorrection, disputed usage among elderly speakers. There is none
here. Informant produced the two attested cases without hesitation and
refused the third without hesitation. That is not a hole in a system.
That is a system with a designed hole.

Note also: she did not say \emph{we do not have a word.} She said
\emph{we do not speak of the middle.} Absence of vocabulary vs.\
prohibition of use. Different phenomena. Different diagnoses.
\end{fieldnote}

``And the middle?'' he asks. ``The water column itself?''

Her eyes snap to his. ``We do not speak of the middle. The middle is
transit. Transit requires silence.''

He smiles, a reflex of intellectual satisfaction, and files it ---
\emph{taboo keyed to verticality} --- and does not notice that the
ash-spirals on the floor have stopped turning in a circle and begun to
draw, very slowly, the same open shape twice, side by side, like a pair
of parted lips.

\sceneline

At half past three she stands up and puts a pan on the stove without
asking him anything, and Rowan understands that he is being fed and that
declining is not among the available moves.

It is fish stew, and it is very good, and it is served in bowls that do
not match. There is rye bread from a tin and butter with too much salt
in it. She eats quickly, the way people eat who have spent their lives
around work, and does not talk while she does it, and the silence in the
room is different from the silence of the interview --- looser, less
load-bearing, the ordinary quiet of two people at a table.

Afterwards she pours coffee and sits back and says, unprompted:

``You want the stories.''

``I want whatever you'll give me.''

``Everyone wants the stories.'' She turns her cup. ``The collectors came
in nineteen sixty-five and they wanted the stories. They came again in
seventy-one. My mother gave them the hidden people for two afternoons
and they went away very pleased.''

``And were they true? The hidden people?''

Edda considers this longer than he expects.

``The country has the stories,'' she says. ``All of it. Every farm from
here to Höfn will tell you not to move that rock because of who lives in
it, and they will tell it to you with a straight face and half of them
half mean it, and it is a fine thing and I would not take it from
anybody.'' She sets the cup down. ``But ask them why. Ask a man in
Selfoss why you do not move the rock and he will say it is bad luck, or
his grandmother said so, or that a bulldozer broke once at Kópavogur.
He has the instruction and he has lost the reason, so he keeps it as a
story, because a story is what an instruction turns into when nobody
remembers what it was for.''

``And here?''

``Here we have the reason.'' She says it without any pride at all. ``So
we do not need the story. We do not have hidden people in Brynjavík,
Dr.~Hale. We have a rule about what may be given a house and what may
not, and the rule is grammatical, and the rock is one of the places we
apply it. The rest of the country turned that rule into a small
frightening family who live under a stone and take offence.'' Something
moves at the corner of her mouth. ``It is not a bad translation. It is
just a translation.''

Rowan writes very fast for about a minute.

``The landvættir,'' he says. ``The land-spirits. Four of them, the
dragon and the eagle and the bull and the giant, one at each quarter,
guarding the country. That's in Heimskringla, it's on the coat of arms,
it's on the coins in my pocket---''

``It is four boundaries,'' says Edda. ``And a rule about who may
approach from which direction, and in what case. Somebody put a face on
each one so it could be remembered by people who were not going to be
taught the grammar.'' She shrugs, a small economical movement. ``That
was well done. It has lasted a thousand years. Ours has lasted a
thousand years too and it has needed a hundred and thirteen people
working at it every month, and theirs needed a picture on a coin. I am
not saying we chose better.''

``Second sight,'' Rowan says. ``\emph{Skyggni.}''

``Noticing which case a thing is in before it has finished arriving.''

``Dreams that tell you a name.''

``The only place a name is safe to say.'' She looks at him. ``Where did
you think the dreaming came from? You may not speak the name of a dead
man on this coast. So he comes into the dream and says it himself, and
you wake with it in your mouth, and you have committed nothing, because
you did not say it, you were told it.'' A pause. ``That is not a
superstition. That is a workaround.''

He has stopped writing. She notices, and does not comment.

``And the naming of children?'' he asks. ``The taboo on repeating a
name?''

``No taboo. The opposite.'' She refills his cup without being asked,
which is the closest thing to affection she has shown him. ``You are
obliged to reuse them, and you must reuse them in order, and you must
not skip. Every child in this village carries a name that has been in it
before.'' She sets the pot down. ``There are eleven names. That is what
there has always been. When one goes out of use we are short, and being
short is a thing we notice.''

``Eleven for a hundred and thirteen people?''

``Eleven \emph{keeping} names.'' She says the word flatly, in the local
form, and it is the first time she has used it in front of him. ``The
rest may be called whatever their mothers like. My mother wanted a
Kristjana and got me instead.''

\sceneline

She tells him about her mother's winter without his asking for it, which
he will realise later was the only unprompted gift of the entire
afternoon.

``Nineteen thirty-four,'' she says. ``I was not born. She was nine.''

She was nine, and in January of that year the distance from the church
door to the harbour wall doubled. Not in one night --- over eleven days,
by degrees, so that a child sent for bread in the morning came back
later than she should have and could not explain why, and was scolded,
and the next day was later still, and by the ninth day the schoolmaster
was keeping the children at the school because it was no longer
reasonable to send a nine-year-old the length of a village that had
become forty minutes long.

``Nobody was frightened,'' Edda says. ``That is what she always came
back to. Nobody was frightened, because everybody was busy. There is a
great deal to do when a place gets bigger --- the water has to be
carried further, the boats are further off, the old people cannot get to
the hall. It was eleven days of extremely tedious inconvenience and
everyone was too tired at the end of it to be afraid.''

``And then?''

``And then it came back. Over about a month. Less neatly than it went.''
She turns her cup a quarter turn. ``Two of the houses did not come back
in quite the same relation to the others, and there was an argument
about a boundary that ran until 1961.''

Rowan sits with that.

``Your mother told you this as a story,'' he says carefully. ``Or as an
event?''

``She told it to me the way you would tell a child about the year the
roof came off.'' Edda's eyes come up. ``Which is what it was. And do you
know what the useful part is, Dr.~Hale? Not the eleven days. The eleven
days are the part that would go in your paper.'' She leans forward
slightly. ``The useful part is the argument about the boundary, because
it went on for twenty-seven years, and because both families were
right.''

\sceneline

The interview pivots. He feels the momentum is his. ``I want to come
back to the water itself. In the village I noticed a reluctance to use
the standard noun \emph{hafið.} Instead I heard \emph{the slip,}
\emph{the edge.}''

``Names are addresses,'' she says. Her thumbs press together until the
nail beds whiten. ``You do not shout an address unless you want a
visitor.''

He nods. \emph{Name taboo. Sympathetic magic.} Standard anthropological
fare, and he pushes past it. ``But syntactically --- the ocean functions
as an agent in your folklore, doesn't it? Does the sea take the boats?
Or are the boats lost to the sea?'' He watches her for the reaction to
the active construction. ``The sea,'' he repeats, clarifying, testing.
``The sea rises. The sea claims.''

The reaction is not emotional. It is electrical.

The paraffin flame gutters --- does not dim, but detaches from the wick,
floating a centimeter as a blue sphere before slamming back down into
yellow. Simultaneously the wind outside stops. It does not fade. It
cuts. One moment the cottage is under a gale, the next the walls are
silent, the roof still, and Rowan's ears pop with the change of
pressure.

Edda's spine locks against the chair. Her eyes narrow, not with fear,
with calculation. ``Stop,'' she says. A flat line.

``I apologize. Did I offend?''

``You constructed it wrong.'' She exhales, a hiss. ``You gave it the
nominative. You made it the actor.''

``I was asking a theoretical---''

``There is no theoretical here.'' The wind returns, changed, a low
rhythmic thud against the north wall, like a heavy hand testing the
turf. ``The grammar is the physics. If you say it rises, you give it the
knees to stand.''

He writes --- \emph{informant displays acute distress at nominative
assignment} --- and marks it with an asterisk, and is aware, distantly,
that his hand is steadier than the rest of him. ``Let me rephrase. How
would you describe a storm surge? Water coming into the village.''

She smooths the wool of her sweater, recalibrating, her lips moving
silently over syllables. ``At the boundary,'' she says slowly, ``there
is an accumulation.''

``Passive voice,'' he notes.

``Locative,'' she corrects. ``We locate the water. We do not let it
move. There is water at the door. Not \emph{the water comes to the
door.}''

``Because if it comes?''

``Then it has intent. And if it has intent, it has memory.''

The recorder crackles, a wet tearing sound, like fabric ripping under
water. The meters redline for a fraction of a second. ``Interference,''
he says, tapping the device. ``RF noise off the coast.''

``It is not noise. It is the unbinding.'' She looks at the black box the
way one looks at a trap that has snapped. ``We use the instrumental case
for the tide. \emph{By means of the tide, the boat is lifted.} Never
\emph{the tide lifts the boat.}''

He nods, engrossed, categorizing --- an ergative shift, agency stripped
from nature to survive it, a trauma encoded in syntax, brilliant,
publishable. ``And if someone breaks the rule? A child? An outsider?''

The static bursts again, louder, a high shriek that cuts off the instant
he finishes the word \emph{outsider.} Edda does not flinch. She watches
the window, where the glass, still clear, has begun very faintly to
breathe a film onto itself, but she does not look at it as if it were
remarkable. ``If the rule is broken,'' she says, without affect, ``the
sentence must be finished by something else.''

He feels a chill that has nothing to do with the draught and attributes
it to the damp.

\sceneline

``There is a form,'' she says, ``used by the line-men when the fog was
thickest. To be sure the voice was still a human voice. It requires the
breath to be held in the middle. You must not breathe between the
subject and the verb. If you breathe, the gap lets the wet in.''

``I'd like to document it.''

She closes her eyes and begins, low, vibrating in the lower throat, a
sequence of long vowels and glottal stops. ``Áaaa\ldots{} rrrr\ldots{}
unn\ldots{} nnn---''

The air in the room tightens. The recorder screams --- not feedback, a
wall of white noise, absolute, the peak-lights locking on, the counter
scrambling forward and back at once. He grabs the device and hits stop,
and the screaming ceases, and the room rings.

``Interference,'' he says, his heart going hard against his ribs.
``High-frequency burst. A trawler's radar.''

Edda opens her eyes and regards the machine with mild curiosity, as one
regards a trap. ``The machine does not like the forward construction.''
He puts on the headphones. Nothing but a roar of static that sounds,
strangely, organic, like the crushing of millions of shells.

``Can you repeat it? Slower --- I need the phonemes.''

``I cannot say it slower. The speed is the lock.'' She considers him.
``But I can say it backward.''

``Retrograde phrasing? Inverted syntax?''

``I mean backward. Behind the verb, where it is safe.'' She sees he does
not follow. ``Forward is the future. To say a thing forward is to invite
it toward you. We do not want the future here; we want the present to
hold. So the dangerous nouns go behind the verb, in the place the
sentence has already passed. What has already passed cannot arrive.''
She says it plainly, and he hears the rule close on itself again,
discoverable, load-bearing: the direction of a phrase is the direction
of time, and time is the thing they will not let move.

She speaks again, the melody inverted, the vowels clipped, the
consonants striking the air like flint. ``Und\ldots{} hrafn\ldots{}
i\ldots{} ver---'' The recorder spins silently, the levels bouncing
green. It captures it cleanly.

He breathes out. ``That got it.''

``Placement,'' she says, and stands, signalling the end, her joints
moving as if calcified. He packs, thrilled --- \emph{temporal
orientation linked to syntactic directionality} --- and at the door his
eye sweeps the window one last time.

The glass is no longer transparent. From frame to frame it is fogged,
but the fog is not random; it has organized into rows of tight angular
script, thousands of tiny weeping characters that mirror the carvings on
the beams. The window is a page, and the text is dense and frantic and
unmistakably grammatical.

He pauses, hand on the latch. ``You might check the insulation on that
frame. The buildup is significant.''

``It is just the residue,'' she says, turning to the stove. ``Of what we
spoke.''

He steps out into the cold, and the transition jars him. The light has
failed. The landscape is a blue-black bruise. He walks down the slate
path, structuring the paper he will write, and reaches the bottom of the
ridge, and stops.

The road is gone.

The tide has not risen so much as surged. Water laps the stone wall of
the nearest sheep pen, fifty meters inland from where the shore lay two
hours ago. It is smooth, glassy, sliding over the grass without a sound.
It does not crash. It occupies. He checks his watch; the charts do not
predict this. He looks toward the village, where the hall lights burn
and a figure crosses a window and a car crawls the upper track. No one
is running. No one is looking at the water that has swallowed the lower
road. The village carries on, ignoring the inundation, the way a family
carries on around a subject that has been agreed, without a meeting, to
be closed.

He stands at the new waterline. The water touches the toe of his boot.
It feels deliberate. He pulls his coat tighter and begins the climb,
stepping carefully, so as not to disturb the syntax of the rising dark.

\chapter{}

The procession does not march; it filters. Like groundwater through
shale, the villagers of Brynjavík seep up the incline toward the basalt
platform that overhangs the bay. There is no coffin. The coffin is
already in the ground in the high churchyard, under a Lutheran service
of unremarkable correctness that Rowan attended an hour ago and could
have attended in any parish in Iceland: the hymn, the reading from
Corinthians, the three handfuls of earth, Séra Hjálmar in his ruff and
cassock saying the words in a warm unhurried baritone while the wind
worked at the pages. Nothing in the churchyard was strange. That is
what Rowan cannot get past. Nothing in the churchyard was strange, and
then the entire congregation walked out through the lych-gate, turned
west instead of down, and climbed toward the water without anyone
appearing to have decided to.

The body has been committed. The grammar of departure is still pending.

He had asked Hjálmar about it on the path, quietly, the way one asks a
question one suspects is rude. The pastor had considered the horizon for
a while.

``The burial puts him with God,'' he said at last. ``That is settled.
That was settled at the graveside and I do not think it can be unsettled
by anything we do on a rock.'' He shifted his grip on the prayer book,
which he carried against his chest the way Rowan has begun to carry the
tin. ``What is not settled is where he is. Where the body is. Where the
name is. Those are different questions from the question of the soul,
Dr.~Hale, and the Church has never claimed the second one.'' A small dry
smile. ``It has occasionally claimed it in other countries. It did not
get much thanks.''

``And the parish permits this? The second service?''

``The parish \emph{is} this,'' Hjálmar said, without any edge at all.
``I bury them. Then I walk up the hill with them and I stand at the back
and I hold my tongue, because the part that happens on the rock is not
mine and I am not good at it.'' He looked at Rowan directly for the
first time. ``I am fifty-one. I have been here nineteen years. I could
not tell you where the one ends and the other begins, and I have stopped
regarding that as a failure of scholarship.''

He counts eighty-nine heads, which is most of the village and all of the
part of it that matters today; the rest are children too young to hold
still and the two Sigurjónsdóttir sisters, who have not spoken to each
other since a dispute over a mooring in 2009 and who cannot be relied
upon to stand in the same line. He adjusts the gain on his recorder,
sheltering the microphones from the salt wind with a fur windshield. The
silence of the group is not emptiness. It is held breath, pressure
behind a valve. This is not grief as he has known grief. No weeping, no
black veils, no bowed heads. The postures are upright, braced. They have
come to see that a debt is settled before the day closes.

He knows a version of this silence. Eight months ago he stood at the
edge of another committal, in a crematorium chapel in Cambridgeshire
that smelled of lilies and warm dust, and did not weep either, and was
privately relieved that no one required it of him. His mother had left
instructions so spare they were almost a grammar of their own: no
service, no address, the ashes to the college garden, and --- in a line
he had read four times --- \emph{nothing said over me that I would not
have said myself.} He had honoured it. He had said nothing. He had
thought, at the time, that this was a kindness he was doing her. He
watches the villagers brace against the wind and revises the thought,
and files the revision, and does not look at it again.

They gather on the shelf of volcanic rock, planed level by centuries of
ice. The sea below is a churning grey slurry, heaving against the cliff
with a rhythm that is dangerously close to a heartbeat. It does not
crash. It thuds. A man steps to the edge, weathered, his knuckles
bleached by cold and work. He does not look at the water. He looks at
the horizon, that smudge where grey meets grey, and raises his hands,
palms down. The suppression gesture.

The speech begins. Vikamál, low and guttural. Rowan's headphones fill
with it, long open vowels stretched thin to accommodate the wind, and he
watches the waveform crawl green across his screen. The speaker is
reciting the Sæfangi, the Sea-Holding, a lattice of datives and
instrumentals, a way of giving without naming, of placing the dead into
the category of matter rather than identity. \emph{From the bone, the
salt. To the salt, the weight. By way of the drift, the return.}

He writes: \emph{Agent removal complete. Objectification of the deceased
prevents nominative claim.} The crowd breathes in sync with the
caesuras. They are the redundancy. If the speaker fails, the weight of
their listening should hold the line.

But there is a gap in the line. Edda stands ten meters off, near a
jagged outcrop, coat buttoned to the chin. She does not watch the
speaker. She watches the water, a patch of foam swirling at the cliff
base. She is the one checking the sum.

The speaker's voice rises toward the knot that ties the ritual shut. He
speaks of the gathering. \emph{``Hérna kemur\ldots{} safnið---''} He
falters.

It is microscopic, perhaps two hundred milliseconds, a stutter on the
suffix. He means to cast \emph{safn} --- the gathering, the collection
--- into the singular instrumental, a mass noun, indivisible. But his
tongue trips and adds a syllable. \emph{``\ldots safnanna.''} Genitive
plural. \emph{Of the gatherings.} He has split the single mass into
countable parts. He has implied there is more than one return, more than
enough to go around.

To an outsider it is a cleared throat. To the system it is a breach.
Rowan sees it before he hears it. Edda goes rigid --- not a flinch, a
seizure, as if a current has run through her spine. The cords at her
throat stand out and vibrate. Her lips shape a correction, a
counter-sound to null the plural, and she does not speak it. She clamps
her jaw. She knows: to overlay the error now would be to confess it to
the listener below, to admit the grammar can be bent. Her silence is
terrible. It is a masonry block dropped into a well.

The speaker pales; he knows too. He rushes the last words, smearing the
consonants, trying to lose the mistake in the noise of the gale, and
steps back from the edge wiping his hands on his trousers as if to clean
off the wrong ending. The villagers do not gasp. They shift --- a
collective adjustment of weight, shoulders closing the gaps between them
--- and they look at their boots, and they do not look at the sea. Rowan
glances down at the churning water. It does not look angry. It looks
attentive. It looks like it is counting.

\sceneline

The wind does not build. It arrives.

At 02:00 the Annex takes a single percussive blow. Rowan wakes --- not
from sleep but from the shallow drift of his own notes --- and goes to
the window, the cold glass shivering against his fingers. Outside, the
storm is taking the village apart.

It is not a rage. It moves with a precision that defies weather. It
bypasses the fragile glass of the greenhouse; it leaves the loose
shingles of the hall. It goes for the storage sheds along the eastern
perimeter. He watches the wind peel the corrugated roof from a
net-drying hut --- not tear it, unstitch it, the sheets lifting in
sequence, one, two, three, folding back like the leaves of a book being
read. A gust hammers the chimney of the old salting house and the bricks
tumble inward, filling the throat of the hearth, while the rest of the
roof stands untouched. A boathouse flush against the tide loses its
eastern wall entire, the frame stripped clean, the boats inside exposed
to spray --- and the oars racked along the far wall do not so much as
rattle.

He counts the strikes. They are rhythmic. They fall on the things that
hold things: sheds, lockers, boundaries. The storm is attacking the idea
of keeping. He remembers the error on the cliff. \emph{Safnanna.} The
plural of the hoard. The sea is correcting the inventory. The word for
what he is watching wants to be a bureaucratic one, and he lets it come
once, because it is exact, and then he lets it go: the sea is not
auditing anything. It is grieving in the only grammar it has, which is
weight and water, and it has been told there is more than enough, and it
does not agree.

He stays at the glass until the wind drops, as suddenly as it came, at
04:15. The silence after is heavier than the noise. The village lies
grey and open under the first light, smoke rising from the fallen
chimney like a flag of surrender.

Then, thin and human, cutting up through the quiet from the direction of
the harbour: a shout. Then another, and the particular ragged pitch of
men calling a name.

\sceneline

He is down the hundred and twelve steps before he has decided to go.

At the harbour mouth the water is still heaving from the storm, and a
small skiff --- someone had gone out too early, or never come in --- is
broadside to the swell, half-swamped, its mooring line parted and
streaming. A figure is in the water between the boat and the quay wall,
and the figure is Bjarni. Rowan knows the fair hair even flattened
black. The young man is not swimming. He is being moved, lifted and
dropped by the swell against the barnacled stone, one arm hooked over
the parted line, his face going under and coming up and going under, and
each time it comes up it comes up lower.

Three villagers are on the quay. One has a boathook. None of them has
gone in. Rowan understands, in a flat instant that he will turn over for
weeks, that they are not frozen by fear. They are calculating grammar.
To call his name across the water is to hand the water a nominative, a
handle. To reach for him and fail is to give the sea a completed action.
They are trying to find the phrase that saves him without naming him to
the thing that has him, and the phrase is taking too long.

Rowan has no such phrase. He has only the old, crude, terrestrial
grammar of a body going toward another body. He goes down the weed-slick
steps, takes the boathook from the man who holds it --- the man lets it
go with something like relief and something like horror --- and wedges
himself at the waterline, and reaches.

``By the line,'' the man behind him says, urgent, in Vikamál. ``Not by
the arm. Give it the rope, not the name.'' And Rowan understands enough
now to obey: he does not shout Bjarni's name, does not reach for the boy
himself, reaches instead for the streaming mooring line and hooks it and
hauls, so that it is the rope that is subject and object of the
sentence, the rope that is pulled, the rope that pulls, and Bjarni only
a weight the rope happens to bear. The distinction is absurd. He makes
it anyway. The swell lifts them both and he leans back against it with
everything in his legs, and the old man is beside him now with a grip on
his coat, and the woman has the line behind, and they draw the rope in
and Bjarni with it, three, four, terrible seconds, and then his other
hand comes up over the stone and closes on nothing and then on Rowan's
wrist.

They land him on the quay like a catch. He is grey-white, his lips gone
the colour of the water, and he coughs up brine in long convulsions,
more of it than a chest should hold, and the woman turns his head so it
runs into the gaps between the stones and not back down. No one says his
name. No one says the word for what nearly happened. The old man says
only, quietly, in the passive that Rowan is learning is a kind of mercy
here, \emph{``He is kept.''} Not \emph{he is saved.} Not \emph{he
lives.} Only that he has, for now, been kept on this side of the
boundary.

Bjarni opens his eyes. They find Rowan crouched over him and hold there,
and something in them has changed --- not gratitude, or not only. There
is a new blankness at the back of the boy who talked too much, a
stillness the water has left in him the way it leaves salt on stone, a
thing that will not wash off. He tries to speak and what comes out first
is not a word but a note, low, held, one of the long resonating vowels
Edda made in her throat, as if the sea has taught him a syllable on the
way down and he has brought it back. Then he closes his mouth on it,
deliberately, the way the others close their mouths, and Rowan watches
the village's discipline enter a person in real time, watches Bjarni
become, in the space of a breath, quieter, older, marked.

``Don't take him to the hall wet,'' the old man says. ``Dry clothes
first. The wet remembers.'' They carry him up between them, and Rowan
carries part of the weight, and the weight is enormous and specific and
entirely of this world, and for once he does not reach for a phrase to
keep it at a distance. He simply carries it, and finds, when they set
the boy down by a stove and the danger has passed, that his own hands
are shaking, and that he is glad of it, glad of the plain animal fear,
because it is his, and it names nothing, and it is the truest thing he
has felt since the crematorium garden.

\sceneline

Later, in the debris field, the air smells of wet soot and split pine
--- the smell of a flooded archive. He keeps his camera in his pocket;
the lens feels too aggressive against the quiet industry. The villagers
work in a silence that is almost structural, a silence that holds the
space up. No calls for hammers, no warnings about nails. A man lifts one
end of a beam; a woman is already there for the other. A nod moves a
load, a glance stops a shovel. They move with heavy deliberate caution,
as if the air itself might shatter under a careless vowel.

He watches three men repairing the dry-stone wall of the salting yard.
The wall was not knocked down; it burst outward from within, the stones
splayed in a fan. He notes the angle. \emph{Expulsive force. Internal
pressure.} He steps closer, meaning to ask about the mortar --- and a
hand closes on his shoulder, heavy, calloused, immovable.

Arnar Jónsson. The harbour master looks as though he has not slept in a
week; the lines around his eyes are cut deep with salt. He smells of
diesel and old wool. He does not look at Rowan. He looks at the wall, at
the men slotting a block back into the breach.

``Do not ask them,'' Arnar says, a rumble that barely moves the air.
``They are concentrating.''

``On the mortar?''

``On the case. They have to close the sentence before the tide turns.''

Rowan adjusts his glasses. ``The damage is very specific. The boathouse,
the shed. Nothing residential. It looks almost curated.''

Arnar takes a pouch of tobacco from his pocket and rolls a cigarette
with shaking fingers. He does not light it. ``The speaker slipped,'' he
says. Not an accusation. A report, a mechanic noting a sheared bolt.

``I noticed. A phonological stutter on the final declension of
\emph{safn.} I have it recorded.''

Arnar turns his head slowly and looks at Rowan with a pity that is
devastating in its mildness. ``It wasn't a stutter.''

``Then what was it?''

``The sea heard the wrong ending.'' He says it the way an engineer says
the load exceeded the rating. ``It heard plural. The sea does not like
plural. Plural means we have enough to spare.'' He pockets a lighter he
never meant to use. ``We do not have enough to spare. We nearly spared
one this morning, in the harbour.''

Rowan says, because his training says it for him, ``You believe the
performance failure produced a collective expectation of punishment, and
the community read the storm damage as targeted---''

Arnar looks at him, and for a moment the distance between them is not
intellectual but nearly biological. ``Dr.~Hale. You think the words
describe the storm. You have it backward.'' He turns to go, boots
crunching wet gravel. ``The words are the weather. Yesterday we asked,
by accident, for the roof to be taken. And the boy this morning --- the
sea was only reaching for the nearest handle we left lying out.'' He
walks toward the quay.

Rowan looks down at his notebook, at the neat rows of phonetic symbols.
His pen hovers over the page where, an hour ago, in a warm room, Bjarni
had coughed up the sea and then swallowed a vowel rather than let it
out. He does not write \emph{ritual entropy.} He does not write
anything. He caps the pen.

The sea is flat now, disturbingly calm, the waves against the jetty
small and evenly spaced, breaking with a soft hiss-thud, hiss-thud. He
listens. The rhythm is the cadence of the funeral chant, stripped of the
error, looping. Hiss-thud. Noun, verb. Hiss-thud. Subject, object. He
reaches for his recorder and touches the cold casing in his pocket and
hesitates. If the sea listens to endings, then recording it is not
observation. It is a dialogue. And he does not know the case for
\emph{listener,} or whether it is safe to be one. He leaves the recorder
in his pocket and steps back from the edge, retreating toward the Annex,
while behind him the villagers seal the wall and lock the silence back
inside the stone.

\chapter{}

The dawn arrives not as an event but as a gradual desaturation of the
dark. The sky over Brynjavík is the colour of wet slate,
indistinguishable from the water except by texture; the sky is smooth,
the water grained with the restless geometry of the tide.

Rowan finds Arnar Jónsson where the basalt shelf descends into the foam.
The harbour master stands ankle-deep in the kelp and drift, his back to
the horizon. He does not look like a man expecting company. He looks
like a piece of the shoreline that has detached itself to run a
maintenance cycle.

Rowan approaches, boots on the scree. He holds his notebook to his
chest, the recorder primed but hidden. He has learned from the funeral.
He does not brandish the device; he lets it listen from the dark.

``Good morning.'' The greeting feels too round, too whole for the hour.

Arnar does not turn. He is working a length of green nylon, his fingers
finding a fray, excising it with a small knife, the motion surgical.
``The wind,'' he says. He does not finish. He does not say \emph{is
cold} or \emph{has died.} He acknowledges the noun and lets it hang,
unanchored.

``I wanted to ask about the repairs. The wall's holding?''

Arnar nods. He lifts a cracked plastic float, turns it to show the
fissure, names nothing, sets it in a crate. Then he says, unprompted,
which he has not done before: ``Not the proper time for it. The wall.''

``There's a proper time?''

``The Keeping.'' He gestures with his chin up the coast, an economical
sweep that takes in the sheds, the wall, the boundary stones, the whole
inventory of the place. ``Fourth Saturday. Everyone. The walls, the
nets, the roofs.'' He works the knife along the fray. ``Doing it now is
doing it twice. But it will not keep until the fourth Saturday, so.'' A
shrug of the shoulders that is almost an apology to the schedule.

``Every month?''

``Every month.'' Arnar says it the way a man confirms that the bins go
out on a Tuesday.

``And that's enough?''

``For a wall.'' He sets the float down. ``Not for the rest.'' He looks
up the coast, and something in the set of his shoulders shifts --- not
solemnity, Rowan will think later, but the particular weariness of a man
contemplating a large annual task in the wrong season. ``The Great
Keeping is at the turn of the year. That one is different. That one is
the whole boundary, every stone, and the reading.''

``The reading of what?''

``The names.'' Arnar picks up the next float. He does not elaborate, and
Rowan has learned enough not to push a man who has already given three
sentences more than his quota.

Rowan writes it down. \emph{Monthly maintenance day, universal
participation --- walls, nets, roofs. Plus annual observance at
midwinter: full perimeter + recitation of names (list? genealogy?
necrology?). Ask Sigrún re: parish overlap.} He underlines
\emph{universal}, and files the whole thing under folk custom, under the
sort of thing a village does because a village has always done it, and
does not yet think to ask what a hundred and thirteen people are
maintaining when they walk a boundary in the dark and read a list of
names aloud to nobody. ``The stone,'' he says.
``The placement holds.'' He picks up a net shuttle, traces its spine
with his thumb, mimes filling it --- does not say \emph{shuttle,} does
not say \emph{needle,} draws its long figure-eight in the air and points
to the pile of mending. ``For the tear.''

Rowan notes it --- \emph{nominal substitution via gesture} --- and
hears, under his own clinical shorthand, something he recognizes and
does not want to. His mother, at the end of a telephone line, describing
a hospital appointment without once using the word for the thing being
appointed. The shape drawn in the air around the absent noun. He had
thought, then, that she was protecting him. He wonders now whether she
was protecting the word, keeping it from fastening to her by refusing it
a nominative. He did not have the grammar for that suspicion until this
coast handed it to him.

``You use a specific weave?'' he asks.

Arnar looks at him. His eyes are pale, rimmed red with sleeplessness.
``The pattern.'' He makes a chopping motion, distinct from the weaving.
``To stop the run.''

They walk north along the shore, away from the quay. The tide is low but
the rock is wet far above the line, black testimony to how high the
water reached in the dark. Arnar keeps to the landward side, his body
between Rowan and the drop. They pass an iron ringbolt rusted the colour
of dried blood. Arnar touches it --- not with affection, with
verification. ``The mooring.'' He looks up the cliff, away from the bay.
``The winter. Fifty years.'' He traces a horizontal line at hip height,
then drops it to his knee. ``The level.''

Rowan writes \emph{subsidence record} and thinks: the man is a living
archive, and the index is encrypted, and he has met this filing system
before, in a woman who kept her whole history in a drawer she would not
open in front of her son.

``Arnar,'' he ventures, stepping over a fissure. ``Thorsteinn mentioned
your wife. That she knew the older weather chants.''

Arnar's stride halts mid-step, then completes with a heavy, deliberate
placement of the heel. He turns --- not to face Rowan, but ninety
degrees to the south, putting his shoulder to the sea and his face to
the granite wall of the cliff. He stares at a patch of yellow lichen.
``She always---'' The sentence hangs. The wind picks at the loose
threads of his sweater. Rowan waits, expecting a memory. \emph{She
always loved the rain. She always sang.}

``The nets would---'' Arnar tries again. He lifts his hands, palms open,
shaping a wide encompassing circle that collapses inward.

``She mended them?''

Arnar closes his hands into fists and drops them. ``The morning,'' he
says. ``The coffee.''

He refuses the possessive. He refuses the past-tense verb. To say
\emph{she drank} would be to place her in time, and time implies an end.
To say \emph{my wife} would be to claim her, and claims are audible to
the water at his back. Rowan reads the side-on stance as a trauma
response, a physical aversion to the environment that took her, and
writes \emph{grief-induced aphasia; emotional load prevents syntactic
completion} --- and then his pen stops, because the phrase is a lie of
the exact kind he specializes in, and he knows it. This is not aphasia.
This is precision. Arnar can say her. He is choosing, syllable by
syllable, not to, because he has calculated what a name costs here.
Rowan has spent his life mistaking that choice for an inability, in this
man, in himself, in the woman who raised him inside a silence he called,
on his grant forms, \emph{the reticence characteristic of her
generation.}

``It's alright,'' Rowan says softly. ``We don't have to speak of her.''

Arnar traces the shape of a flower in the air, then erases it with a
flat palm. ``The voice,'' he whispers. ``Sometimes. The glass.'' He
turns back to the path, shoulder locked as a barrier, and nods --- not
points, a chin-lift, a finger might read as a vector --- at a jagged
spire of rock fifty meters out. ``The break. The water.'' He spirals his
index finger, a vortex in the empty air.

``The currents are dangerous there.''

``The water,'' Arnar repeats, firm, correcting the attribution of
danger. ``Does not---'' He stops. He swallows the verb. \emph{Does not
kill. Does not forgive.} He walks on. The silence he leaves is not
empty. It is load-bearing. It holds the cliff up; it keeps the tide from
rushing into the gap where the sentence should have ended. Rowan
follows, recording the sound of their boots on stone, and does not yet
know that the rhythm of their footsteps is the only grammar keeping them
dry.

\sceneline

They go out at six the next morning, because Arnar is going out anyway
and does not object when Rowan appears at the quay, which is as close to
an invitation as the man has in him.

The boat is a nine-metre plastic thing with a wheelhouse the size of a
telephone box, and it is called \emph{Vonin}, the hope, and Arnar has
never once said the name aloud in Rowan's hearing; when he refers to it
at all he says \emph{the boat.} The season's quota is on a laminated
card taped inside the wheelhouse door, and the figure on it is, Rowan
gathers over the next four hours, an outrage, a personal insult, and a
subject on which the harbour master of Brynjavík holds views of some
duration.

He gets more words out of Arnar in the first hour at sea than in the
previous three days on land.

Not about the water. About the system. About the quota being tradeable,
which means saleable, which means gone. About the four boats that were
here in 1994 and the one that is here now. About a man in the south who
owns the right to catch fish in this bay and has never seen it and never
will, and about how that is entirely legal, and is not theft in any
sense a court would recognise, and is.

``Half the coast is the same,'' Arnar says, adjusting the throttle. ``It
was not our weather that took the boats.''

It is the most sustained piece of ordinary political grievance Rowan has
heard in this village, and it arrives in complete sentences, with
subjects and active verbs and no hedging whatsoever, because the
Ministry of Fisheries is not listening from below.

Then he cuts the engine on the drift, and the talking stops, and does
not start again for two hours.

Rowan watches him read the weather. There is no instrument involved.
Arnar looks at the water four times in twenty minutes --- not at the
sky, at the water --- and on the fourth look he starts the engine and
turns for home, and at that moment the bay is flat and the sky is the
same nothing it has been all morning.

The squall arrives forty minutes later, when they are tied up and Rowan
has both hands round a mug in the shed.

``How?''

Arnar coils a line. ``The chop was wrong for the wind.'' A shrug. ``More
coming in than the wind was making.''

Rowan writes it down. Then he stops, and looks at what he has written,
and crosses out \emph{intuition,} and writes \emph{differential} ---
because that is what it is. A man reading the gap between two
quantities and inferring a third from the size of it. It is not
mysticism and it is not instinct. It is the identical operation Rowan
performs on attested and unattested forms, and Arnar does it on open
water at six in the morning with one hand on a throttle, and has been
doing it since he was a boy, and would not have a word for it if asked.

\sceneline

Arnar's cottage is an exercise in subtraction. The exterior is rough
basalt, indistinguishable from the scree; the interior is planed smooth,
scrubbed to the grain. The heavy door seals with the finality of a
vault. The room smells of scalded milk and dried peat. Rowan stands
dripping on the mat, cataloguing. It is furnished for one, yet the
geometry insists on two. The table is pushed to the wall, one chair
tucked in, the other rigid in the corner holding a stack of charts.
Above the hearth, two hooks: one holds a damp oilskin; the other is
empty, the iron dark against whitewashed stone. He looks at the empty
hook a beat too long. He has an empty hook of his own --- a coat still
on the back of a door in Grantchester that he could not make himself
give away, that his rational mind knows is only wool and cannot possibly
still hold her shape, and does.

Arnar does not ask what he wants. He lifts the kettle --- the gesture is
the noun --- and sets it on the burner, takes down two mugs. Clink.
Clink. ``The cold,'' he says, and taps a mug.

``Thank you.'' Rowan sits. The room is not just quiet; it is sanitized
of pronouns. No photographs. No knick-knacks. Only tools.

``I've been studying the village records,'' he says, opening to a fresh
page. ``The census shows decline, but the physical footprint seems
fluid. Buildings on the old maps that aren't here, and the reverse.''

Arnar watches the steam. He stands with his hip to the counter, blocking
the line of sight between the table and the window. ``Before the
last---'' He stops. He makes a downward sweeping gesture. \emph{Storm?
Surge?} He leaves it open.

``Before the event,'' Rowan supplies.

Arnar nods, a concession. ``Houses stood.''

``Your houses?''

``Houses,'' Arnar corrects. He looks at the floor. ``The foundation. The
clay.'' He locks his fingers and pulls them slowly apart. ``The
ground---''

``Subsidence. The ground gave way.''

Arnar shakes his head. He picks up a length of twine, holds it up. ``The
ground,'' he says, and twists it into a loop, ``did not hold.'' He
refuses the active verb. The ground did not collapse; it failed to
maintain the state of holding. A subtle distinction --- to Arnar, the
difference between a tragedy and a sentence.

His fingers move with terrifying precision, looping the end over the
standing part. ``For the wind,'' he says, tightening. A modified
bowline, the loop doubled. ``When the air---'' He gestures at the
window. ``When the air pulls.'' He ties another, a spiral of hitches
around the rope's core. ``For the net. When the weight---'' He leaves it
dangling. \emph{When the weight is a body.}

He does not stop there.

He ties eleven, and he ties them in an order, and it takes about twenty
minutes, and it is the longest sustained act of communication Rowan has
had from him. There is no commentary. Each knot goes onto the table in
sequence, left to right, and where two are related the second is set
slightly below the first, so that what accumulates on the scrubbed pine
is not a row but a shape --- a branching, with a trunk and two limbs and
a place near the end where three of them cluster very close together.

Rowan watches the shape arrive and feels the back of his neck go cold,
because he has drawn this. Not this object; this diagram. He has drawn
it forty times on whiteboards in three countries. It is a stemma. It is
what a morphologist draws when they are showing which forms descend from
which, and where the branches split, and which cousins sit so close
together that students confuse them.

``Arnar.''

``Mm.''

``Who taught you the order?''

Arnar's hands do not stop. ``The order is the order.''

``No --- I mean, when you were taught. Was it one man? Over how long?''

A pause of three or four seconds while a hitch is set.

``My father,'' Arnar says. ``Eleven winters. One a winter.'' He seats the
knot and looks at it. ``You do not get the next one until the last one
holds in weather.''

\emph{One a winter.} Rowan does the arithmetic without meaning to: a boy
starting at seven has the full set at eighteen, which is when he would
first be put in a boat alone, which is precisely how long a
morphological paradigm takes to become automatic in a first-language
speaker. This is not a folk pedagogy. It is a syllabus, and it has been
costed, and somebody a long time ago worked out exactly how much a child
can hold per year and then held to it for four hundred years.

``And if a boy is quick?'' he asks. ``If he could take two in a winter?''

``Then he waits.'' Arnar puts down the twine. ``The quick ones are the
ones you lose.''

Rowan leans forward. He recognizes the spiral; it matches the diagram he
drew of the Sæfangi, the encirclement. ``You're mapping the tension. The
knot is a grammatical modifier. It distributes the load.''

Arnar looks at him. ``The knot stops the slip.'' He sets the knotted
twine on the table between them like a specimen.

Rowan glances at the window. The stove's heat has begun to work against
the chill of the glass, and fog gathers in the corners and creeps inward
--- not uniformly. It streaks. The condensation forms short horizontal
dashes. Hyphens. A page of text in which every word has been redacted. A
single droplet, heavy, rolls down the pane tracing a shape too much like
a comma, a pause where none should be.

``The insulation's poor,'' Rowan notes. ``You get a lot of buildup.''

Arnar turns and sees the glass. He does not reach for a rag. He crosses
to the window with the measured step of a priest approaching an altar
and lifts his hand, palm flat and dry. He wipes the glass --- does not
scrub, erases. One long horizontal stroke, left to right. Then another
below it. He is completing the lines the moisture is trying to write; he
is finishing the redactions himself. He is censoring the window. ``The
view,'' Arnar says, his hand pressed to the cold pane. ``Needs---''
\emph{Clearing? Closing?} ``Needs salt.'' He turns back, wipes his damp
hand on his trousers.

Rowan does not write \emph{compulsive cleaning ritual.} He watches a man
he has begun, against every rule of his profession, to love a little,
hold a page of the sea's handwriting flat under his palm until it gives
up, and he understands that this is what the whole village is: a hand
held against the glass, erasing faster than the water can write. He has
done it himself. He did it over his mother's grave, where he said
nothing, and called the nothing a kindness.

Arnar pours the water and sets the pot down but does not pour the tea.
He waits, looking at the empty hook. ``The tea,'' he says finally.
``Wait.'' Rowan realizes he is waiting for the boil to subside, for the
water to go passive before he will let it be consumed.

``You're very precise, Arnar. With your words. With your hands.''

Arnar touches the knot. ``The sea,'' he says, barely audible. ``Does
not---''

``Does not what?''

Arnar looks at the empty chair. ``Forget.'' He says it without a
subject. Just the verb. He pours the tea, and the sound of the liquid
hitting the ceramic is loud, and it sounds like a sentence ending.

\sceneline

He learns about the wife in pieces, across that week, and never once
from Arnar.

Her name is on the harbour rota from 2011, still pinned in the shed
because nobody has taken it down: \emph{Kristbjörg,} in the Saturday
column, eleven Saturdays, in a round unhurried hand.

She is in a photograph in the community hall, in a group of fourteen at
the opening of the slipway, second from the left, laughing at something
outside the frame.

The mugs in the cottage are a set of six, and four of them are chipped
in the same place, which means they have been washed by the same hands
in the same sink in the same order for a very long time, and two of them
are not chipped at all.

And there are the two hooks by the hearth.

Rowan noticed the empty one on his first visit and read it, correctly,
as the obvious thing. What he notices on his third is that Arnar hangs
his own oilskin on the \emph{left} hook, every time --- and that the left
hook is the bad one. It sits behind the door swing. It is four inches
lower than it should be. A man reaching for it has to lean and turn his
shoulder, and Arnar leans and turns his shoulder, every day, twice a
day, for however many years it has been.

The right-hand hook is at shoulder height, perfectly placed, and empty.

He is not preserving it. Preservation would be sentiment, and there is
no sentiment in this room --- no photograph, no keepsake, nothing that
could be called a shrine, nothing an outsider could point at. He is
simply not using it. Day after day, deliberately, at a small daily cost
to his own back, the way a man does not use a word.

And Rowan, who has stood in this room before and looked at the empty
hook and thought about a navy coat on a door nine hundred miles south,
revises the whole reading, because he had it sentimental and it is not
sentimental. He had been telling himself for eight months that what he
could not do was give the coat away. That is not what he could not do.
He could not take it down and hang his own there.

It was never grief. It was the same operation. Leave the hook empty,
leave the noun out of the sentence, and pay for it yourself, twice a
day, quietly, forever, rather than let the thing be handled.

He does not mention the hook. He hangs his coat over the back of a
chair.

Arnar sees him do it. Something in the man's face moves, once, and
settles.

``The chair is fine,'' he says.

\sceneline

The walk back to the guest house is a navigation of small accumulated
errors. Rowan measures his progress not by distance but by discrepancy.
The village has not changed, yet it does not fit the map in his head.
The tide line on the harbour wall, two meters below the coping
yesterday, now crusts the concrete a half-meter from the top, dry,
white, as if the water held that height for hours unseen. At the lower
path a volcanic boulder, pocked with gas bubbles, sits in the center of
the gravel track. He stepped around this rock yesterday, on the left
verge. Now it blocks the way, having moved three meters inland without a
furrow behind it.

He looks back. Arnar stands in his doorway, a silhouette in the dark
stone, watching Rowan, or watching the space between Rowan and the
water. ``Arnar,'' Rowan calls, pointing. ``This has moved.''

Arnar leans against the frame, half in shadow. ``The location,'' he
calls back, carried effortlessly by the damp air, ``is not fixed.''

``But the rock. It's heavy. It implies force.''

``The sea,'' Arnar says. He looks up, checking the grammar of the
clouds. ``Remembers what we---'' He stops. He closes the door. The latch
is a period at the end of a paragraph Rowan has not finished reading.
\emph{Remembers what we forgot? What we took?}

He climbs to the Annex. The wind is low in the dried grass, sibilant.
Sss. Sss. Inside, he does not light the lamp. He stands in the dark,
breathing stale air, the room pressurized. He is bone-tired, the fatigue
of a man who has been treading water for days. He sleeps in his clothes.

The dream is not a narrative. It is a diagram. He hovers above the
coastline. Below him the ocean is not water but syntax. The waves are
long rolling clauses that gather from the deep and approach the shore in
the nominative --- subjects seeking an object. He watches a wave curl:
the crest is the verb, propelling the mass; as it breaks it declines,
giving its energy to the land in the dative, the foam hissing up the
beach in spiral genitive markers, possessive loops that try to circle
the rocks and claim them. \emph{Of the stone. Of the sand. Of the man.}
The tide withdraws in the passive. \emph{The sand was taken. The stone
was eroded.} He sees the pressure gradients as lines of force that
answer to the intensity of grief; where the water is deepest the grammar
is most compressed, dense with unsaid nouns.

He wakes to the grey luminescence of pre-dawn fog against the glass. The
silence is absolute, yet the room is vibrating. He sits up and looks at
the window.

The glass is opaque. The condensation is thick, white, organized --- not
random moisture but rows of jagged angular script, lines that start
strong and trail into nothing. Hooks without bait. Stems without
flowers. \emph{Th. Wh. Sh.} He crosses to it, floorboards cold under his
socks. The patterns are intricate. He sees the shape of Arnar's knot
repeated a hundred times in the droplets. He sees the spiral of the
genitive. He sees the sharp vertical stroke of the command. He raises
his hand --- not for the camera, not for the notebook. The need to
document has been replaced by the need to answer.

He presses his fingertips to the glass. The vibration travels up his
arm, low and rhythmic, not the wind, not a truck on the road. A cadence.
Da-dum. Da-dum. Da. The rhythm of a heart, or of a sentence waiting for
a full stop that never comes. The glass is singing the part of the
conversation Arnar refused to speak. Rowan stands with his hand on the
cold wet text, feeling the pulse of the sea against the thin barrier of
the room, and understands. He does not know the words yet. He recognizes
the syntax. It is the language he was raised in. It is waiting for him
to answer, and he has been not-answering it his whole life.

\chapter{}

The schoolhouse sits on a rise the village has quietly receded from,
stranded in withered heather. It is not a ruin. Ruins imply a surrender
to nature, a softening of lines. This building is a rejection --- black
basalt blocks squared with punishing precision, hunched against the
skyline, its slate roof pulled low like the brim of a hat worn by a man
who does not wish to be recognized.

Rowan climbs, boots slipping on wet scree. The wind here is different.
Down in the village the gale shoves, physical and boisterous. Here the
air streams around the stone corners with a thin high whistle that
sounds disturbingly like inhalation. Sigrún walks two paces behind, a
scarf of grey wool swallowing the lower half of her face, leaving only
her eyes, rimmed red, moving from the clouded horizon to the heavy oak
door. She does not look at the windows, whose glass is mostly gone,
leaving dark rectangular sockets that observe their approach with
bureaucratic indifference.

``The key,'' Rowan says at the threshold.

She fumbles in her coat with hands that shake at a specific frequency
and produces an iron key, heavy, pitted with orange rust. She holds it
out but does not move to insert it. ``It has been closed since
ninety-four,'' she says, muffled by the wool. ``The Ministry
consolidated the district. They said the heating was unsustainable.''

``And the curriculum?''

``Unsustainable,'' she repeats.

The mechanism grinds --- metal teeth on grit --- but the bolt throws.
The door swings inward on hinges that scream and then settle into a
sullen silence, and they step into the air of 1994. Dry, cold, chalk
dust and dead flies. The silence is not empty; it is suspended. Dust
motes turn in slow hypnotic gyres in the shafts of grey light. Rowan
moves forward, disturbing a layer of silt that has lain unbroken for
thirty years, leaving dark accusatory prints. Sigrún stays near the
door, arms crossed tight, compressing herself into the smallest possible
volume.

``We shouldn't be long,'' she whispers.

He unzips his jacket, the sound tearing the quiet, and moves down the
corridor, walls the peeling institutional green of sea-sickness. He
stops at the first display, a chart pinned with rusted tacks.
\emph{Declensions of the Strong Nouns.} Yellowed, curling like dried
skin. The standard four cases --- Nominative, Accusative, Dative,
Genitive --- faded but legible. But below the Genitive the paper has
been surgically removed. Not torn. Cut, with a blade, perfectly
straight. He moves to the next: \emph{Usage of the Definite Article,}
the bottom third --- where the irregular, ritual forms would be listed
--- excised, the wall behind it a brighter green, a phantom limb.

``Sigrún,'' he says, not turning. ``Who cut these?''

A pause. The wind finds a gap in the roof tiles and moans, dropping an
octave, a vowel trapped in a throat. ``The inspection board,'' she
answers from the shadows. ``In eighty-two. They said the walls were
cluttered. They said the charts should reflect the examinable
material.''

He touches the edge of the paper. ``They removed the Sæfangi. And the
Hrimfall.''

``They removed the clutter.'' Flat, rehearsed.

Rowan walks deeper into the hall. The amputation is systemic. A painted
mural of the local fauna --- puffins, arctic foxes, seals --- has been
altered: the seals painted over with grey blocks, leaving only the rocks
they rested on. A rule for the address of weather phenomena has been
whitewashed, the paint thick and hurried, a bandage over a wound that
would not stop. He stops outside a classroom door, its glass pane
shattered, the shards inside. The wind through the gap is making a
resonance, the draught catching the jagged edge and producing a sibilant
hiss that modulates with the gusts. \emph{Sssss. Uuu. Thhh.} It sounds
like a child sounding out a word they are forbidden to speak. On the
floor the dust is disturbed --- not by feet, by the draught, swirled
into perfect tight spirals, the genitive loops of his dream. The
building is breathing, and in its exhalation trying to decline the air.

``This isn't a ruin,'' he murmurs, photographing the sliced chart.
``It's a scoured field.''

At the end of the corridor there is one more.

He almost walks past it. It hangs where the light is worst, between the
cloakroom pegs and a fire notice from 1988, and it is a chart of the
same species as the others --- ruled columns, a heading in the same
institutional hand, the paper gone the same nicotine colour.

The bottom third has been removed.

He puts his thumb on the cut and stops, because the cut is wrong.

The others are thirty years old. Their cut edges have aged with the
rest of the sheet: the same yellowing, the same brittleness, a faint
brown line of oxidation along the exposed fibre where three decades of
damp air has worked at it. This edge is clean. The exposed fibre is
paler than the sheet around it. The paper has been opened recently
enough that it has not yet caught up with itself.

He gets the loupe out and puts the light of his phone on it at a raking
angle, and the fibres along the edge are crisp and standing.

Not 1982. Not 1994. Two years, perhaps three.

Someone has been in this building --- which is locked, which has one key,
which Sigrún carried up the hill in a shaking hand this morning and had
to be asked twice for --- and has stood in the dark at the end of this
corridor and cut a section out of a chart on a wall in a school that has
had no children in it for thirty-one years, and has taken the piece
away, and has not mentioned it.

He photographs the cut. He photographs the wall behind it, which is a
brighter green.

He does not call out to Sigrún, and he is careful, walking back down the
corridor, to arrange his face.

Because there are two readings and he cannot get between them. Either
somebody in this village is still, quietly, thirty years after the
Ministry finished, removing the dangerous parts of the grammar from
where a person might see them --- or somebody came up here in the dark
and cut out the one panel that still had a form on it that nobody living
could otherwise reproduce, and took it home, and put it somewhere safe.

Suppression or rescue. The same blade, the same wall, the same two
inches of missing paper.

``Dr.~Hale,'' Sigrún calls, her voice cracking. ``The office is to the
left. Please.'' Her silhouette against the daylight of the open door
looks terrifyingly small, a figure in the mouth of a cave she knows
holds a bear.

\sceneline

The office smells of mimeograph ink and damp wool, dominated by a desk
of oak veneer peeling back to reveal the particle board beneath. On the
surface, preserved in a stasis of neglect, sit the Skólanámskrá --- the
curriculum ledgers, bound in black oilcloth, the years painted on the
spines in white correction fluid. He approaches as if it were an altar.

He pulls 1965. The hand is copperplate, severe and beautiful. \emph{Week
42: Dative Declensions of Coastal Nouns. Focus on current and drift.
Week 43: The Aflending (Unbinding). Usage in calm weather only.} The
grammar is complete, the lessons practical and integrated, the teacher's
notes mundane. \emph{Jónas struggles with the vocal length. Gudrún is
excellent.}

He reaches for 1978. The hand has changed --- hurried, slanted.
\emph{Week 42: Common Coastal Terms. Standard Icelandic Forms. Note from
Ministry: Ritual cases designated Archaic/Folklore. To be moved to
supplementary reading. Not for examination.} A red pen has struck out
\emph{Aflending.} In the margin, a different hand: \emph{Hard to explain
to parents. Dropped.}

He opens 1982. Stiffer pages, standardized curriculum. \emph{Week 12:
Modern Icelandic Syntax. The Passive Voice. Note: Remove reference to
Sæfangi. Confuses the students regarding subject-object relationships
with inanimate nature.} And then the margins begin to speak. In 1965
they were blank. In 1982 they are crowded. \emph{Oct 14: Windows rattled
all lecture. Hard to hear. Nov 2: Roof leak in north corridor. Wind is
loud. Nov 10: Fog inside the vestibule again.}

He arranges the books in a chronological arc across the dusty desk and
pulls the last, 1991. The curriculum is clean, modern, safe; Vikamál
reduced to a single hour of ``Local Heritage,'' the grammar drills gone.
And the margins are screaming. \emph{Sept 4: Storm took the playground
fence. Sept 12: Damp is ruining the books. The ink runs. Oct 1: Cannot
keep the classroom warm. The draught speaks. Oct 20: Stopped the lesson.
The noise was too much.} He runs his finger down the page. The
correlation is linear. The week the Hrimfall was struck from the plan
--- a heavy black \emph{OMIT} --- the margin note three days later
reads: \emph{Boiler failure. Frost inside the walls. Pipes burst.}

``It tracks,'' he whispers. ``The erosion isn't only cultural. It's
physical.''

Sigrún grips the doorframe, white-knuckled, her other hand picking at a
loose thread on her cuff, a darker grey than the wool, unravelling with
each twitch. ``They wanted modernization,'' she says, thin and reedy.
``They wanted the children to get jobs in Reykjavík. To speak like
people on the television.''

His finger rests on an entry from 1988. \emph{Sea water in the basement.
Source unknown.} ``Sigrún. Look at the dates. Every time a form was
removed, the building took damage. Did no one see it?''

She looks at the ledger, then away, to a calendar still turned to
November 1994. ``We saw the weather. We thought it was the climate
changing. We thought it was bad luck.''

``But it's one-to-one. Here. You stopped teaching the containment forms
for the tide in March eighty-nine. The quay flooded in April.''

``It was a wet spring.'' She pulls the thread. It snaps.

Rowan hears the academic in his own skull reciting its catechism.
\emph{Correlation is not causation.} The teachers, stripped of a
cultural identity, projected anxiety onto the weather; they noticed the
storms because they felt exposed; a psychosomatic loop. He has spent his
career building rooms like that, rooms with a locked interpretation, so
that nothing gets in that he has not first labelled. But his finger
rests on a note from 1993, in a shaking hand: \emph{I can no longer
teach the words. When I say them, the windows crack.} He closes the
book. The dust puffs up in a small dry cloud.

``Let's look at the classroom,'' he says.

Sigrún does not move. She stares at the desk as if the books are
radioactive. ``Yes,'' she says. ``The classroom is where we finished
it.''

\sceneline

The classroom is a box of grey light. The west windows face the bay but
the salt spray has rendered them a cataract. Rows of desks sit in strict
formation, laminate scarred by generations of bored carving ---
initials, hearts, crude ships. Rowan sits at the teacher's desk on its
small dais; Sigrún takes a student chair in the front row. The
difference in height makes him uncomfortable, re-enacting a dynamic he
does not want. Behind him the chalkboard is a landscape of erased
equations and sentence diagrams, faint white ghosts of letters clinging
to the slate, wiped but not gone.

``The ledger said you were head teacher in the final years.''

She has unwound her scarf. Her face is pale, the lines around her mouth
cut deep. She looks not at him but at the blank of the blackboard. ``I
signed the orders.''

``The Ministry sent the directives. You had to comply.''

``The Ministry sent paper.'' She clasps her hands in her lap. ``I did
the teaching.''

The wind thuds against the glass. ``Why did you stop? Really. You knew
the forms. You knew their history. You could have kept them as oral
tradition, off the exam.''

She laughs, dry and brittle. ``History. Tradition.'' She shakes her
head. ``You think of it as heritage, Dr.~Hale. Something to put in a
case in a museum.'' She looks at the empty chair beside her. ``They
asked me,'' she whispers.

``The Ministry?''

``The children. Twelve years old. Sharp. Modern. Satellite television.
Science textbooks that explained the tides were the moon's gravity, not
the breath of a giant.'' She leans forward, her eyes wet but not
weeping, reliving a specific defeat. ``I would put the Sæfangi on the
board. \emph{Hafinu} --- to the sea. \emph{Hafsins} --- of the sea. And
then the ritual form, \emph{Sævangur.} And a hand would go up. `Teacher,
why do we use that one? It doesn't follow the rules. It makes the
sentence longer. What does it do?'\,'' She spreads her hands. ``What
could I tell them? Could I say, `If you don't use this suffix, the water
will come into your bedroom'? They would laugh. They would tell their
parents I was teaching ghost stories. I was a modernized educator,
Rowan. I had a degree from Reykjavík.''

She looks at her hands. ``So I lied. I told them it was just an old way
of speaking. Poetic. Optional.''

``And then?''

``And then they asked, `If it's optional, why do we have to learn it?
It's hard. It's boring.'\,'' She slumps, the energy leaving her. ``And I
didn't have an answer. Not a modern one. So I stopped. I stopped because
I was tired of trying to explain a lock they couldn't see to children
who only believed in keys they could hold.''

``Was it one child?'' Rowan asks. ``Or the room.''

Sigrún is quiet long enough that he thinks she has not heard.

``Ásdís Helgadóttir,'' she says. ``Second row, by the radiator, because
she felt the cold.''

She says the name the way one produces a document from a drawer one
visits often.

``She was not a difficult girl. That is the thing people always assume
--- that it was some clever awful child who liked to catch teachers out.
She was kind. She was a serious child, and serious children ask the
serious question, and the serious question was: \emph{Miss, if it worked,
wouldn't somebody outside know?}''

The wind pushes at the window.

``I want you to understand how good that question is,'' Sigrún says. ``I
have had thirty-one years to find a flaw in it and there is not one. If
the forms hold the water back, then Reykjavík would know. The university
would know. Somebody would have measured it. There would be a paper, and
a grant, and men in coats.'' A short breath. ``And I stood in front of
twenty-two children and I could not tell her the true answer, which is
that they \emph{did} know, once, in six places on this coast, and that
five of them stopped, and that we have all the evidence anyone could
want and it is a hundred and fifty years of nothing happening here and
things happening elsewhere, and that no method on earth can turn that
into a proof.''

``What did you tell her?''

``I told her it was a very good question.'' Sigrún's mouth twists.
``Which is what a teacher says when she has lost.''

``Where is she now? Ásdís.''

``Reykjavík. She is a radiographer.'' The twist becomes almost a smile,
and it is terrible. ``She comes at Easter. She is a lovely woman. She
brings her children and they run about on the slipway and she has no
idea, none at all, that she is the reason, and I would sooner cut my
hand off than tell her.'' Her hands tighten in her lap. ``It was not her
fault. She was eleven. She asked the only question there was.''

Rowan stares at her, and it is the banality of it that lands. The
defence of the village was not dismantled by a cult or a monster. It was
dismantled by pedagogical convenience, by the exhaustion of a teacher
facing a room of bored children asking \emph{is this on the test.} And
under that, quieter, a recognition that turns his stomach: this is
precisely how his mother's history was lost. Not to catastrophe. To the
small daily refusal to explain a thing that cost too much to say
plainly. She, too, had decided that some grammar was better dropped than
defended in front of a child who would only ask what it was for.

``You removed the grammar,'' he says, ``to save yourself the
embarrassment of explaining it.''

``I removed it,'' she says, ``because I couldn't bear to teach them a
language I wasn't allowed to tell the truth about.''

She looks up at the chalkboard, and for a moment Rowan thinks she sees
something written there. ``And now,'' she whispers, ``the weather asks
the questions. And no one knows the answers.'' He wants to tell her it
was not her fault. But the deficit is absolute, and he has spent his
life learning that some silences are the only honest response, and he
keeps this one. The only sound in the room is the wind, which is no
longer whistling. It is listening.

\sceneline

They exit into a silence that feels stunned, the wind dropped, a vacuum
pressing on the eardrums, the grey afternoon flattened of depth. He
feels a desperate need to catalogue the confession before it evaporates,
and reaches for his notebook.

``Dr.~Hale.'' Sharp. She is five meters off, staring at the ground near
the north wall's foundation. There is a pool of water there. When they
entered, the ground was dry --- frozen scree, dead heather. Now a
depression is filled with dark glassy water, a perfect circle a meter
across, sitting in the earth like a lens dropped from a telescope.

``Where did this come from?'' He crouches. ``A burst pipe?''

``Don't touch it,'' Sigrún says, backing away.

He leans over. The surface is unnervingly still --- no tremor, no
ripple. It reflects the grey sky, the basalt wall, his own face. He
stares at his reflection. The reflection stares back. Then the
reflection blinks. Rowan has not blinked. He freezes. The face in the
pool shifts, the mouth opening slightly as if about to speak, the water
perfectly flat while the image inside it moves independently of the
source. The reflection's eyes drift left, toward the village. Vertigo,
the earth tilting. The pool is not reflecting him; it is observing him.
On the surface a film of oil or dust coalesces into lines, curves ---
genitive markers. The water is trying to spell \emph{Min---}

``Sigrún?'' He turns. She is walking away, rigid, frantic, hands pressed
over her ears. As she passes the corner she speaks --- not a scream, a
mutter into her collar. \emph{``Vatnið heldur\ldots{} vatnið
þegir\ldots{} engin orð\ldots{} engin nöfn---''} \emph{The water holds.
The water is silent. No words. No names.}

The syllables hit him like a pressure change in a descending cabin. His
ears pop. A high whine drills into his skull, instant and painful. The
air around the pool vibrates, blurring the sight. He looks back and for
a second the reflection is gone; the pool is just dirty water in a hole.
But the silence has broken. The sound of the sea, distant and heavy,
comes crashing back into the world. He scrambles up, backing away from
the circle, heart hammering, and runs after her, leaving the schoolhouse
and its new unblinking eye to watch the empty moor.

He catches her at the foot of the path. She is not running now. She has
stopped at the cairn marking the schoolhouse boundary --- a waist-high
stack of basalt, each stone placed by some predecessor whose name is not
on it. She is breathing through her teeth, her hands shaking so the
keychain rattles in her pocket.

``Give me your notebook,'' she says.

He hesitates --- Cambridge instinct, the document is the proof --- then
hands it over. She does not open it. She turns to the page she knows is
there: his transcript of her confession, the columns of dates and
amputated suffixes, the line where he wrote her name at the top.
\emph{Sigrún Ásgeirsdóttir, head teacher, 1989--94.} Below it the dotted
line he prepared for her signature. She holds the pen out, and for a
moment he thinks she means to sign. She does not. She tears the page out
cleanly, along the perforation. She does not crumple it; that would
still be an object. She folds it twice, kneels at the cairn, slides it
into the gap between two stones at the base, and seats a fist-sized rock
across the gap. The page is gone. No one will read it. The cairn now
holds a name that was almost spoken aloud.

``You don't write it,'' she says. ``Not with my hand on the page. The
eye looks for signatures. Don't feed it mine.'' She hands the notebook
back. The remaining pages feel lighter than they should. ``You should
burn the recorder,'' she adds, calmer now. ``You won't. So at least
don't play it back here.'' She walks the rest of the way down alone,
slower, picking her steps as if the path were something she has just
renegotiated.

\sceneline

The community hall is a bunker of light against the gathering dark.
Inside, the air is thick with wet wool, strong coffee, and the metallic
tang of fresh blood. Yellow lamps cast a jaundiced glow over the long
tables where the day's work is being finished. He enters still shaking
from the schoolhouse, needing normalization, needing to see ordinary
people at ordinary things. He finds them. Twelve villagers stand around
the central table, sorting the catch. Metal bins in a row. Fish sliding
across the gutting table --- cod, haddock, pollack. Knives flashing wet
and efficient. Slit, gut, slide.

He nods to Arnar, who is sharpening a fillet knife with a rhythmic
shhh-shhh against a stone. Arnar nods back, eyes on the blade. ``Good
catch?'' Rowan asks.

``Variable,'' Arnar says. Not \emph{the sea gave us a lot.} ``The volume
is high.''

A woman at the end of the table picks up a haddock and turns it over.
Rowan stares. The fish has two dorsal fins --- not side by side,
stacked, one perfectly replicating the other, as if the noun had been
pluralized by mistake. The woman does not gasp. She clicks her tongue
--- mild annoyance, like finding a bruise on an apple --- and tosses it
into a red bin marked with a black X. ``Too many doubles today. The
count is off.''

He looks into the red bin. It is a charnel house of grammatical errors.
A cod whose scales run backward, tail to head. A flounder gone perfectly
spherical, swollen with fluid. And on top, a large grey pollock, its
mouth open. He grips the edge of the table. The teeth in its jaw are not
needle-sharp piscine teeth. They are flat, white, rectangular. Human
molars.

Bile rises in his throat. ``That---'' He points, his finger trembling.
``That fish has human teeth.''

The table does not stop. A knife comes down, severing a head. Thwack. An
older fisherman, his face a map of broken capillaries, wipes a bloody
hand on his apron and looks at the fish, then at Rowan. ``The sea gives
back what it remembers.'' He says it the way one says \emph{the postman
brings the mail.} A statement of logistics.

``But that's---it's a biological impossibility. It's mutation. You have
to report this. The environmental agency---''

``We sort,'' the woman interrupts, lifting the pollock by the tail
without looking at the teeth. ``We burn the ones that remember too much.
We eat the ones that are quiet.'' She drops it into a smaller bucket by
the door with a heavy dull thud.

No one is horrified. They are resigned, checking gills, eyes, fins,
quality-controlling a reality that is coming apart at the seams. He
listens to their low talk. \emph{The west side brought the heavy ones.
The drift was full of echoes. This one is mis-set. Burn it.} They never
say \emph{caught.} They never say \emph{ocean.} They never say
\emph{monsters.} The horror is not in the bin. The horror is the lack of
surprise. The village has adapted. Without the grammar to hold the
boundary, the definitions of species have gone porous, and they are
filtering the sea's nightmares out of their food supply, one fish at a
time.

He steps back, away from the blood and brine, and watches Arnar slide a
knife into a cod, opening it to see whether the inside is written in the
language of fish or in something else. And Rowan knows, with a quiet
sickening certainty, that he can no longer only observe. The schoolhouse
is open. The sea is sending its memories back to be eaten. And somewhere
under all of it, unwatched, is the recognition that he did not come here
to study any of this --- that he came carrying a sealed envelope and a
photograph of a black coast, and that the debt he is standing in the
middle of has his own name somewhere on the ledger.

\chapter{}

The list goes up on the hall wall at six in the morning and Rowan's name
is on it.

He has come early because he assumed there would be a beginning --- a
gathering, an announcement, some point at which the day is declared to
have started --- and there is not. There is a hall with the lights on
and the urns already going and forty people in work clothes reading a
sheet of A4, and the sheet has been typed on a computer, badly, in a
font somebody chose in about 2004, and it is headed \emph{FJÓRÐI
LAUGARDAGUR} and then the date.

Two columns. Names on the left. Tasks on the right.

\emph{Hale, R. --- north wall w/ Ingi \& Steinunn.}

He reads it three times. Nobody has mentioned this to him. Nobody has
asked whether he is free, or whether he would like to help, or whether
a visiting academic would prefer to observe; he has simply been put on
the list between Guðjón and Halla, in alphabetical order, in the same
font as everyone else. He is aware of a feeling in his chest that he
does not immediately identify, and it takes him most of the walk to the
north wall to recognise it as the specific warm foolish pleasure of
being picked for a team.

He does not, at any point that day, examine what it means that he did
not have to be asked.

\sceneline

The sisters are at opposite ends of the village.

He works this out at about half past seven, in the way one works out a
piece of local physics, by noticing an absence. Guðrún Sigurjónsdóttir
is on the north wall with him; Kristín Sigurjónsdóttir is, per the
sheet, on the roof detail at the curing house, which is nine hundred
metres away and around a headland. He mentions it to Steinunn, who is
seventy and has hands like a bundle of dry sticks and who is currently
outworking him three stones to one.

``Mm,'' says Steinunn.

``Is that deliberate?''

``It's the list.'' Steinunn seats a stone, sights along the course, and
takes it out again. ``It has been the list since two thousand and
nine.''

``What happened in two thousand and nine?''

``A mooring.'' She sets the stone back, differently. ``I do not have an
opinion.'' A pause of perhaps four seconds. ``Kristín was wrong.''

Rowan laughs before he can stop himself, and Steinunn's face does not
change at all, and he understands that he has been given something.

\sceneline

The north wall is a dry-stone field boundary a hundred and forty metres
long and it is, Rowan discovers over the following four hours, an
absolute bastard.

There is no mortar. There is no cheating. Each stone has one orientation
in which it will sit and forty in which it will not, and the difference
between a course that will stand for sixty years and a course that will
come down in March is invisible to him and apparently obvious to
everyone else. Ingi --- forty-ish, enormous, a man built like a door ---
watches him wrestle a block for two minutes and then reaches over and
turns it ninety degrees and it drops into the gap with a small final
click, like a word arriving in the right case.

``Ah,'' says Rowan.

``You will get it,'' Ingi says, which is generous and untrue, and goes
back to his own section.

By nine his back has opinions. By ten his hands have split at the
knuckles in two places and he has stopped noticing. The wind is coming
off the water in long cold pushes and the work is warm and stupid and
entirely absorbing and he has not thought about morphology, or his
mother, or the tin, since about a quarter past seven.

They talk while they work. Not about the wall.

The talk is about quota, which is a disaster; about the price of diesel,
which is a scandal; about a man in Húsavík who sold Ingi a winch; about
whether the school bus should be made to wait the extra four minutes at
the junction, which is apparently a live constitutional question in this
village and on which Steinunn holds views of some force. Somebody's
daughter has got into a programme in Akureyri. Somebody's boat needs a
gearbox. There is a long, detailed, and to Rowan completely opaque
account of a dispute about a fence in 1998 which ends with Ingi saying
``and that is why we do not ask Þórarinn about fences,'' and both women
nodding as at a settled point of law.

At eleven a boy comes along the wall with a thermos and a bag of
kleinur, distributing them to the work parties with the self-importance
of an appointed official, which he is; his name is on the list too, near
the bottom, \emph{Kári --- kaffi, allar sveitir.}

He is nine. He has a cowlick that will not lie down and a card trick he
performs for Rowan while the adults drink, which does not work. He
performs it again more slowly. It does not work again.

``That is the best one I have seen today,'' Rowan says.

``It works at home,'' says Kári, unbothered, and moves on down the wall
with his thermos.

\sceneline

At noon they eat in the hall, all of them, at trestle tables, and it is
the loudest room Rowan has been in since Reykjavík.

There is lamb soup and rye bread and butter and a vat of coffee, and
there are dogs under the tables, and someone has brought a cake because
it is Ólöf Magnúsdóttir's sixtieth and the cake says so in wobbling blue
icing. Ólöf is made to stand up. Ólöf does not want to stand up. Ólöf is
made to stand up anyway, and a version of the birthday song is sung with
what Rowan can only describe as aggressive incompetence, four different
keys and at least two different sets of words, and Ólöf sits down going
red and pleased and immediately begins criticising the cake.

Rowan sits between Ingi and a woman whose name he does not catch and
eats two bowls of soup and is asked, at length and with genuine
curiosity, about English weather, English house prices, and whether it
is true that in England you must pay to see a doctor. He explains the
National Health Service. There is a general and forceful opinion about
America. Someone wants to know if he is married and he says no and the
follow-up question does not come, and he notices the follow-up question
not coming, and is grateful.

Across the hall Sigrún sits with her soup untouched, wrapped in the
grey scarf, watching the room.

Their eyes meet once. She looks away first.

\sceneline

After the meal he is reassigned. The sheet has a second column for the
afternoon and he is now \emph{Hale, R.\ --- afritun}, which he has to
ask about, and which turns out to mean the back room.

Four people sit at a table under a strip light with a stack of paper
between them, copying.

They are copying by hand. There is a photocopier in the corner of the
room with a cover over it, plugged in, working; Rowan checks later, and
it works. What they are doing is transcribing --- from soft old
notebooks and loose sheets and two things that are unmistakably
nineteenth-century --- onto new lined A4, in ballpoint, slowly, checking
each other's lines.

And they are not throwing the old ones away. The old ones go into a
second stack, squared, in order.

``What happens to those?'' Rowan asks.

``They go in the cupboard.''

``With the others?''

``With the others.''

He looks at the cupboard, which is an ordinary steel-doored office
cupboard, and he thinks about how many fourth Saturdays there have been
since 2004, when this photocopier was presumably bought, and about how
many there have been since ballpoint, and paper, and the hall.

``Can I ask why you don't photocopy them?''

He asks it lightly, of the table generally, and he gets four answers.

Guðjón, who is sixty and does not look up: ``The machine does not make a
copy. It makes a picture of a copy.''

The woman opposite Guðjón, without hesitation and slightly impatiently,
as though the question were a bit slow: ``Because then nobody has read
it.''

The young man at the end, who is perhaps twenty-five and whose
handwriting is the neatest at the table: ``I don't know. This is how we
do it.'' He shrugs, entirely without embarrassment. ``It's four hours.''

And Halla Kristjánsdóttir, at the head of the table, who does look up,
and who considers Rowan for a moment over her glasses: ``Because the
hand has to go the same way the hand went.''

Rowan waits. Nobody amends anyone. Nobody says \emph{well, what Halla
means is.} Four mutually incompatible accounts of the same practice have
been laid on the table and all four of them are apparently fine, and the
copying continues, and someone asks whether there is more coffee.

He picks up a pen and is handed a notebook, and copies for two and a
half hours, and finds --- to his considerable irritation --- that the
woman opposite Guðjón was right. He has read every word. He could not
have told you a line of it if he had photographed it.

\begin{fieldnote}
Monthly communal maintenance. Universal participation, allocated in
writing, no exemptions observed --- including for the visiting
researcher, who was rostered without consultation and did not think to
object.

Manual transcription of textual material, retention of superseded
copies. Four informants, four incompatible rationales, no observable
discomfort at the incompatibility. Cf.\ orthodox practice in most
literate ritual traditions, where the \emph{explanation} is precisely the
part that is not standardised, because it is not the part that is being
transmitted.

The practice is the transmission. The reasons are decoration, and the
village knows it, and I am the only person in this room who needed to
ask.
\end{fieldnote}

\sceneline

The boundary stones are last, and they are done in an order, and the
order is not geographical.

Rowan works this out following a party of nine up the eastern slope in
the failing light. They do not walk the perimeter as a perimeter; they
go to the third stone, then the first, then the seventh, then the
fourth, in a sequence that everybody knows and nobody consults a list
for, and when he asks Ingi why, Ingi says ``that is the order,'' which
by now Rowan has learned to hear as a complete answer.

At each stone: clear the growth back a hand's width. Check the set with
a boot. Repack if it has lifted. Somebody says a short thing --- four or
five syllables, flat, unceremonious, in the local grammar, delivered in
the tone of a man reading out a serial number. Then the next stone.

The children come. That is the part Rowan does not expect. The whole
afternoon's children, six or seven of them, are up on the slope in the
cold with the adults, and they are not being minded; they are being
included, in the specific way that means somebody has decided which
child is old enough for which stone.

At the north gatepost, Kári is put in front.

His grandmother stands behind him --- Halla, out of the copying room now
and up the hill in a coat like a tent --- and she puts her hand over his
hand on the cold stone, and says the form once, and waits.

Kári says it. He gets the ending wrong.

``Again.''

He gets it wrong the same way.

Halla does not correct the word. She moves his hand slightly on the
stone, an inch to the left, and says the form again herself, and this
time when Kári repeats it, it is right, and Rowan --- standing four
metres off in the wind with his hands in his armpits --- understands
that he has just watched a case ending taught as a physical location.
Not \emph{say it like this}. \emph{Stand here and it will be like
this.}

``Góður,'' says Halla. Good.

Kári goes scarlet to the ears and turns away fast and pretends to be
extremely interested in the fence.

Rowan writes nothing down. He is not sure, afterwards, whether that was
respect or whether he simply did not want to be a man with a notebook in
that particular minute.

\sceneline

The eleventh stone is on the western spur, and when they reach it they
do the growth and the set and the repacking, and then they stand for a
second, and then they move on.

Rowan is three paces past it before it registers.

He stops. He looks back at the stone, and at the party going on ahead
of him in the dusk, and he counts back through the afternoon --- third,
first, seventh, fourth, ninth, second, and so on, ten stones, ten short
flat unceremonious things said in ten places --- and at the eleventh
nobody said anything at all.

He catches up to Ingi. ``The last one.''

``Mm.''

``Nobody spoke.''

``No.'' Ingi keeps walking.

``Is there --- '' Rowan picks his way, because he has been here long
enough now to know that the wrong question closes a door for a month.
``Does that one not take a form?''

``It takes one.'' Ingi steps over a rill. ``We do the stone.''

``But not the form.''

``Not that one.'' And then, because Ingi is a kind man and can hear that
Rowan is not being a journalist: ``Nobody has it. Þórunn had it. She
died in --- '' he thinks, ``--- ninety-one? Ninety-two.'' A shrug that
is not indifference; it is the shrug of a man discussing a part that is
no longer manufactured. ``So we do the stone.''

``And that's --- all right?''

Ingi considers this seriously, walking, for about ten metres.

``The stone is most of it,'' he says.

Rowan looks back once more. In the last of the light the eleventh stone
is indistinguishable from the other ten: cleared, checked, repacked,
sitting exactly where it should sit, maintained to a standard that would
satisfy any inspector in Europe. There is nothing wrong with it. There
is nothing wrong with any of it. Forty people have spent eleven hours
doing careful, skilled, unpaid, uncomplaining work on a village that
they are keeping in excellent order.

And a woman died in ninety-one or ninety-two --- nobody on this hill can
tell him which, and it does not appear to matter to any of them --- and
took four syllables with her, and the
village noted it, and adjusted, and did not stop, and has been doing the
stone without the form for thirty-four years, and there is not a person
on this hill under the age of fifty who has ever heard the eleventh
stone spoken to.

Sigrún is at the back of the party, twenty metres down the slope, not
climbing well. She has stopped, and she is looking up at him, and she
has seen exactly which stone he is standing next to.

Neither of them says anything. There is nothing to say that either of
them does not already know, and they both have the discipline for it
now, and that shared discipline is the most intimate thing that has
passed between them since the schoolhouse.

She turns and starts down. After a moment, so does he.

\sceneline

There is coffee in the hall afterwards, and there is a great deal of it.

The urns come back out. The kleinur reappear. The dogs, who have had an
outstanding day, lie flat under the tables like dropped coats. Somebody
has put the radio on. Ólöf's cake is finished, aggressively, by
committee. There is an argument at the next table about the gearbox
which appears to have been running since roughly lunchtime and which
both parties are visibly enjoying.

Séra Hjálmar comes in at seven with paint on his hands and paint in his
eyebrows, having spent the whole day on a ladder doing the hall's south
elevation, and is heckled about the colour. He defends the colour. The
colour is, Rowan gathers, the same colour it has been for forty years,
and the heckling is also the same, and the pastor's defence of it is
word-for-word the defence he mounts every year, and everyone including
the pastor knows this, and it is done with great relish on all sides.

``Dr.~Hale.'' Hjálmar drops onto the bench opposite with a mug. ``They
put you on the wall.''

``They put me on the wall.''

``How are the hands?''

Rowan turns them over. Hjálmar laughs --- a real one, loud, from the
chest --- and Rowan finds that he is laughing too, properly, for the
first time in he does not want to calculate how long, at nothing, at
the state of his own hands, in a hot noisy room that smells of wet wool
and coffee and paint.

It is the happiest he has been in eight months.

He is aware of it while it is happening. That is the part he will come
back to, later, on the rock, in other weather. He is aware of it while
it is happening --- that he is sitting in a hall in a village of a
hundred and thirteen people at the end of a road that maps disagree
about, with his hands wrecked and his back seized and a name typed on a
list on the wall in a font somebody chose in 2004, and that he is
happy, and that nobody in the room requires anything of him except that
he pass the milk.

\sceneline

He walks back to the Annex at nine with Kári's grandmother's voice still
somewhere in his ear.

The night is clear and enormous and the cold has real teeth in it now.
Below him the village is lit up in patches, the hall windows yellow, a
few kitchens, the pastor's ladder still against the south wall because
he will do the second coat tomorrow. From up here the whole place could
be held in two hands.

He is tired in a way he has not been tired in years. His shoulders have
stopped being shoulders. He is going to sleep for nine hours and wake
up unable to move.

And somewhere on the climb, without deciding to, he starts going back
over the day the way he goes over an interview, and he arrives, near the
top, at the thing that has been sitting under all of it.

Nobody said any of the words.

Not one. In eleven hours, across forty people and two meals and a hill
in the dark, nobody said \emph{tradition.} Nobody said \emph{custom} or
\emph{heritage} or \emph{the old ways.} Nobody said \emph{ritual} ---
not ironically, not defensively, not once. Nobody explained anything to
anybody. No adult took a child aside and told it why. He has done
fieldwork in eleven communities and in every single one of them somebody
found him within the first hour and started explaining the culture to
him, proudly or apologetically or with an eye on the recorder, and here
he has spent an entire day inside the thing and not one person has
described it, because there is nothing to describe; there is a list on
a wall and a wall on a hill and a stone that gets cleared and checked
and repacked, and four hours of copying, and a birthday cake.

He stops on the path with his breath going white.

The list. \emph{Fjórði laugardagur.} Forty people, every month, for as
long as anyone can name --- and he thinks, with the particular chill of
a man doing arithmetic he was not looking for, about what a village
would have to be maintaining to justify eleven hours a month, every
month, for three hundred years, from every single adult, without
exemption, without complaint, and without anyone ever once being told
what it was for.

Then he thinks about the eleventh stone.

He stands there for a while in the dark, above the yellow windows, with
his ruined hands in his pockets.

Then he goes in, and eats bread and butter standing at the counter
because he cannot face cooking, and sleeps for nine hours, and wakes up
unable to move.

\chapter{}

The maps are the only thing in Brynjavík that anyone has ever tried to
explain to him.

He notices it on the Sunday, climbing the hall steps very slowly, with
three folded sheets under his arm and a back that has strong opinions
about yesterday, and the noticing has an edge on it that it did not have
a week ago. Eleven hours on a hill and nobody described
anything. Two minutes in Thorsteinn's office on his first evening and
the man was already explaining --- subsidence, erosion, other things,
the whole apparatus of causal account, offered unprompted to a stranger
who had not asked.

Rowan had taken it, at the time, for administrative courtesy.

He is beginning to suspect it was loneliness.

\sceneline

``Hale.'' Thorsteinn does not look up. ``Sit. Move the folder.''

The office is warm and smells of the same cleaning solution, and it is,
Rowan sees now, the only room in the village organised on principles he
recognises. Box files labelled on the spine in a consistent hand. A wall
planner. A stapler in a drawer with other staplers. It is the office of
a competent man in his forties who is, by any standard the outside world
would apply, extremely good at his job, and who has not had a promotion
in sixteen years because there is nowhere for him to be promoted to.

Rowan puts the three sheets on the desk and unfolds them, one over
another, and does not say anything.

Thorsteinn looks at them for perhaps four seconds.

Then he says: ``Which figures.''

``Population. And the connector road.''

``One-one-three, ninety-six, forty-one.'' He says the numbers the way a
man says his own telephone number. ``And the road is solid, dashed, and
absent.'' He puts down his pen. He is, Rowan realises, not surprised at
all; he is something much worse, which is pleased. ``You went to the map
shop on Bankastræti.''

``I did.''

``Everyone goes to the map shop on Bankastræti.'' He stands and crosses
to a filing cabinet, and it is not a dramatic movement, but there is a
quality of appetite in it, of a man reaching for something he has been
waiting a long time to hand to somebody. He comes back with a lever-arch
file four inches thick and drops it on the desk between them.

The label on the spine says \emph{KORT --- FRÁVIK.}

Maps. Discrepancies.

\sceneline

There are three hundred and eleven entries and they go back to 2011.

They are on a form. Thorsteinn made the form. It has fields: source
document, publication date, publishing body, feature, expected value,
observed value, first noted, re-checked, notified, response. The
handwriting in the early entries is round and rather generous. By about
2016 it has tightened. In the last two years it is small and upright and
extremely regular, the hand of a man who has decided that his
handwriting is one of the things he can still control.

``Start at the tabs,'' Thorsteinn says. ``The tabs are the ones I have
been able to get an answer on.''

There are nine tabs.

Rowan reads for forty minutes and Thorsteinn lets him, and makes coffee,
and does not hover, and it is the most professional courtesy he has been
shown by anyone in Iceland.

The file is superb. That is the thing that keeps arriving. It is
genuinely, unambiguously superb --- the observations are dated, the
sources are cited to edition and printing, the re-checks are performed
at intervals and logged, the negative results are recorded as carefully
as the positive ones. There are photographs, numbered, cross-referenced.
There are two entries where Thorsteinn has written, in the response
field, \emph{my error --- misread scale bar,} and left the entry in.

Rowan has sat on three funding panels. He would have funded this.

\emph{Entry 44.} District sheet, 1987 impression. Connector road from
the pass: not shown. Expected: shown. Observed: absent. Notified:
National Land Survey, 2012. Response, 2013: \emph{thank you for your
correspondence; the 1987 sheet is superseded and no correction will be
issued to a withdrawn impression.}

\emph{Entry 45.} Query to district roads office: date of construction,
connector road. Response: no construction recorded 1970--2003.
Re-queried 2014, 2017. Same response, twice, from two different
officials.

\emph{Entry 46.} Query to district roads office: if not constructed,
when first maintained? Response: resurfacing budget lines exist
continuously from 1961.

Rowan puts his finger on that one.

``They resurfaced a road that was never built.''

``For fifty years.'' Thorsteinn sits back. ``Every year. There is a
figure. It is a small figure and it is in the accounts and nobody has
ever queried it except me.''

``What did they say when you queried it?''

``That it is a legacy line and the coding predates the current system.''
He says this without bitterness, which is somehow worse. ``Which is
true. It does predate the current system. Everything here predates the
current system.''

\sceneline

Tab six is the stake.

\emph{Entry 208.} Kilometre marker 47, connector road. Surveyed
position established by GPS 14/05/2019, cross-checked against 1:50,000
sheet and against the 1961 alignment plan. Stake set 14/05/2019.

\emph{Entry 209.} 22/05/2019. Stake found 780 m south of set position.
No disturbance to ground at set position. No drag marks. Ground at new
position undisturbed except by stake.

\emph{Entry 211.} Re-set 23/05/2019, concrete collar.

\emph{Entry 214.} 09/06/2019. Stake found 810 m south of set position.
Concrete collar found at set position, intact, empty.

Rowan reads that twice. \emph{Intact, empty.}

``I photographed it,'' Thorsteinn says. ``Photograph two-two-one. The
collar is intact. You cannot pull a stake out of a set collar without
breaking the collar or the stake. The stake is also intact. It is
eight hundred metres away in undisturbed ground.''

``What did you conclude?''

``That somebody in this village is very good at concrete.'' He shrugs.
``I did consider it seriously. There are perhaps four men here who could
do it. It would take a night and a mixer and a reason, and I have never
been able to find the reason, because it does not benefit anybody. The
stake is not a boundary. It is not planning. It is a distance marker on
a road with nine vehicles on it.''

``And the third time?''

``The third time I reported it.'' Something in his face closes. ``An
inspector came from Akureyri in September. A good man. Thorough. He
spent two days. He found the stake at kilometre forty-seven, correctly
positioned, in a sound collar, and he measured it three times, and he
wrote a report which said that the marker was correctly positioned and
that the reporting officer might be under strain.''

Rowan says nothing.

``It was not correctly positioned,'' Thorsteinn says. ``I measured it
myself the morning he left. It was eight hundred metres south.'' He
turns his coffee cup a quarter turn on the desk. ``That is in the file
too. Entry two-two-six. I wrote it down and I did not notify anybody,
and that is the last entry I notified anybody about. There are eighty-
five entries since.''

\sceneline

They walk down to the harbour at three because Thorsteinn wants to show
him the tide gauge, and on the way, without any particular ceremony, he
gives Rowan his history.

Hafnarfjörður. A two-year contract in 2009, taken because there were no
jobs and because it was a two-year contract. He met Vigdís in the second
month, which is a thing that happens, and married her in the second
year, which is a thing that happens less.

``And she is not here,'' Rowan says.

``She is in Egilsstaðir. She has a boy who is eleven and is not mine and
that is entirely all right.'' He says it evenly, and Rowan believes the
evenness about eighty per cent. ``She left in 2016.''

``Can I ask why?''

Thorsteinn walks a few paces.

``Because I would not learn it,'' he says.

It is the first crack. It is a hairline.

``The speech?''

``She tried for three years. Her mother tried. They are patient people
about this, you will have noticed --- they do not push, they do not
lecture, they simply keep leaving the door open, month after month, and
it is unbearable.'' He puts his hands in his pockets. ``I have Icelandic.
I have good Icelandic. I run this district in Icelandic. And they wanted
me to sit in a kitchen at forty years old and be taught to say a word
wrong on purpose.''

``That's not quite what it is.''

``I know what it is.'' Sharp, and then immediately regretted; he lifts a
hand. ``Forgive me. I know what it is. It is a dialect with an unusual
case system and a body of oral formulae, and it is interesting, and you
are the right man to write about it and I hope you do.'' He stops at the
top of the harbour steps. ``What I would not do is pretend that saying
it correctly keeps the water out. Because that is what she wanted. Not
the words. She did not care about the words. She wanted me to
\emph{believe} it, and I would not, and after seven years that turns out
to be the whole marriage.''

The bay below them is grey and flat and perfectly ordinary.

``I do not think she was stupid,'' Thorsteinn says. ``I want to be clear
about that. She is the least stupid person I have ever met. She has a
master's degree.''

\sceneline

The tide gauge is bolted to the harbour wall and it is a good one, and
its record does not agree with the tide tables, and Thorsteinn has
sixteen years of that too.

He shows Rowan the graphs on his phone, standing in the wind. The
divergence is not random. It is not noise. It correlates with nothing
meteorological --- he has checked barometric pressure, wind fetch, the
lot --- and it is small, four or five centimetres, except on about nine
occasions a year when it is not small at all.

``What's on those nine days?''

``Nothing.'' Thorsteinn puts the phone away. ``I have been through the
parish diary, the school calendar, the weather, the fishing log. There
is no pattern. Some of them are funerals. Some of them are Tuesdays.''

Rowan thinks about a fourth Saturday and does not say it, and the not
saying it is so easy, so smooth, so exactly the shape of every silence
he has been trained to admire, that he stands in the cold with his ears
ringing slightly.

He could tell him. It would take one sentence. \emph{Cross-reference
your nine days against the fourth Saturday of the month and against the
turn of the year, and tell me what you get.} Thorsteinn would do it that
night. He would do it properly, with the rigour of the best empiricist
in the district, and by morning he would either have nothing --- or he
would have the thing he has been building toward for sixteen years,
handed to him by an outsider, in a form he could not dismiss.

Rowan says: ``It's a good dataset.''

``It is a very good dataset,'' Thorsteinn agrees, ``of an effect that
does not exist,'' and laughs, briefly, and it is the worst sound Rowan
has heard in Brynjavík.

\sceneline

He asks about the Keeping on the way back up, because he cannot not.

``Oh, I go.'' Thorsteinn seems mildly surprised at the question. ``I was
on the roof at the curing house. Every month.''

``You do the work.''

``Of course I do the work. It is eleven hours of the entire adult
population doing free maintenance on the public building stock.'' He
gestures at the hall as they pass it, at the south elevation with its
second coat drying. ``Do you know what that is worth to a district
budget of this size? I costed it. It is eleven million krónur a year of
labour I do not have to find.'' He almost smiles. ``I would fight
anybody who tried to stop it. It is the single most functional
institution in this village and it is the only reason the fabric is
still standing.''

``And the boundary stones?''

``The boundary stones are a folk custom and they take forty minutes and
they keep the old people happy.'' He holds the door. ``I have no
objection to folk custom. I have an objection to folk custom being cited
at me in a planning meeting.''

Rowan goes through the door.

Sixteen years, he thinks. Three hundred and eleven entries. A tide gauge
and a form he designed himself and a photograph of an intact empty
collar. This man has assembled, alone, unfunded, and against active
institutional discouragement, the most complete evidentiary record of
the phenomenon that exists anywhere in the world --- and he goes up on
the roof every fourth Saturday and comes down again and files the
paperwork, and in sixteen years the one connection he has never made is
the one between the two halves of his own life.

He is doing the Keeping. He has been doing it the whole time, in a
register that has no power in it: maintaining a boundary by documentation,
month after month, meticulously, without exemption, because he cannot
bear the alternative.

Rowan does not say this either.

\sceneline

The map comes out at half past four.

It is the erosion map, and it is three metres wide, and it goes on the
long table with weights on the corners, and it is beautiful work.
Thorsteinn has drawn it himself over two winters: the coastline at
1912, 1961, 1987, 2009, and now, in four colours, with the survey
control marked and the sources keyed in a legend in the bottom right.

``The inspector is Tuesday,'' he says. ``It needs a quorum by then.
Nine signatures.''

``How many have you got?''

``Four.'' A pause. ``Two of those are the Sigurjónsdóttir sisters, which
took eleven days and cannot be repeated.''

He gets a pen out of the drawer. He does not push it across. He sets it
down beside the map, which is worse.

``You have a doctorate,'' he says. ``You are on a national research
grant. You have been here five weeks and you have walked every metre of
that shoreline with a survey sheet in your hand, which I know because I
have watched you do it.'' He straightens a weight that does not need
straightening. ``An independent academic signature is not a quorum
signature. It is better. If a Cambridge linguist signs as an observing
witness to the survey positions, the funding office in Akureyri stops
treating this as a village complaint.''

Rowan looks at the map.

``It would be improper,'' he says.

``It would be---''

``No, listen. I'm not being delicate.'' He puts his hand flat on the
table beside the coastline. ``I'm not a surveyor. I have no
qualification in geodesy or coastal process. If I sign that as a witness
to survey positions and anybody in Akureyri looks up what I actually
am, the document becomes weaker than it is now, not stronger. They will
see a linguist certifying a landform survey and they will discount the
whole file. Including your file. Including sixteen years.''

Thorsteinn is quiet.

``That is a good argument,'' he says at last.

``It's a true one.''

``I did not say it was not true.'' He looks at Rowan with a directness
that is entirely new, and Rowan understands that he has miscalculated,
that this is the best-trained mind in the village and it has spent
sixteen years learning to hear exactly this, the fluent correct
unanswerable reason. ``I said it was good. There is a difference and you
know it better than anybody who has ever come here.''

The room is very quiet. Outside, gulls.

``Is there a second reason?'' Thorsteinn asks.

And there it is, on the table, between the 1912 coastline and the 2009
one, and all Rowan has to do is pick it up.

The second reason is that he does not believe the map. Not the
measurements --- the measurements are impeccable. He does not believe
the \emph{premise}: that there is a line of any date, drawn in any
colour, keyed to any control point, which describes where this coast
ends, because he has stood on a hill and watched a village hold a
boundary in place with eleven hours a month and a set of syllables, and
he has watched them do the eleventh stone without the form and he knows
what that will cost, and to put his name on a piece of paper asserting
that the edge of Brynjavík is a measurable and stable feature of the
physical world would be, at this point, a lie he can no longer tell.

Five weeks ago he would have signed it in a second and thought nothing.

``No,'' Rowan says. ``There's no second reason.''

Thorsteinn nods slowly. He picks the pen up off the table and puts it
back in the drawer, and closes the drawer.

``All right,'' he says.

\sceneline

Rowan leaves at ten past five and stops in the corridor to put his coat
on, and through the half-open door he can see Thorsteinn at the wall
behind the desk, where the overlays are.

He has a fresh sheet of tracing paper. He is pinning it over the others
--- there are eleven or twelve under it now, the whole coastline in
strata, each one a little further in --- and he is squaring it up, and
his hands are steady, and he takes a pencil and begins the new line
along the top from the northern headland, working down, the way he has
done every autumn for sixteen years.

Adding a layer. Nine signatures needed and four in hand and an inspector
on Tuesday, and a man alone in an office at five in the evening drawing
this year's edge over last year's edge over the edge before that, in a
hand that will be smaller than last year's, and the whole stack getting
thicker, and every single line on it correct, and not one of them
holding anything.

Rowan buttons his coat and goes out into the dark.

He tells himself, on the path, that he was right. That he had no
qualification and no standing and that a false signature would have done
the man active harm. Every word of it is true and he assembles it
carefully as he walks, the way he has assembled such things his whole
life, and it holds all the way to the Annex door.

Then he thinks: \emph{cross-reference the nine days.}

One sentence. He could have given him one sentence.

He stands with his hand on the latch for a while in the cold, and then
he goes in and does not light the lamp.

\chapter{}

He spends the morning collecting hymnals, and tells himself it is
fieldwork.

It is not difficult. The village lends books the way it lends salt ---
without ceremony and without asking why. He knocks on eleven doors
between nine and one, and in every kitchen he is given coffee he does
not want and a hymnal he does, and in most of them a child is sent to
fetch it from wherever hymnals live, which is variously a dresser, a
windowsill, and in one house a biscuit tin on top of a refrigerator. At
Halla Kristjánsdóttir's he is fed dried fish and asked, at length and
with real interest, whether it is true that in England the water comes
out of the tap already warm. Kári has cracked the card trick since
Saturday. It now works, twice, and the second time is better than the
first, and he accepts Rowan's applause as no more than the settlement of
an outstanding debt. On the wall behind him hangs a school photograph:
fourteen children on the steps of a building in Húsavík, ninety minutes
each way on the bus that comes at seven --- the bus about which
Steinunn, as the whole village is aware, holds views.

By afternoon he has eleven hymnals on the Annex table, spanning a
hundred and six years of printings: Akureyri 1928, Reykjavík 1945,
Reykjavík 1972, a 1991 reissue, two undated, one with a Danish
imprimatur so old the paper has gone the colour of weak tea. Three
different presses. Four different centuries of orthographic convention.
Two of them were bought in Reykjavík within the last decade, by people
who drove to the city and walked into a shop and paid for them off a
shelf.

He opens them in sequence to page nine, or its equivalent, and lays them
in a row.

\emph{Verði að þínum vilja.}

\emph{Verði að þínum vilja.}

\emph{Verði að þínum vilja.}

Eleven for eleven. Including the two bought in the city, which cannot
possibly have left the press that way, and which he checks twice, and
then a third time with the loupe, looking for the erasure, the scraped
fibre, the pen-stroke, and finds nothing --- the paper is undisturbed,
the ink is press ink, the kerning is machine-set and correct for the
longer line. These books were not altered after printing. They were
printed like this. Books that were bought in a shop in Reykjavík, off a
shelf, from a stack of identical books, and carried three hundred
kilometres north, and opened here.

\begin{fieldnote}
11/11. Presses: Akureyri, Reykjavík (×3 houses), Copenhagen (?), 4
unattributed. Span 1928--2019.

\struck{Local overprint. Someone in the parish orders a custom run.}
--- No. Two acquired retail, Reykjavík, 2014 and 2019, by Halla K.\ and
by Þ.\ Sveinsson, independently, neither of whom has any reason to lie
about where a hymn book came from and neither of whom found the question
interesting.

\struck{I am misremembering the standard text.}
--- No. Verified against my own recollection, against Sveinsson's
account of hearing it said otherwise on the radio, and against the fact
that Sveinsson, asked to recite it, produced the dative and then
laughed and said \emph{that is not how they say it on the radio, is it,}
with the mild amusement of a man noting a regional word for a bread
roll.

Remaining hypothesis, which I do not like: the alteration is not in the
books.
\end{fieldnote}

He sits for a long time in the failing light with eleven open books in
front of him, and does not write the next line, because the next line is
that a text does not have to be edited if the coast will simply not
permit the other reading to arrive intact --- that the books came north
in the nominative and were, somewhere along the road, declined.

He thinks of a fluorescent stake set eight hundred metres early. He
thinks of a GPS pinning a village half a kilometre out to sea. He closes
the books, one after another, and stacks them, and squares the stack,
and finds that his hands want to align it parallel to the desk edge, and
lets them.

\sceneline

The fluorescent tubes overhead hum with a frequency that scratches at
the back of the skull. It is 19:00. The community hall is at capacity,
though the capacity has shrunk with the years: a hundred and ten chairs
in concentric semicircles facing a heavy oak lectern, set out from a
stack that used to run to three hundred and now stops, unremarked, at a
number nobody has had to reconsider in a decade. Outside the wind is a
theoretical pressure against the north wall, muffled by the insulation
Thorsteinn commissioned three years ago against a cold that was never
the real problem.

Rowan sits in the back row, notebook open on his knee, recording the
ambient detail: the smell of damp wool drying in the heat, the scuff of
boots on linoleum, the way the villagers hold their bodies, tight and
contained, knees pressed together as if to minimize the air they
displace. Thorsteinn stands at the lectern looking eroded, his jacket
over the back of his chair, his shirt wrinkled at the elbows. He holds a
sheaf of government forms --- Relocation Grant Application, Form 4B ---
and a red pen, tapping it against the paper, dry and percussive.

``We need to finalize the zone definitions,'' he says, his voice pitched
slightly too high for authority. ``The inspector arrives Tuesday. If the
erosion map isn't signed by a quorum, the funding for the retaining wall
is frozen.''

A thick-set man in the second row clears his throat. ``The line isn't
where the map says it is, Thorsteinn. The survey's from July. The
shingle has moved.''

``The survey is the legal document.'' Thorsteinn runs a hand through his
hair. ``It doesn't matter where the shingle is today. It matters where
the surveyors put the stake. We just need to sign that we agree the
stake exists.''

``But the stake is gone,'' an older woman says from the corner, knitting
something grey and endless, not looking up. ``The tide took it on
Thursday.''

He closes his eyes. ``Then we assume it is still there. For the
paperwork.''

Rowan watches the tension gather in the administrator's shoulders. Here
is a man trying to impose a grid on a fluid, to conjugate the village
into a stable tense, and Rowan feels for him with an intimacy that
surprises him, because he has spent his own life doing the same thing to
a grief that would not sit still and let itself be surveyed.

``We cannot sign a lie,'' the woman says. Her needles click. ``To sign a
lie is to invite a correction.''

``It is not a lie. It is bureaucracy. It is how we get the money to pour
the concrete. Do you want the concrete, or do you want to keep watching
the path fall into the---'' He stops, catching the word before it slips.
``---into the bay.'' The room shifts, a rustle of coats. \emph{Bay} is
permitted, barely. A euphemism, a polite fiction.

An elderly fisherman stands near the front, leaning on a driftwood cane
polished smooth by salt and palm-oil. Rowan recognizes him as the man
who would not speak to him on the quay that first day. ``The concrete
will not hold,'' the fisherman says, gravel rolling in a drum. ``You are
trying to bribe the winter with receipts. You are filing forms against
hunger.''

``I am filing forms against the incompetence of the district planning
office!'' Thorsteinn shouts, and the outburst cracks the room's
controlled fatigue. He slams his hand on the lectern; the red pen jumps.
``I am trying to keep the road open. To keep the power on. While you sit
here and talk about hunger and correction and refuse to tell me where
the water line is because you're afraid of a grammar rule from the
nineteenth century!''

``It is not a rule,'' the fisherman says, calm. ``It is a negotiation.
And you are shouting at the other side of the table.''

``There is no other side!'' Thorsteinn steps out from behind the
lectern, small against the vast dark map of the coastline pinned behind
him. ``There is us, and there is the physics of the North Atlantic, and
that is all there is. Just because old women whisper to the fog does not
mean the fog is listening.''

Rowan stops writing. The ink bleeds into the paper in a dark bloom. The
air thins, as if the oxygen were being drawn quietly off.

``Sit down, Thorsteinn,'' the knitting woman says. ``You are tired. The
tongue slips when the body is tired.''

``My tongue is fine.'' He is shaking now --- not with fear but with the
vibration of a structure pushed past its rating. ``I am the only one
here speaking clear Icelandic. The only one trying to save this place
from its own silence.''

``The silence saves us,'' the fisherman says, pointing his cane at the
map. ``You mark the boundary. You do not name it.''

Thorsteinn rips the map from the wall. Thumbtacks ping against the
floor. He shakes the paper at the room, the inked coastline dancing.
``It's water!'' he yells.

The fisherman freezes. The needles stop. Thorsteinn steps forward,
panting, and looks at the fisherman, then at the window, where the black
reflection of the room hangs against the night. ``It is just water,'' he
says, and he uses the nominative, and Rowan hears the case land like a
struck bell. \emph{Þetta er bara vatn.} He names it. He gives it a
subject. He reduces it to a neutral substance, a thing that can be held
and measured and contained. ``It rises because of the moon. It falls
because of gravity. It does not think. It does not remember. It is just
water.'' He drops the map, and it flutters face-down to the floor.

Silence. Not the silence of a pause but the silence of a deletion.
Rowan's ears pop. The fluorescent hum continues, distant now, as if
transmitted from another room, and the pressure in the hall drops the
way it drops in a house when a door somewhere has been opened onto a
much larger space. No one moves. The fisherman stares at the air in
front of Thorsteinn as if expecting it to split. The knitting woman
studies her hands and waits.

Thorsteinn looks around, expecting rebuttal, argument, the ordinary
friction of a meeting. ``Well?'' he demands. ``Correct me.''

No one speaks. To correct him would be to repeat the error. The
fisherman sits and lays his cane across his knees, aligning it
carefully. The knitting woman folds her wool and stands. The villagers
look past Thorsteinn now, at the walls, the ceiling, the shelter itself,
the way people in a listing ship look at the hull. Thorsteinn wipes his
mouth, and his eyes find Rowan, and confusion flickers across his face
as he understands, all at once, that the silence is not agreement.

``The meeting is adjourned,'' he whispers. He reaches for the red pen,
misses the cap, tries again, misses, and leaves it uncapped. Rowan
closes his notebook. He does not record the breach; to write it would be
to archive it, and he has learned, this week, the difference. He stands,
his knees unreliable, and the villagers file out without speaking,
moving like water toward a drain, and Thorsteinn remains, staring at the
bright rectangle on the wall where the map used to hang, a window onto a
time before the boundary was drawn.

\sceneline

Rowan opens the door onto a night that has stopped entirely. No surf. No
flagpole rope. The darkness toward the harbour is thick and waiting. The
sentence has been spoken; now the listener will parse it.

At a quarter past eight he stands on the concrete promenade that skirts
the inner harbour, and he switches on his flashlight, and the harbour is
moving. Not violently --- there is no cracking pavement, no scream of
rebar. The quay is simply descending. As he watches, a galvanized bench
sinks; the concrete beneath its bolts loses its solidity and the surface
closes over the seat, and then the backrest, and then it is gone. The
bollards retract into the ground like pistons. And where the modern
surface sinks, something older rises to meet it: dark timber, tarred,
rough-hewn, joined with wooden pegs, so that wood and concrete occupy
the same coordinate in two different drafts of time. This is not
erosion. It is revision. The quay is reverting.

Salt bleeds from the retaining wall and crystallizes in the cold air,
and white geometries bloom between the stones --- not random
efflorescence but ordered, angular, the shapes he has learned to read.
Case endings. The wall is diagramming, in salt, the sentence that used
to bind it, and the sentence is being un-said, and so the wall is going
back to what it was before the sentence held it. He reaches for his
recorder. The red light blinks once and dies into silence.

The modern breakwater dissolves. Rocks surface where the charts have
shown open water since 1890. He understands, with the slow cold clarity
that has replaced his curiosity, that the old map was not a plan. It was
a memory, and the sea is choosing to remember it.

He tries to write, and his hand forms \emph{The---} and locks. The
definite article. He crosses it out and tries \emph{It,} and \emph{it
moves,} and the nausea recedes, and he understands that even his own
language has begun to take sides, that the reaching, agentive
constructions his discipline is built on will no longer pass his hand
without a toll. He puts the notebook away.

Arnar Jónsson steps into the edge of the flashlight's beam. ``It goes
back,'' he says.

``The concrete liquefied,'' Rowan whispers.

``It was not remembered.'' Arnar does not say the word Thorsteinn said.
``To name it is to deny that it has a will. And what has no will has no
boundary. It does not want to be kept.'' He watches the black timber
surface out of the mud. ``That is my grandfather's rack. Lost in
forty-two.'' A pause. ``You should go to Thorsteinn. He will need
witnesses.''

``For what?''

``For the end of the sentence.''

``When does it stop?''

``It does not reach the present,'' Arnar says, and there is no drama in
it, only the fatigue of a man reporting a tide. ``The present falls back
to meet it.'' He walks toward the oldest rocks, and the ground hums
underfoot, a low seismic note, a system quietly overwriting itself, and
Rowan stands at the edge of it and feels, for the first time, that he is
not in the file at all --- that he is watching a page be rewritten and
is himself only a smudge of ink in the margin that the new draft has not
yet decided what to do with.

\sceneline

At eight in the morning the village remains, but the angles are wrong.
Rowan enters the community hall to find the villagers standing outside
the office, listening. From behind the frosted glass of the door comes a
voice, broken into pieces. \emph{``Nom\ldots{} nom\ldots{} in\ldots{}
a---''}

Thorsteinn sits at his desk, immaculate, the papers squared, every sheet
inked solid black. \emph{``Sub\ldots{} ject---''} He looks up. There is
terror in his face, but it is a terror without connection to anything, a
man frightened by the inside of his own skull. \emph{``The---''} he
tries. He cannot reach the noun. The pages under his hands are filled
edge to edge with suffixes. Nominative endings, endless, a subject
rehearsed a thousand times without ever arriving at the thing it was
meant to name.

Rowan reaches for the word \emph{diagnosis,} and it slides out from
under him, and the loss is sudden and physical, a stair that is not
there in the dark. The contagion is close. Thorsteinn writes \emph{W}
again and again until the paper tears. \emph{``Is---''} He is trying to
rebuild the sentence from its smallest parts, the way a man with a
shattered hand tries to make a fist. The villagers wait, not unkindly,
the way people wait outside a room where something is being finished.
Rowan backs out. He needs to hear a complete sentence, spoken by
someone, before he forgets how one is made. Behind him Thorsteinn keeps
carving the nominative into the desk with the uncapped pen, trying to
name the thing that has already, patiently, finished naming him.

\chapter{}

In the morning he goes to look at the house where his mother was born,
and there is grass.

Edda told him where it stood, weeks ago, in a sentence so flat he did
not understand at the time that he had been given anything: \emph{past
the western markers, where the ground goes soft.} He walks out at eight
with the survey sheet and a compass, over the boundary line, past the
stone Arnar planted, into the part of the map that the village has
stopped speaking.

The ground is soft. The air is colder here by a degree he could measure
if he had brought a thermometer, and he did not, and he notes the
omission the way he notes everything.

He finds the site by triangulation, from the church tower and the
northern headland, and stands on the coordinate.

Grass. Not a footprint of stone, not a hollow, not the rectangular
darkening in the sward that a foundation leaves for two centuries in
any field in England. Not rubble. Not the nettles and elder that grow
on old human ground everywhere in the world because human ground is
rich. Just the moor, green-grey and uninterrupted, running unbroken over
the place where a kitchen was, in which a woman was born in March 1948,
and in which she presumably learned to walk, and at some table in which
she was taught to decline a coastline.

He crouches and puts his hand flat on it.

He is a rational man in a cold field with his palm on some grass. He
knows it. But he has excavated at four sites and he knows what land does
with a house, which is that it remembers it badly and forever, and this
land has not remembered it at all.

He stays about ten minutes. He does not take a photograph, and cannot
afterwards say why, and suspects it is because a photograph of nothing
is a document that proves the wrong thing.

\sceneline

Séra Hjálmar has the registers out on the vestry table before Rowan
finishes asking.

``I wondered when.'' He is unembarrassed about it. ``You have been here
six weeks and you have not once asked me for a name. Every researcher
asks for a name in the first fortnight. It is always a grandmother.''

``I wasn't sure I wanted it.''

``No.'' Hjálmar opens the volume. ``Which is why I did not offer.''

The register is a beautiful object, half-leather, the ruling still crisp,
and the hand of the 1940s incumbent is a small upright copperplate that
must have taken him twenty minutes a line.

\emph{3.~mars 1948. Bryndís Ásta Guðmundsdóttir.}

Rowan sits down.

It is not a shock. He has known the shape of this for eight months and
the fact of it for three weeks. What it is, is the first time his
mother's name has been written down in front of him by somebody who was
not her, and there is a difference between knowing a thing and seeing it
ruled into a column between two other children, and the difference goes
through him at the sternum like cold water.

\emph{Bryndís Ásta.} Two names. She used neither of them for fifty-seven
years.

``Confirmed,'' Hjálmar says, turning forward, ``1962. Fourteen. Here
--- with eleven others, which was a large year.'' His finger moves down.
``And she is in the communicant roll for sixty-three, sixty-four,
sixty-five.'' He turns another leaf. ``And then here.''

The 1966 page. Halfway down, in the margin beside her name, in the same
small upright hand, four words and a date.

Rowan reads it and does not recognise the construction.

It is not \emph{brottflutt} --- moved away, the standard notation, which
appears twice on the same page against two other names, both with
destinations. It is not a death entry; deaths are ruled in red and there
is no red here. It is not a transfer of parish. It is a phrase he can
parse word by word and not as a whole, and after fifteen seconds of
trying he does what he has spent his life training himself not to have
to do, which is to ask.

``What is that?''

Hjálmar looks at it for a while.

``It says the line is not finished,'' he says.

``Whose line?''

``Hers.''

``Meaning she had no children? She had a child. She had me.''

``Then it does not mean that.'' Hjálmar closes his hands on the edge of
the table. ``Dr.~Hale, I will tell you exactly what I know, which is
less than you want. That notation appears in this book nine times in two
hundred and ninety-three years. I have checked. I checked in my second
year here, because I found one and could not read it either.'' He turns
the volume slightly so the light falls better. ``It is not
ecclesiastical. It is not in any handbook. Three of the nine are in
hands from before 1800. It is not something a pastor invented; it is
something pastors were told to write.''

``By whom?''

``That is the part I do not know.'' He says it simply. ``I have written
the phrase in this register twice myself, in nineteen years, both times
because I was asked to, both times by someone who did not explain, and
both times I did it, and I would do it again, and I could not give you a
defence of that which would satisfy a synod.'' He almost smiles.
``Fortunately the synod has never asked.''

Rowan looks at the four words against his mother's name.

``What happens to the other eight?'' he says. ``The other eight people.
Where do they go?''

Hjálmar turns the pages and finds them, one after another, and reads
out the subsequent entries in a level voice, and there is nothing
sinister in a single one of them. Two died here, old, decades later,
with ordinary red-ruled entries. One married in 1811 and is in the
register four more times for four children. One emigrated to Canada in
1888. One is a man who lived his whole life in this parish and died in
1954 and whose funeral, per a marginal note, was very well attended.

They are ordinary lives. All of them.

``It is not a curse,'' Hjálmar says. ``I am going to say that plainly,
because you have gone white. It is not a sentence passed on anybody. Whatever
it records, the people it was written against went on and had lives and
several of them had long ones.'' He closes the book gently. ``It records
something that was left open. That is all I can tell you and it is all
I have ever been able to tell myself.''

\sceneline

The school ledgers are still on the desk in the schoolhouse office where
Rowan left them, arranged in their chronological arc, gathering a
fortnight of new dust.

He came here the first time for the curriculum. He read 1965 for the
structure of the syllabus and the teacher's marginalia and the shape of
the attrition, and he transcribed four pages of it, and he noticed the
personal notes only as texture --- \emph{Jónas struggles with the vocal
length. Gudrún is excellent.} --- and he did not read them, because they
were not what he had come for, and because he did not know there was
anything in this building with his own blood in it.

He opens 1965 and goes through it line by line with his finger.

Week 9. Week 14. Week 22.

\emph{Week 31: Dative declensions, coastal nouns. Bryndís Á.\ has the
length. Better than mine at that age --- she does not push it from the
throat, it simply arrives. Þórdís says the same and Þórdís does not say
such things.}

Rowan stops. He reads it four times.

\emph{Better than mine.} The head teacher of this school, a woman with
thirty years in the room, writing in a working document about a
seventeen-year-old.

\emph{Week 36: Bryndís Á.\ took the second clause on Saturday, at the
north stones, in front of everyone. Fifteen and a half seconds. Nobody
had to carry her. Guðmundur cried, which he will deny.}

\emph{Week 38: The broadcast people come Thursday from Reykjavík.
Þórdís has said Bryndís is to do the paradigm. She is the only one who
can do the whole of it without stopping and it is right that it should
be her. I am proud enough to be ashamed of myself.}

\emph{Week 39:}

The entry for week 39 is four words long.

\emph{Þórdís has withdrawn her.}

Nothing else. No explanation, in a ledger which for thirty-eight weeks
has explained everything, including the weather and Jónas's throat and
the price of chalk. The hand does not shake. It simply stops, and the
next entry, week 40, is about a broken window and is entirely normal,
and Bryndís Á.\ does not appear in the ledger again, not in week 41, not
in week 47, not in the summary at the year's end.

She was seventeen and she was the best in the village and eight weeks
later she was gone from the country.

\sceneline

He finds the reel because Sigrún tells him where to look.

She is in the schoolhouse, in the blanket, at the child's desk by the
window, and she is worse than she was a fortnight ago in a way that
neither of them will mention. He tells her what he has found and she
listens with her eyes closed.

``The broadcast,'' she says. ``Yes. They came about every four years,
the collectors. Dialect survey. We were on the list because we were
remote.'' A dry breath that might be amusement. ``We were extremely
cooperative. We gave them sheep words for twenty years.''

``There's a copy here?''

``There is a copy of everything here. The school kept the second reel;
that was the arrangement, they took one and left one.'' She lifts a hand
an inch off the blanket and points it at the corridor. ``Store cupboard.
Bottom shelf, under the county maps. And there is a machine, and it will
not work, because there has been no power in this building since two
thousand and four.''

``I have a converter and a battery pack.''

``Of course you do,'' says Sigrún.

It takes him ninety minutes. The reels are in flat card boxes labelled
by year in a hand he now recognises from the ledgers, and 1965 is where
she said it would be, and the machine is a Danish thing the size of a
suitcase which has been maintained --- he notes this, he cannot help
noting it --- with a small can of oil that is on the shelf beside it and
which has been used, at some point, by somebody, in a building nobody
has taught in for thirty-one years.

He carries it all to the classroom and sets it on the teacher's desk and
runs it off the battery pack, and Sigrún makes him lace the tape twice
because he does it wrong the first time.

\sceneline

There is four minutes of a man from Reykjavík establishing the date, the
place, the equipment, and his own name, in the flat confident radio
Icelandic of 1965.

Then a girl says her name.

\emph{``Bryndís Ásta Guðmundsdóttir. Sautján ára.''}

Rowan puts both hands on the desk.

It is not that it sounds like her. It does not sound like her. The voice
on the tape is high and quick and completely unguarded, a girl's voice
with a coast in it, and his mother's voice was low and English and
sanded flat and he never once in fifty-five years heard her produce a
vowel like the one in \emph{sautján}. There is no recognition in it at
all.

Except that the pause is his mother's pause. Before she says \emph{ára}
--- years --- there is a fifth of a second that he has heard ten
thousand times across a kitchen table, and it is not hesitation, it is
placement, a woman putting a word down exactly where she means it to go,
and he sits down slowly on a child's chair with his hand over his mouth.

The recordist asks her to give the paradigm.

She gives it.

She is astonishing. Rowan has listened to perhaps nine hundred hours of
elicited speech in his career and he knows within eight seconds that he
is listening to something he will not hear again --- not the fluency,
fluency is common, but the \emph{placement,} every ending landing in the
same acoustic position, the vowel lengths so regular he could set a
metronome against them, and no strain anywhere. She is not performing.
She sounds bored, in the specific way of a very good seventeen-year-old
doing something easy in front of adults.

Nominative. Accusative. Dative. Genitive.

And then she keeps going.

The fifth column. The sixth. He does not have names for what she is
doing but he can hear the architecture of it, the same stem carried into
cases that his discipline does not admit exist, and the recordist says
nothing because the recordist has no idea, and the tape hiss goes on,
and Rowan sits in a cold classroom in the dark with his eyes shut and
listens to his mother do at seventeen the thing he has spent thirty
years building instruments to detect.

Then the seventh.

She goes wrong.

It is not a stumble. That is the thing he will not be able to put down
afterwards. There is no catch in the breath, no false start, no
correction, no rise in pitch. She produces a form, cleanly and in
tempo and with the same bored competence as everything before it, and
the form is wrong --- and he knows it is wrong the way he knows a bell
is cracked, from the shape of the sound and not from any rule he could
cite, because he has no rule, because there is no published description
of the seventh case of anything.

She finishes the column. She says \emph{takk.} The recordist thanks her
warmly and says it was very good.

The tape runs on for another forty seconds of room tone.

And in the room tone, faintly, at the edge of the microphone's reach,
somebody --- a woman, older --- says one short word, and it is not on
the transcript sheet in the box, and Rowan rewinds it eleven times and
cannot get it, and on the twelfth he stops trying, because Sigrún has
opened her eyes.

``Did you hear that?'' he says.

``I heard it.''

``What is it?''

Sigrún looks out of the window at the flat grey. ``I was three years
old,'' she says. ``I am not the person to ask.'' A long pause. ``But
that is Þórdís, and she is not surprised.''

\sceneline

He sits with it until the battery pack dies.

He can construct two accounts and they fit the evidence equally.

In the first, a girl of seventeen at the end of a long recording, under
the eye of a man from the capital and a room of adults who are proud of
her, on the seventh and hardest column, at the outer limit of a system
that perhaps four people alive can hold, makes an error. It happens. It
happens to everyone. Þórdís withdrew her from the broadcast two weeks
later because a custodian does not put a girl on a public reel making
errors in the deep cases, and the entry in the ledger is four words long
because the head teacher was disappointed and disappointment is not
minuted. And a girl who had been told her whole life that she was the
best, and who found out on tape in front of the village that she was
not, left the country eight weeks afterwards and did not come back, and
that is not a mystery, that is the oldest story there is.

In the second, she did it on purpose.

Rowan turns the two of them over for a long time in the dark classroom.
The evidence does not choose. The tape does not choose; a deliberate
error, produced by someone that good, would sound exactly like an error
produced by someone that good, and the whole point of being that good is
that nobody in the room can tell. Þórdís's word in the room tone is
consistent with dismay and consistent with recognition and he cannot
hear which it is and never will.

What he has, at the end of it, is not an answer. It is a demonstration.

His mother, at seventeen, in this village, could hold the whole of a
grammatical system that no linguist has ever described, and had been
handed it by the woman who owned it, and was, by the assessment of every
adult who heard her, the best in the line. And then something happened
in the seventh case, and she stopped, and a pastor wrote four words in a
margin that mean \emph{something was left open}, and she went to England
and became a woman who worked for a county council drawing boundaries
on maps of a country she had no history in, and married and had a son
and answered his questions with a warm grey surface for thirty years,
and died, and left him a tin with forty torn drills in it and a sealed
envelope addressed to the only other person alive who was in that room.

And on a table in a cottage on the north ridge, six weeks after he slid
it across the wood between the salt and the tea, there is a sealed
envelope with four characters and a full stop on the face of it, in that
same hand.

He wants it. That is the plain fact of the evening and he sits with it
in the cold and does not dress it up. He wants to walk up the ridge now,
tonight, in the dark, and ask for it back, and take it somewhere with a
lamp, and open it.

Edda has not offered. She has not mentioned it once in six weeks. She
did not move it toward herself on the first day and she did not move it
away, and Rowan understands, sitting in a dead classroom with a dead
tape machine in front of him, that this was never a woman deciding
whether to accept a letter. It was a woman deciding whether the person
who brought it was going to be able to hold what was in it, and it is
six weeks later and she has not decided yet.

Tonight he can hear the fifth of a second before \emph{ára} sitting
inside that handwriting. The same placement. The same woman putting a
thing down exactly where she meant it to go, and then not sending it for
forty years.

He does not go up the ridge. He packs the machine instead, and helps
Sigrún down the hill in the dark, which takes a long time and which
neither of them remarks upon, and at her door she says, without any
preamble:

``She was kind to me. When I was small.'' Her breath is short. ``I do
not know what happened to her either. Nobody would ever say. But she
used to carry me on the boundary walk when my legs gave out, and she
sang, which nobody does.'' She fumbles the latch. ``I thought you should
have one thing about her that is not in a book.''

\chapter{}

The fog does not lift; it only dilutes, thin enough to admit solid
bodies while keeping the light for itself. Rowan stands at the Annex
threshold, watching the coastal track. The landscape is monochrome ---
black basalt, grey scree, white mist --- and then a colour error
appears. A figure climbs the incline, not with the villagers' low
hip-swaying negotiation of gravity but with vertical purpose, stepping
over stones rather than around them. The jacket is chemical yellow.
Gore-Tex rasps against itself, a sound too sharp for this place.

He wipes the condensation from his glasses and measures the arrival. Two
minutes. The figure resolves into a woman with a padded camera bag heavy
with lenses and a smartphone held aloft, hunting a signal that died ten
kilometers south. She stops ten meters off and exhales, and the vapour
disperses at once, where the locals' breath clings.

``Dr.~Hale?'' Her voice is clear and high, stripped of hesitation, the
voice of a presenter. It hits the turf wall and does not resonate; it
bounces.

``Yes. Hale.''

She bounds the last steps and offers a hand in a fingerless glove.
``Brynja Þórhallsdóttir. Thorsteinn said you were up here. I parked by
the quay, or what's left of it. The road situation is a nightmare.'' Her
grip is dry and firm.

He switches to Icelandic, the local dialect weighted on his tongue now
without his choosing it. ``You are Edda's granddaughter.''

``Yes.'' She answers in Reykjavík standard --- clipped vowels, sharp
consonants, clean nominatives without a second thought. ``Haven't been
back since I was twelve. Place hasn't changed, has it? Just sunk a
bit.'' She raises her phone and photographs the grey void where the bay
should be, the artificial shutter absurdly loud. ``I'm here because the
project moved up. The council thinks your village will be gone by next
winter. They want a full package before the budget closes.''

He hears the pressure under the professional tone: if she comes back
empty-handed, it will not only be the village that disappears from
someone's ledger. ``What kind of package?''

``Oral histories, dialect features, environmental footage.'' She taps
the bag. ``And---'' She hesitates, the first hesitation he has heard.
``---Edda. I need her voice on record. I promised them. I promised
myself, too.'' There it is, he thinks: not grant-speak but obligation
dressed as work, and he knows the costume well, having worn it across an
ocean.

``You mean to film.''

``High fidelity. The stories, the songs, the layout. It's cultural data.
If we don't hold it, it evaporates.'' He looks at the yellow jacket,
fabric designed to shed weather, and thinks that she is dressed like an
astronaut and believes the suit is protection rather than insulation
from consequence.

``The elders are protective of the spoken forms,'' he says carefully.
``They do not welcome copies.''

Brynja laughs, light and airy. ``Oh, I know. Grandma Edda is the worst
for it. `Don't write it down.' `Don't fix it.' Quaint. But folklore dies
when it stays oral. If we don't record it, it's just noise in the
wind.''

The mist curls around her boots and seems to recoil from the synthetic
yellow, as if the fabric refuses to be absorbed. ``The language here
isn't representational,'' he says, lowering his voice. ``It acts.''

She waves a hand, an efficient dismissal. ``Semantics, Doctor. You
linguists romanticize syntax. I need footage and interviews. If I can
get Edda to recite one of the storm forms on camera, the council signs
off and I'm done.'' \emph{Done.} The word lands wrong in his body.
Nothing is done in Brynjavík. Only held.

``I'll walk you down. But I can't promise they'll speak.''

``They'll speak,'' she says, already turning for the path. ``Everyone
wants to be remembered. That's human nature.'' She descends quickly, her
silhouette cutting a hard line through the soft grey. Too fast. Too
clear. A sharp bright thing entering a place that has spent a century
teaching itself to blur.

\sceneline

On the basalt shelf she levels her tripod with a spirit bubble and
mounts a sleek black mirrorless camera, the lens wide enough to swallow
the bay. She checks the exposure. ``The light is flat,'' she murmurs.
``But the texture is incredible. Look how it doesn't quite reflect the
sky.'' He stands three paces back, hands in his pockets, fists clenched.
The tide slides over rock with a silence that feels enforced.

``Rolling,'' she says, and pans from cliffs to horizon. On the little
screen the image is crisp --- basalt sharp, foam clean --- until the
frame crosses the line the villagers will not name. The image does not
blur. It tears. Ribbons of ragged colour shred the horizon, sky and
water swapping places in jagged blocks, as if the machine cannot decide
which is subject and which is object.

``Damn it.'' She stops recording. ``Magnetic interference off the
volcanic rock. Fixable.''

``It's the line,'' Rowan says. ``It isn't fixed.''

``It's a camera, Rowan. It doesn't have a philosophical crisis. It has a
signal-to-noise problem.'' They move inland. A fisherman --- Einar, a
glass eye catching the flat light --- stands by the drying racks. Brynja
smiles, clips a small microphone to his lapel. ``Tell me about the
winter of seventy-nine. In your own words.'' Einar's eye flicks to the
red tally light. He speaks, but his syntax contracts, articles dropping,
verbs shortening, speech pared down to slip past a recording head
without catching. ``Cold came. Sheep died. Ice lay.''

``And the sea?'' she asks, too casually. ``How was the sea?''

Einar freezes, his gaze going past the camera. ``The west,'' he
whispers. ``Was heavy.''

``I'm getting hum,'' Brynja frowns into her headphones. ``Say that
again? The bottom just dropped out of the frequency.'' Rowan hears it
without headphones --- not electrical but subterranean, something
dragging across the floor of the spectrum, lower than any diaphragm
should be able to reach.

``Let him go,'' Rowan says. Einar unclips the microphone with shaking
fingers and walks away as if leaving a line of fire.

``Wind shear,'' Brynja mutters, making a note. ``Tricky environment.''
She heads for the lowest cottage in the scree. Edda's door is shut, the
windows dark though it is midday. ``Grandma will be the centerpiece,''
she says softly, almost pleading now. ``If I can get her to recite the
Sæfangi on record, the model is complete.''

The air tightens. ``Do not ask her for that case,'' Rowan says.

``It's for the archives. It's for her.'' She raises a hand to knock, and
the door opens. Edda stands compressed in the frame --- not frail, dense
--- her eyes fixed on the lens, not on her granddaughter. ``Hi, Amma,''
Brynja says, gentling. ``Just a short session. Basics.''

Edda's gaze stays on the red light. ``You bring a bucket with no
bottom,'' she says in Vikamál, pure and hard. ``You think it holds.''

``It's digital. It holds everything.''

``It holds nothing. It repeats.'' She lifts a hand and draws a single
vertical slash in the air, cut short. \emph{``Slökkva á auganu.''}
\emph{Extinguish the eye.} The door slams.

Brynja exhales and turns back, already reaching for competence. ``She's
dramatic. We'll try again tomor---'' She stops. The camera screen is
black. A moment ago the battery read eighty-four percent; now it flashes
empty red and dies, the lens retracting with a thin exhausted whir.
``Dead,'' she says, popping the battery, touching the contacts. ``Stone
cold. It just drained.''

``The circuit broke,'' Rowan says, looking at the shut door.

``Bad cell. Grid fluctuations. I'll charge everything tonight.'' She
slots a spare and the screen wakes at sixty-eight percent. She points
the lens at her boots, the path, her own hand. The frame holds. She pans
to the dark window and speaks into the microphone, level: ``Test. Sound
check. The sea is heavy.'' The percentage drops as she speaks ---
sixty-eight, fifty-one, thirty-three --- and it is the vowels that do
it; Rowan watches the meter dive on the long open syllables and recover
on the consonants, as if the cell is being asked to hold a note it
cannot describe and is paying for the description in charge. Brynja
tries again, in English this time: ``The water is large. The water is
wet.'' Sentences with no agentive grammar, no force. The percentage
holds at thirty-three. Stable. She powers the camera down before she has
to admit what she has just measured, and packs up. Around them the
curtains settle back into place. The village is not hiding from the lens
because it is shy. It is hiding because once a thing is recorded it can
no longer change its mind, and here, changing one's mind is how you
survive the weather.

\sceneline

Dinner is in the annex of the curing house, the beams low enough to
graze a tall man's crown, the air dense with boiled haddock and rye and
wet wool, thirty people trying not to breathe too deeply. Rowan sits
near the end of the trestle; Brynja sits mid-table, a bright singularity
in a room of earth tones, talking fast --- Ministry stories, funding
loops, the absurdity of the paperwork --- her Icelandic fluent and
unbroken, nominatives and active verbs delivered like small
declarations.

He watches the window behind her. It is old single-pane glass, wavy,
cold, and it does not fog. Instead, along the black iron of the window
latch and the nailheads that hold the frame, a fine white rime begins to
grow as she speaks --- salt, sweating out of the cold metal into the
warm room, crystallizing in threads and hooks. It does not spread
evenly. It builds on the long open sounds and pauses on the stops, a
slow frost of salt writing itself along the ironwork in the exact rhythm
of her sentences, and no one but Rowan is looking at it, and he cannot
say whether it is a warning being written or a receipt.

Brynja laughs. ``The funding model is circular. We apply for a grant to
study the application process. It's---'' A fork clatters to the floor.
The sound severs the room, and the silence drops instantly. Einar
stands, his chair scraping back. He does not look at Brynja; he looks at
the space in front of her mouth, as if tracking a thrown object. ``My
apologies,'' he says, tight. ``The air is full.'' He leaves, and latches
the door.

``Touchy,'' Brynja murmurs. ``I think he hates the funding talk.''

``He broke the sentence,'' Rowan says. ``You were about to give the
water agency.''

``I was using a metaphor.''

``Metaphors here act like instructions.''

After the meal they step out into a cold that holds its breath, and walk
to the tide pools by the breakwater, basins black and still, reflecting
stars without a shimmer. ``Let's get this on record,'' Brynja says,
pulling out her phone, a red dot on the voice app. ``Dr.~Hale. The
synthesis. What's the core mechanic?''

He hesitates. The pool reflects his face. ``It's a system of
avoidance,'' he says. ``Grammar removes the speaker from the path of the
subject.'' He hears the phrase twice --- once from his mouth, once from
the water. \emph{\ldots path of the subject\ldots{}} He freezes. The
second voice did not come from the phone. It came from the pool.
Concentric rings spread across the surface though nothing has struck it,
and the voice rises flat and wet, perfectly timed.

Brynja laughs a little too fast. ``Echo off the cliff. Wild.''

``That's not an echo. The timing is wrong.''

``It's acoustics.'' She steps closer. ``Hello?'' she calls. The pool
ripples again. \emph{\ldots Hello\ldots{}} Identical. Not mimicry. The
sound returned whole, unopened, as if it had been kept a moment and
handed back. Her face loses its confidence for the first time, and she
looks down as though the pool has grown an eye. ``Okay,'' she says,
forcing the smile back. ``Let's check the video from today.'' She opens
her gallery and plays the drying-racks clip. The edges smear where the
compression failed; the villagers stand tense and unhappy; and behind
them, blurred by the depth of field, other figures stand --- pale,
indistinct, like smudges on the lens that have somehow learned posture.

``Zoom,'' Rowan says. She pinches the image, and the figures mirror the
villagers exactly, hunched shoulders echoed, crossed arms echoed. Not
ghosts. Not people. Grammar given a body, the other half of the verb
standing patiently in the frame, waiting to be conjugated.

``Glitch,'' Brynja says, too loudly. ``I need to reformat the card.''
She shoves the phone away as if hiding it makes the image less true.
``I'm going to the hall. I need to upload. I need signal.'' She walks
off fast, not looking back, and behind her the pool repeats
\emph{signal, signal, signal} in smaller and smaller versions into the
empty rock. Rowan watches her go, and notices, as she climbs, that the
path behind her looks subtly flatter than the path in front, less
textured, the way an over-exposed edge of a photograph loses its detail.
Not gone. But edited.

\sceneline

Sleep is a hypothesis he cannot prove. The Annex is silent, but the
silence presses, active and physical, and at a quarter to one he leaves
the bed and steps into the night. Fog hangs low, bandaging the village.
He navigates by his memory of the ground, and one light burns on the
north ridge: Edda's window. He stops at the dry-stone wall five meters
from the glass. The curtains are drawn but not closed, a narrow slit
showing candlelight, and Edda at her table --- not the iron pillar of
their first interview. She is trembling. Her hands are raised, fingers
splayed, moving through complex tracing geometries, syntax drawn into
the air. He recognizes the motions: the downward stroke, the enclosing
circle. She is building a sentence to brace the room. Her lips move; he
cannot hear the words through the glass, but he can see the effort, the
jaw tight, the neck cords straining, the labour of lifting an invisible
weight.

She reaches for the binding, the Sæfangi. Her hands move to form the
knot. And she falters. Her left hand spasms; the circle breaks; she
tries to retrace it and her fingers are stiff, arthritic, exhausted,
overpowered. Her mouth shapes the long \emph{Á.} Her breath catches. She
collapses forward onto the table. The candle flame lengthens, turns
blue, and snaps audibly sideways, as if pressed by a heavy hand. Edda
does not rise. Her cheek is on the wood, her hands open. The sentence
remains unfinished.

Rowan backs away and looks toward the sea. The tide is high --- too high
--- lapping at the base of the ridge, meters above any mark that should
exist. And it is not making the sound of water. It is making the sound
of articulation. \emph{Thhhh. Uuu. Mmm.} Breath through teeth,
consonants grinding against stone vowels. The sea is rehearsing, trying
to pronounce the thing it touches. He turns and walks faster toward the
village, the gravel loosening underfoot as if definition itself were
draining out of the ground, and passes the community hall, where light
glows from the office, and stops.

Brynja is inside, hunched over a laptop in a nest of drives and cables,
the blue screenlight hollowing her face, sweat shining on her upper lip.
She is pressing Delete, over and over. He moves to the glass. The files
corrupt on the screen, but not into random noise --- the faces melt into
long jagged lines that resemble the scripts carved into the schoolhouse
beams. A frame freezes: the shoreline, but the rocks have rearranged
themselves into text, the pixels declining into cases. Brynja's mouth
forms a small sound he cannot hear, but he can read it: refusal, denial
breaking. She scrubs the record as if scrubbing could undo the act of
having looked.

Rowan stands between the two windows. Behind him the sea practices its
new vocabulary against the ridge. To his left, the last speaker lies
defeated by the grammar she spent her life holding. To his right, the
archivist is being overwritten by the very data she came to save. There
is no neutral ground on which to stand and watch. The translation has
begun, and the language being spoken no longer requires a human tongue
--- though, he understands, standing in the cold between the windows, it
would still very much prefer one.

\chapter{}

The body of Magnús Jónsson lies on the shingle three meters above the
high-tide mark, arranged in a posture of delivery rather than collapse,
arms extended toward the village, palms open as if offering an
explanation. The gravel around him is dry. The wool of his sweater is
dry, still holding the smell of lanolin and stale tobacco. Even his
boots --- heavy leather workwear --- show no darkening from brine. He
has not been in the water.

Yet when Arnar kneels and presses two fingers to the sternum, the sound
that comes from the open mouth is the gurgle of a rising tide. Arnar
withdraws his hand and wipes his fingertips on his trousers. He does not
close the eyes, which are fixed on the grey zenith of the sky, the
corneas clouded with a white crystalline precipitate. ``Full,'' Arnar
says. He does not say \emph{he is full.} He leaves the subject out,
keeping the condition from fastening to the man.

The villagers stand in a semicircle, heavy coats, silence, a wall placed
deliberately between the body and the waterline. They do not look at the
sea. Their backs are rigid barriers shielding the evidence from the
plaintiff. They cluster tight, closing the gaps where a line of sight
might slip through. Rowan stands at the periphery, observing the
geometry. In any other jurisdiction this is a crime scene, or a medical
emergency. Here it is a matter of contested boundaries.

He knew Magnús a little. Not well --- no one lets an outsider know them
well here --- but he had sat across a table from the man twice, had
watched him mend a creel with the same economy Arnar brought to a knot,
had been given, once, a piece of dried fish and a nod that meant
\emph{you are permitted to be here.} Now he looks at the salt crust
forming on the dead lips and finds that the clinical phrase will not
come clean. He writes it anyway, because the writing is the last
handhold he has. \emph{Internal inundation. No external contact.} And
then, because it is true and the true words are the local ones now, he
writes what the village would write. \emph{He was called, and the water
in him answered.}

A gap opens in the defensive line, just wide enough for Edda. She moves
with a friction that suggests the air resists her passage. Her stick ---
driftwood capped with iron --- strikes the stones in a slow click. She
stops at the body's feet, looks at Magnús, at the villagers, briefly at
the horizon they are trying to obscure. She leans down, touches the salt
on his lip, rubs the crystal between thumb and forefinger, testing the
grit. ``He was named,'' she says. Not a whisper. It lands like a gavel.

Arnar nods. ``In the night. We heard nothing.''

``The sea does not need to shout when it has the paperwork.'' She
straightens, her spine locking into place with audible effort. ``It
found a handle. It called him, and the water inside him went to meet the
water outside, the way a level finds a level.''

Rowan steps closer, pulled by the logic of the pathology despite
himself. ``It filled him from the inside.''

Edda turns her head. Her eyes are opaque, reflecting the dull dawn.
``The water is not a substance, Dr.~Hale. It is a location. Magnús was
reclassified. He was moved to the location of the sea without his body
leaving the stones.'' She looks back at the corpse. The death is not
tragic to her because Magnús suffered. It is tragic because the boundary
of the village has been successfully contested, and everyone standing
here now lives a little closer to the edge than they did yesterday.

``The claim is invalid,'' Arnar says, though his voice lacks conviction.
``He was on the roster for the winter shear.''

``The roster is paper. The sea has precedent now. It has taken a subject
without wetting the ground. If we do not answer the claim, it will take
the next one by simply thinking of him.'' She taps her stick against the
stones. One. Two. Three. ``We must unbind it. The Aflending. Before the
next tide makes the record permanent.''

A ripple of unease moves through the cluster. The Aflending is not a
prayer; it is a retraction, a dangerous maneuver that requires the
speaker to stand in the void between noun and verb and hold it open.
``Is the voice strong enough?'' Arnar asks, his gaze flicking to her
throat, the cords standing out like rigging under strain.

She does not answer directly. She looks at the water lapping innocently
at the quay. ``The voice is the tool we have. We do not ask whether the
hammer is heavy. We ask whether it strikes.'' She turns her back on the
body. ``Cover him. Do not take him to the church yet. To bury him now is
to accept the death as natural. We contest the cause first.''

Séra Hjálmar has come down the shingle in a waxed jacket over his
collar, and he stands at the head of the body with his hands at his
sides, and he does not touch it and does not pray over it, and when Edda
gives her instruction about the church he does not contradict her.

Rowan moves along the line to stand beside him. ``You're not
objecting.''

``To what? She has not said he may not be buried. She has said not
today.'' Hjálmar's voice is very low. ``I have known Magnús Jónsson for
nineteen years. He mended my gutter twice and would not be paid. I will
bury him properly on a day of my choosing and I will use every word of
the office.''

``And in the meantime he lies under a tarp.''

``In the meantime he lies under a tarp,'' Hjálmar agrees, ``with stones
on the corners, in the parish of which I am the incumbent, and I have
been asked to allow it by a woman of eighty who has never once asked me
for anything, and I have allowed it.'' He looks at the covered shape.
``You will want to know what I think happened.''

``I would.''

``I think a man died in the night and there will be a physical cause and
a doctor from Húsavík will find it or not find it, and I think whatever
he writes on the form will be true, and I think it will not be the
whole of the truth, and I have no method for the remainder and neither
does he.'' A gull goes over. ``That is not evasion, Dr.~Hale. That is
the actual position and I have held it for nineteen years and it is more
uncomfortable than either of the two comfortable positions available.''

The villagers move at once. A tarp --- heavy boat canvas --- is drawn
over Magnús, hiding the salt-filled eyes and the dry wool, and basalt
rocks weight the corners, pinning him to the land, preventing any drift
into the domain of the sea. There is no weeping. They are securing a
perimeter. As Edda passes Rowan she smells of ozone and old paper, and
she does not acknowledge him, but he hears her muttering, rehearsing,
testing the tensile strength of the syllables she will need.
\emph{``\ldots unda\ldots{} und\ldots{} undi\ldots{}''} She ascends the
path to the headland alone, and the villagers watch her go with the
terrified hope of people who have handed a losing case to the only
advocate they have.

\sceneline

The basalt shelf thrusts out into the grey slurry like a lectern, the
rock slick with spray, black and geometric, a platform built for
address. Rowan wedges himself into a fissure twenty meters back, hidden
by dried kelp and rockfall, holding the recorder close, shielding the
microphone with a gloved hand. The device is running. The levels are
low, capturing the ambient hiss.

Edda stands at the edge. She does not raise her arms; she does not
kneel. Feet shoulder-width, she braces against the drag of the current,
a vertical line drawn against the horizontal mass of the water. She
begins with no preamble. \emph{``Haf\ldots{} unda\ldots{} straum\ldots{}
und---''} The sound is mechanical. It grinds. The suffix \emph{-und} is
a cancellation, a brake applied to the noun's momentum, and she speaks
it with flat percussive force, hitting the final consonant hard.
\emph{``Salt\ldots{} unda\ldots{} dýpi\ldots{} und---''}

Rowan watches the water. It does not crash. It listens. The waves
approaching the shelf do not break; they pause, the foam dissolving back
into the dark mass, the level at the base of the rock beginning to
lower. It is not a tide going out. It is a margin being redrawn. The sea
pulls back from the stone in a clean concave arc that mirrors the radius
of her voice, and the wet rock dries instantly, the moisture sucked back
into the body of the water as if ordered home. Edda presses forward,
stepping onto newly dry stone, reclaiming territory.
\emph{``Nafn\ldots{} unda\ldots{} Magnús\ldots{} und---''} She fixes the
unbinding suffix to the dead man's name, annulling the citation.

Her voice holds, but it lacks resonance now, metal on metal. The effort
of holding the grammatical negative against the physical positive of the
whole ocean is enormous; Rowan sees her shoulders shaking, not with cold
but with load. \emph{``Alda\ldots{} unda\ldots{} blá\ldots{} undi---''}
The water recedes further. A jagged reef, unseen for decades, breaks the
surface, black mussels clicking softly in the air. The sea concedes the
point, backing down before the precedent she cites.

Then fatigue arrives --- not a stumble of the feet but a stumble of the
tongue. She reaches for the compound \emph{strand-lengja,} coastline,
needing to unbind the coast from the claim. \emph{``Strand---''} Her
voice catches. The long vowel fractures. She tries to engage the suffix.
\emph{``Strand\ldots{} un\ldots{} daa---''} The rhythm breaks. The final
vowel lengthens too much, the precise termination going soft, almost a
question. The sea stops retreating. The arc of water freezes, holding
its position, verifying the silence.

Edda knows. She tries to correct, to force the consonant back into
place. \emph{``Und. Und.''} But the structure has failed; the lattice
that held the water at bay snaps under the fatigue of seventy years of
maintenance. Her mouth opens and no sound comes. The muscles of her jaw
work furiously, trying to shape the case, but the bridge between intent
and execution is gone. The sea waits one beat, confirming the gap, and
then resumes --- water sliding back over the reef, covering the dry
stone, filing back into the emptied space without hurry, approaching her
boots. She stumbles back, her stick clattering, one hand to her throat,
looking at the water with the confusion of a mathematician who has
forgotten how to count. The water stops. It does not reclaim everything;
it halts two meters short of its old high mark, conceding the immediate
surge while keeping the jurisdiction.

Rowan stops the recording. Sickness rises in him, a sympathetic
resonance with the failure he has just watched, and under it something
worse: he has watched a woman spend the last of a voice, and he did
nothing, and could have done nothing, and part of him was taking levels.
Edda turns and sees him, or the shape of him in the rocks, and does not
scold, and does not speak. She walks past, her gait uneven, her mouth
slightly open. As she passes he hears her trying to whisper, trying to
say \emph{the water,} and cycling the endings, mismatching them, calling
the water \emph{waters,} calling it \emph{to the water,} unable to find
the subject. She has lost the ability to name the thing that defeated
her. He stays crouched until she is gone, watching the sea settle into
its renegotiated line, lapping against the stone with the soft wet sound
of a page being turned.

\sceneline

The doctor comes out from Húsavík on the Thursday and is on the shingle
for forty minutes and gives the cause as drowning, and on the evidence
in front of him there is no reason to give any other, and Rowan --- who
stands at the edge of the examination with his hands in his pockets ---
would have written the same.

What Hjálmar does with it is the thing Rowan will remember.

He shows him that evening in the vestry, without being asked: the parish
return, the same form filed in the first week of every January since
1732. A column for the year. Columns for baptisms, confirmations,
marriages, burials. And a narrow column at the right, headed in an
abbreviation four centuries old, for the circumstance.

``I will put him in the burials,'' Hjálmar says. ``That is a fact and I
will state it.''

``And the circumstance?''

``I will leave it blank.''

Rowan looks up.

``I cannot write \emph{drowning.}'' Hjálmar sets the pen down. ``He was
dry. I was there. I looked at his boots and at the gravel under him and
there was not one wet stone, and I am not going to put a word into a
book that will sit in that cupboard for three hundred years and tell
somebody in 2325 a thing I know to be false.'' He caps the pen. ``And I
cannot write the other thing, because I do not know the other thing ---
and because if I write it, the diocese sends somebody. That somebody
will be kind. He will conclude that a rural incumbent has become
isolated. He will be gentle about it and he will be entirely wrong, and
I will be replaced by a man who drives out on alternate Sundays.''

``So you leave a hole.''

``I leave a hole.'' He turns the volume toward Rowan and puts a finger
against the column. ``It will be the fifth in two hundred and ninety-
three years. There are four already. Seventeen eighty-three. Nineteen
eighteen. Nineteen thirty-four. Nineteen ninety-one.'' He takes the
finger away. ``I have looked at all four and I have looked hard. In
every case the burial is recorded. In every case the circumstance is
not. And in every case the hand is perfectly steady --- no blot, no
hesitation, no scratching out. Four men who knew exactly what they were
declining to write, and wrote everything else beautifully.''

Rowan looks at the four years.

Seventeen eighty-three is the ash, and a harbour that walked east.
Nineteen thirty-four is the winter Edda's mother was nine and the
village doubled. Nineteen ninety-one is Þórunn, and four syllables, and
an eleventh stone that has been cleared and checked and repacked every
month since by people who cannot speak to it.

He does not say any of this.

``I find I am not lonely about it in the way I expected,'' Hjálmar says.
``That is the odd part. Four men across three hundred years, none of
whom I can ask anything, and I know precisely what each of them was
doing on the day he picked up the pen.''

Rowan looks at the space that is not yet blank because it is not yet
written, and understands that he is watching this village's record being
kept the way everything else here is kept: not by what is set down, but
by a three-hundred-year refusal to put the wrong word in.

\sceneline

The village contracts. Not a retreat --- an editing. By noon the three
families on the western perimeter, nearest the old salting sheds, begin
to move. No trucks, no shouting. They carry boxes and furniture by hand,
walking the centerline of the road toward the cluster of houses near the
hall. A man carries a kitchen chair. A woman follows with folded linen.
They do not look back. They do not speak of the house they leave; to
speak of it would draw attention to a vacancy, and a vacancy is a space
the sea can fill.

Rowan watches from the Annex window. As he watches the abandoned house,
its edges lose definition. The roof gutter blurs into the sky. The
corners soften, less stone than smudge. It becomes difficult to hold in
focus; his eye slides off it, drawn back to the occupied houses at the
center. The abandoned house is becoming a dropped subject, slipping
quietly out of the village's sentence.

Brynja sits on the hall steps in a triage of hardware --- a satellite
phone, a laptop, her mobile. He approaches. ``The uplink is solid,'' she
says without looking up, her fingers hammering keys. ``Four bars on the
sat-phone. Everything works.''

``But no connection?''

``Oh, it connects.'' Her voice is tight with the rage of an
administrator denied a receipt. ``It pings the satellite. The handshake
initiates. And then---'' She turns the screen to him. A single line:
DESTINATION UNKNOWN. ``It can't route the call. Not a dead line. The
routing tables can't find Brynjavík. I tried emergency services. I tried
the coast guard. The automated system asks for my location, and when I
give the coordinates, it hangs up.'' She dials again, listens.
``\,`Invalid format,'\,'' she narrates, repeating the robotic voice.
``\,`Please enter a valid geolocation.'\,'' She walks to the edge of the
gravel, holding the phone up as if height could cure the problem. ``We
are right here,'' she tells the device. ``Sixty-four degrees north. We
are on the map.'' The silence from the speaker is absolute. ``It's a
database error,'' she snaps. ``Some server in Reykjavík wiped the
district codes. It has to be.'' She lowers the phone and looks at him,
pale, her makeup stark against her skin. ``I can't file the evacuation
request if the system doesn't recognize the origin point. It's like we
don't exist anymore.'' She says it as a clerical grievance. She does not
hear how accurate it is. She throws the phone; it arcs and thuds into
the grass without shattering, the screen dark. ``I'll try the radio in
the harbour master's office,'' she says, smoothing the yellow jacket.
``Analog might bypass it.'' She marches off, sure this is a procedure
and not an erasure.

Rowan walks the new perimeter. The boundary has moved inward by fifty
meters. Arnar kneels in the dirt near the path to the Annex, setting a
stone into the earth, rough basalt, the face freshly chiseled. Rowan
reads the carving: \emph{-um.} The dative plural. ``A new declension?''

Arnar packs soil around the stone. ``We are shortening the sentence. The
previous clause was too long. Hard to defend.''

``You're ceding the outer ring.'' Rowan looks past the marker to where
the path that led to the western ridge used to be. It is not overgrown,
not blocked. It is simply not there. The eye expects a track, a
depression in the grass, and the grass is unbroken. The landscape has
healed over the memory.

``We do not cede,'' Arnar says, standing, testing the stone with his
boot. ``We punctuate. This is where the village ends now.''

``And what lies beyond it?''

Arnar looks at the space where the path used to be, at the blurred
outline of the abandoned house. ``There is nothing beyond,'' he says.
Flat, final. Not describing a void but enforcing one; to say \emph{there
is nothing} is to make it so. Rowan looks at the nothing and tries to
retrieve yesterday's sheep pen --- grey wood, rusted wire --- and cannot
place it in the current terrain. The reference points are dissolving.
``It is safer,'' Arnar says, lifting his tools. ``A small room is easier
to heat. A short sentence is easier to speak.'' He walks away, keeping
strictly within the new markers. Rowan stares at the invisible line,
then steps one foot over the stone, onto the grass of the nothing. The
ground feels spongy, insufficient. The air is colder here, and a silence
presses against his ears like deep water. He pulls his foot back, and
retreats to the side of the stone that still has a name.

\sceneline

The desk lamp casts a cone of yellow on the blank page. Rowan grips his
pen, meaning to write a summary of the morning. \emph{The discovery of
the body,} he writes, and stops on the word \emph{body.} It feels
insufficient --- light, hollow, a shell with the meat taken out. It does
not carry Magnús's weight on the stones; it does not account for the
salt in the lungs. He crosses it out. \emph{The discovery of the
vessel,} he tries. No --- \emph{vessel} implies utility, and Magnús was
not used, he was misfiled. His hand hovers. A tremor starts in his
wrist. He searches the whole of English for a word that means a
container emptied of its self and filled with a hostile geography, and
there is no such word. But there is a sound rising unbidden at the back
of his throat. \emph{Sæfangi.} The sea-prisoner. The one held by salt.
He writes it down. \emph{Sæfangi found at 06:00.} The Icelandic
characters look sharp and distinct on the page; the English around them
looks vague, the ink bleeding into the cheap stock.

He tries to continue. \emph{Edda performed the ritual.} His hand locks.
\emph{Ritual} is wrong; it implies symbol, and what happened on the rock
was mechanical. He searches for the word for the legal retraction of a
weather event, and it is \emph{Aflending,} and the pen moves almost on
its own. \emph{Edda initiated Aflending.} He drops the pen and pushes
the chair back, and the panic that arrives is not adrenaline but a slow
cold understanding: the shield of English, the professional distance he
came here wearing, is failing, because it cannot describe Brynjavík ---
it lacks the cases. And he is not surprised, and the lack of surprise is
the most frightening thing of all, because it means some part of him
always knew that English was a language he had chosen, and that
underneath it, waiting, was a first language he had spent thirty years
refusing to learn --- a mother tongue, in the most literal sense, the
tongue of a mother who kept her words behind a wall.

He paces. He needs an anchor. He needs to say one complex English
sentence to prove he still owns his own mind. ``The sociolinguistic
parameters of the village are undergoing a rapid contraction,'' he says
aloud, and stumbles on \emph{contraction,} the word going mushy in his
mouth. ``The --- the shrinking,'' he tries, and the Vikamál arrives with
a physical rightness, clicking into the slot \emph{shrinking} could not
fill: \emph{sam-drættir,} the drawing-together, the pulling-in of nets.
The displacement is aggressive. The dialect is not merely vocabulary; it
is colonizing his thought, replacing abstraction with the concrete and
the ritual. He is no longer thinking \emph{about} fear. He is thinking
about vibration against glass, about a wall of salt writing itself along
an iron latch.

He turns to the window. The night is not uniform; there are grades of
blackness. At the property edge stands the turf shed where he keeps his
bicycle, and he can see its outline, but the outline is flickering. One
moment a solid block of shadow; the next, the darkness flows through it,
uninterrupted. He blinks and tries to focus. \emph{It is a shed,} he
thinks, deliberately, in English, in the nominative. \emph{The shed
stands.} And the shed vanishes. No collapse, no fire. It simply ceases
to be a reference point, the space it occupied continuous now with the
empty field, the eye sliding over it without purchase. He gasps, and
feels the hole in the landscape as a hole in himself. He backs into the
desk and the notebook falls open to his scrawled notes --- English stems
weak and dying, Vikamál endings sharp and black and authoritative, a
page at war with itself.

He needs to name it before he loses the ability to name it. ``Linguistic
erosion,'' he whispers, forcing the phonemes, and his tongue rebels,
curling back, refusing the soft Latinate flow, wanting the grinding
consonants of the coast. He tries again. ``Linguistic---'' His mouth
shapes the vowel, his throat constricts, and what comes out is deep and
guttural and terrifyingly precise. \emph{``Tungu-taka.''} The
Tongue-Taking. The word tastes of brine and resonates in his chest,
vibrating his ribs, heavy and solid and true in a way \emph{erosion}
never was. He closes his mouth. He stands in the room's silence, feeling
the vibration settle into bone, and he knows now that he is not
observing the collapse. He is inside it. And under the fear, quiet and
unwelcome, is something almost like relief --- the relief of a man who
has been speaking a foreign language politely for thirty years and has
just, at last, been addressed in his own.

\chapter{}

The village of Brynjavík is smaller today than it was yesterday. This is
not a trick of light, nor of fog pressing on the cliffs. It is
measurable. The path to the northern ridge, once wide enough for a
truck, now barely admits the width of Rowan's shoulders. The turf houses
on the periphery have been edited out, their mass absorbed into the grey
vacancy of the moor, and the remaining structures huddle together like
teeth in a shrinking jaw. He walks up the incline with his eyes on his
boots, measuring the distance between stones, because the geography is
under revision and he no longer trusts it to hold still. To his left
Arnar kneels in the mud, wrestling a boundary stone into a fresh
depression five meters inward from where it stood, packing wet soil
around its base with the flat of a shovel, punctuating the village,
shortening the sentence to keep it defensible.

Edda's door is unlatched. It swings open with a weight that suggests the
hinges are rusting in real time. Inside, the electric lights are dead;
the room is lit by tallow candles set on the floor, their flames
perfectly vertical in the still air. The iron stove throws intense heat,
yet she sits in her chair wrapped in layers of heavy wool, and she looks
like a bird dried in salt wind --- light, brittle, made entirely of
intent. She does not greet him. She lifts a hand, the skin translucent
as parchment, and beckons him to the stool beside her. Her eyes are not
cloudy. They are terrifyingly clear, stripped of nostalgia, the eyes of
someone preparing to close an account.

``Sit,'' she says. The verb is stripped of courtesy, a placement
command. He sits, and out of reflex reaches into his bag for the
recorder, checks the card, sets it on the small table, presses power.
The screen flickers and fails to wake. The red light pulses once, a weak
starving heartbeat, and dies. He reaches for the backup camera. Nothing;
the lens stays retracted, the shutter locked. The devices sit on the
wood, inert. They have not broken. They have been dismissed. The room
has narrowed its jurisdiction to exclude them, because they are not
permitted to witness this, and so they cease to be witnesses.

``Leave them,'' Edda says, not looking at the machines. ``Paper cannot
hold this. Silicon cannot hold this. Only the meat holds it.'' She leans
forward, smelling of old books and dried kelp, and catches his gaze and
holds it. There is no affection in the look, only assessment --- she is
checking that the vessel is sound before she pours. ``Before I give you
the weight,'' she says, ``I will give you the reason. You have carried
it long enough not to know it.''

She reaches into the folds of her shawl and brings out the envelope. His
envelope. His mother's hand on it, \emph{for E.B. only,} the paper
softened now from being held. She has kept it, all these weeks,
unopened, against her skin.

``You slid this across my table the first day,'' she says. ``You thought
you were being a scholar. You were being a son.'' She turns it over in
her fingers. ``Elizabeth. Beth. She was Bryndís, before Cambridge,
before she made her name flat enough to fit on an English door. She was
born in the house that used to stand where the western nothing is now.
She and I learned the grammar at the same table, from the same woman, in
the same winter. We were not sisters. This place has no word for sisters
that is safe to say near the water. But we were as close as two people
are allowed to be named here, which is to say we were never named at
all. We simply held the same silence for twenty years.''

Rowan does not move. The candle flames stand straight up. Somewhere
under his sternum a door he has kept shut his whole life is being opened
by a stranger with a key he did not know existed.

She is quiet a moment, gathering something that costs her.

``You have been asking the wrong question since the road,'' she says.
``You have been asking what this place is. You should have asked what it
is the last of.''

Rowan says nothing.

``There were others.'' Her hand moves on the arm of the chair, laying
them out, one finger at a time, the way a woman counts stitches. ``Not
many. Six that were spoken of when my grandmother was taught, and she
was taught by a woman born in 1851, so take it as six and do not press
me on the number. Six harbours along this coast and the coast opposite,
where the grammar was kept. Not the same grammar. Cousins of it. A
speaker from here could have made herself understood at Vestrahöfn in a
morning and would have needed a winter to be trusted.''

``What happened to them?''

``What happens.'' She says it without weight. ``One went under the ash
in 1783 --- not the fire, the ash; the ash came and the sheep died and
the people walked east and the grammar does not travel, it is fixed to a
water. One modernized. That is the word they used at the time and it is
still the correct one. They put in a road and a cannery and a school
that taught from the capital, and they were prosperous for forty years
and the bay took them in a single night in 1934, and it is in the
records as a wave, and it was a wave.'' A pause. ``One emigrated
entire. Manitoba. Every soul. They wrote back for two generations and
then they stopped, and what they were writing by the end was Icelandic
and not the other thing, so it does not matter when they stopped. And
one --- '' and here her mouth thins --- ``one simply put it down. No
disaster. No emigration. A generation decided it was a superstition and
their children agreed with them and their grandchildren did not know
there had been anything to agree about. The place is still there. You can
drive to it. There is a petrol station. The sea took the harbour wall in
the eighties and they applied for a grant and rebuilt it and it went
again, and they applied again. They are still applying.'' She looks at
him. ``That one frightens me more than the ash.''

``And Brynjavík.''

``Brynjavík is the last one.'' She says it flatly, without any pride in
it at all, and that is what makes it land. ``Not because we were
faithful. Not because we were chosen. Nobody chose us and there was
nothing to be faithful to --- it is a maintenance schedule, Dr.~Hale,
not a covenant. We are the last one because we did not stop. That is the
whole of the distinction. Five harbours had reasons to stop and took
them, and every reason was a good one, and here we simply did not have a
generation that had a reason good enough at the moment it would have
counted.'' Her hands go still. ``It is luck. It is a hundred and fifty
years of a kind of luck that looks exactly like stubbornness from the
outside, and I have spent my life being thanked for it, and it is luck.''

Rowan thinks of the district survey with its cream void, of a library
catalogue in Reykjavík listing seventeen holdings, of the six that could
be pulled from the shelf. He had assumed, for weeks, that Brynjavík was
an anomaly. It is a survivor. Those are entirely different objects, and
he has been describing the wrong one in every note he has taken.

``The old grammar goes down through the women of one line,'' Edda says.
``It was to have been hers. She had the throat for it --- better than
mine. The teacher chose her. And she stood on that shelf one autumn,
seventeen years old, and she looked at what the case does to the one who
speaks it, at the hands, at the years, at the cost, and she decided she
would rather have a life than a boundary.'' Her mouth moves, not quite a
smile. ``She chose the wall over the net. She left, and she never came
back, and she never wrote --- except once.'' She sets the envelope on
the table between them. ``She wrote this, and did not send it, because
to send it would have been to admit she had a debt here, and she had
spent her whole second life not admitting it. She raised you inside that
not-admitting. You think you became a man who studies withheld speech by
accident. You became it because you were raised in the purest example of
it in the world.''

``Why did she keep it,'' Rowan says, and his voice does not sound like
his. ``If she never meant to send it.''

``Because a debt does not disappear when you refuse to pay it,'' Edda
says. ``It only waits for someone in the line to carry it back. She
could not come. So she made you, and she wrote your vocation into you
without telling you, and she died, and you looked in a tin box and found
a coastline you had never seen and thought it was your own idea to
come.'' She looks at him with something that is almost, but is not
quite, pity. ``Your mother sent you, Dr.~Hale. Forty years late. She
sent you to finish the sentence she walked out of.''

He should feel used. He waits to feel used. What he feels instead is the
strange loosening of a knot he has been carrying so long he mistook it
for the shape of his own spine --- the first time in his life that his
mother's silence has had a reason he can hold in his hands, the first
time the wall has had a door in it, even a door opening onto grief.

``You will not open it now,'' Edda says, and it is not a suggestion.
``You are not finished enough to read it. When you can hold the boundary
alone, you will be able to hold what she wrote. Not before. Keep it
against you.'' She closes his fingers over the envelope with a hand like
cold twine. ``Now. Listen. This is the part the machines wanted.''

She begins.

The sentence starts in the nominative --- a naming of the self, but not
her self, the Speaker as a position, a location in the coastline's
geometry, the sound low and vibrating in the chest, establishing the
subject that will bear the load. Rowan's breath catches. He tries to
catalogue the phonemes and they slide past analysis and hook directly
into the nervous system. She moves to the accusative; the direct object
is the coastline, and she binds the verb to the rock, defining the
action of holding, drawing tension between speaker and land --- not
description but a bolt tightened, keeping cliff attached to earth. She
shifts to the dative, the recipient, and hands the indirect object, the
village's whole debt, to Rowan, and he feels the pressure settle onto
his shoulders, a sudden gravity that makes the stool creak. The genitive
follows --- possession, of the sea, of the storm, of the dark ---
defining what owns the village, the vowels lengthening, circling,
enclosing him inside the claim.

Then the forbidden cases. The Aflending, the Unbinding, in which she
speaks the cancellation of her own claim, dissolving the contract that
has held her body in this chair for eighty years, the syllables sharp,
cutting the threads that tie breath to lung. The Sæfangi, the
Sea-Holding, the lock --- a sequence of glottal stops and hard
consonants that verify the boundary and seal the gate and place the key,
not as metaphor but as grammar, in Rowan's hand; he feels it turn shut
somewhere behind his eyes. And finally the Hrimfall, Winterfall, the end
of the cycle --- the last syllable a whisper that carries the mass of a
collapsing glacier, a full stop set at the end of a lifetime of
sentences.

The room goes silent, instant and absolute. Rowan sits paralyzed, the
whole sequence burning in him, not as words he could repeat but as
instructions he cannot forget. He knows the shape of the sentence. He
knows it is perfect. He knows it cannot be written.

Edda exhales --- a long steady release, no rattle, no struggle. Her
chest lowers and does not rise. Her eyes stay open, but the intent
behind them is gone. Her face is smooth; the lines that defined her, the
furrowed brow, the set jaw, have slackened. She does not look like
someone sleeping. She looks like a woman who has set down, at last, a
weight she was never permitted to describe.

Rowan sits in the candlelight, the heat ticking in the stove metal, the
envelope in one hand and the grammar in his chest, and he does not weep.
He is too heavy to weep. He looks at the dead recorder and the dead
woman, and he thinks, with a clarity he will spend the rest of his life
failing to soften: she is the second woman he has watched pour herself
out entirely so that a room could keep standing, and he is beginning to
understand that the rooms are always kept standing by the ones who are
used up, and that he is the one still breathing, and that this is not
the same as being the one who was strong.

\sceneline

The basalt platform is a stage set for a tragedy the audience has
watched too many times. The villagers assemble under a sky the colour of
bruised iron. There is no coffin --- Edda's body has been committed to
the earth in the high cemetery --- but the grammatical committal happens
here, facing the plaintiff. Rowan stands at the rear and counts
sixty-eight heads, and knows the figure without having to work it,
because he has been counting this village since the road, and because
the number that is not here is the one he counts by now: eighty-nine at
Magnús's committal, sixty-eight at Edda's, and between the two only a
season. No one has left by the road. They have left by the other
door --- the houses on the western perimeter that stopped being spoken,
and then stopped being houses, and then stopped. They do not form the
usual crescent at the cliff lip;
they cluster near the path, leaving a twenty-meter gap between
themselves and the drop. A tactical retreat. They know the defense has
weakened.

Arnar stands alone in the buffer, in a suit that fits poorly, the fabric
strained across the shoulders, the cuffs frayed. He does not have Edda's
posture. He stands like a man expecting a blow. He begins the
recitation, his voice rough, abraded by salt and sorrow, and he rushes
the vowels, clips the long sounds the sea requires, a shorthand version
of the rite, summary and not performance. \emph{``From the bone\ldots{}
the salt. From the breath\ldots{} the wind.''} The villagers bow their
heads and accept the abridgment. Arnar is a harbour master, not a
speaker; he knows knots, not declensions. Rowan hears the error as a
physical dissonance in his chest --- Arnar uses the standard dative
where the instrumental is required, treating the sea as destination
rather than vehicle.

Arnar reaches the crux, the Sæfangi, the binding that keeps the name
from being absorbed. He inhales, opens his mouth to shape the archaic
guttural sounds Edda spoke in the cottage. \emph{``Hérna fer\ldots{}
Edda\ldots{} til---''} He stops. His throat works. The silence
stretches, taut and vibrating. The wind picks at his jacket. He cannot
find the ending. The shape of the word is gone, washed out by grief, or
by simple lack of capacity. He looks at the grey expanse and lowers his
hands. \emph{``Hvíldu í friði,''} he says. \emph{Rest in peace.} Modern
Icelandic, a card-shop phrase, no traction, sliding off the reality of
the place like oil off water.

A ripple moves through the crowd. A woman clutches her shawl; a man
shifts his weight, ready to run. They wait for the water to answer, for
a wave to climb the cliff and erase the error. Rowan watches the bay. It
does not rise. It laps gently against the rocks, rhythmic and
indifferent. Hiss-thud. Hiss-thud. The sea does not strike. It accepts
the modern phrase with a terrifying politeness, the patience of a
creditor who knows that interest compounds. The lack of reaction is
worse than a storm; a storm would mean the conversation was still going
on. This means the sea has stopped negotiating and is content to wait
for the default.

Arnar stands a moment, staring at the calm water, understanding the
reprieve and its price. He turns his back on the ocean. ``It is done,''
he says, and the dismissal is instant. The villagers turn and walk away,
heavy and fast, no lingering, no shared memory, evacuating the platform
as if the air had gone toxic. Rowan remains, and watches Arnar climb,
and the harbour master stops when he reaches him. He does not offer a
hand. He does not say \emph{I am sorry} or \emph{thank you.} He studies
Rowan's face --- the set of the jaw, the stillness in the eyes ---
looking for the mark of the sentence Edda passed on. He finds it. He
nods, short and sharp, not approval, only the look a man gives when the
heavy tool has been picked up by the one other person in the shop.
\emph{You have it now.} Then he walks past, toward the harbour, to check
lines he knows will not hold forever, and Rowan stands alone on the
rock, feeling the deferred weight of the ocean pressing at the back of
his neck, the envelope stiff against his chest under two layers of wool.

\sceneline

The dream is not a narrative. It is a lesson. Rowan sits at a wooden
desk that smells of pencil shavings and chalk dust. The room is the
schoolhouse, but as it was meant to be --- warm, dry, the windows sealed
against draught, the charts on the walls intact and displaying the full
declensions for tide and wind and silence. Edda stands at the blackboard
with a piece of limestone chalk and writes a root word. \emph{Haf.}
Ocean. ``Decline,'' she says, not looking at him, looking at the word.

He opens his mouth expecting academic hesitation, the hunt for the
suffix, and instead his tongue moves with a muscular memory he does not
possess. \emph{``Haf\ldots{} Hafið\ldots{} Hafinu\ldots{} Hafsins---''}
He moves through the standard four cases at the speed of breathing.
``The binding,'' she commands, tapping the board so that the chalk snaps
and leaves a white dot. \emph{``Haf-und\ldots{} Haf-ing\ldots{}
Haf-hrafn.''} He speaks the ritual cases and the sounds are not foreign;
they fit his throat as if they had always been there. He understands
that the Aflending lives in the diaphragm, the Sæfangi in the back of
the palate, the Hrimfall in the teeth. He is not learning a language. He
is learning where to place a sound so that it lands on a target that has
no mass. Edda nods and erases the board. ``The lesson is concluded,''
she says. ``The recess is over.''

He wakes. There is no confusion, no disorientation. He opens his eyes to
the grey light of the Annex, the air cold and motionless.
\emph{``Frostið bindur gluggann,''} he says, and the sentence leaves his
lips before thought. \emph{The frost binds the window.} He lies still,
listening to his own voice --- the pronunciation perfect, the vowel
length exact --- and feels no thrill, only a dull acceptance, the way a
man accepts that he has walked with a limp his whole life and that the
bone has finally, quietly, set.

He sits up and looks at the window. It is only fogged, unremarkable. But
on the sill outside, and along the top of the boundary stone he can see
from the bed, and creeping up the whitewash of the Annex's northern
wall, salt has bloomed overnight --- not the ordinary crust of spray but
a fine ordered growth of crystal, looping and interlocking into a
lattice of sharp angles and closed curves. It is the Sæfangi ending. The
Sea-Holding case, written in salt across the skin of his own house. The
sea acknowledging the new signatory, laying its counter-signature on the
stone.

He dresses quickly. The buttons feel different under his fingers; he is
aware of the verb of buttoning, the noun of the cloth, the world
resolving into grammatical components. He steps outside. The gravel is
wet, the fog rolling in from the bay and erasing the horizon, and he
looks toward the spot where Brynja parked her jeep. The space is empty.
Tire tracks score the mud, deep and frantic, spinning and veering,
leading away toward the coastal road. He follows them with his eyes.
Fifty meters up, they reach the point where the road was edited out, and
they do not stop, and they do not turn. They fade into heather and moss.
Mud becomes grass; impression becomes texture. Brynja did not crash. She
drove out of the sentence and the narrative let her go.

He understands, dimly, why it could let her go. She had spoken Reykjavík
standard the entire time she was here --- clean nominatives, active
verbs, no ritual cases. The village never declined her. To its grammar
she was a foreign substance, never grafted in, only suspended in the
text like an unincorporated quotation, and a quotation can be lifted out
without tearing the page. She was never bound, and so the cases that
would have held a villager in place had nothing to grip. It is not a
consolation. It is the opposite. Because he knows he is no longer in her
category. Somewhere in these weeks his own speech began to take the
local endings, and now she is the third --- Edda, Sigrún, and the
granddaughter with the yellow jacket, three women gone from the sentence
by three different doors --- and he is the one still standing in it, and
the thought arrives cold and exact: the village keeps whoever it has
managed to bind, and it has bound him, and being kept is not the same as
being spared.

He turns back to the village. The fog thickens, swirling around the
eaves, and he feels the pressure building behind the mist, an
accumulated memory pressing on the thin membrane of the land. He walks
to the boundary stone Arnar planted and places his hand on the cold
basalt. There is a vibration in the rock, a low hum, a question waiting
for an answer. Rowan stands at the edge of the text, no longer a reader.
He inhales, holds the air in the center of his chest, and waits for the
first line of the storm to be spoken.

\chapter{}

Rowan packs the Annex with the detached precision of an archivist
closing a wing of a library that no longer receives visitors. He does
not rush. To rush would be to admit a fear that requires a verb of
motion, and he is holding himself, this morning, in the nominative. He
places the recorder --- its screen dead and grey --- into the padded
compartment of his bag. He stacks the notebooks, the early ones full of
clean English syntax, the later ones scarred with the heavy angular
script of Vikamál, and slides them into the main pocket. He checks the
room: the bed stripped, the desk clear, the space anonymous again, a
vacuum waiting for the next noun to occupy it. He shoulders the bag. The
weight is significant, a physical reminder of the data he tried to
extract and could not translate. Against his chest, under the wool, the
envelope. He has not opened it. He is not finished enough to read it. He
believes her.

He steps out into the morning. The fog hangs ten meters from the door, a
white wall that swallows the world's context. He walks north toward the
coastal road that once connected Brynjavík to the arteries of the
continent. His gait is steady. His boots find purchase on gravel that
feels disturbingly loose, as if the stones were merely resting on the
surface rather than set into the earth. He passes the marker stone Arnar
planted, the rune for boundary wet with condensation, and steps past it.

He expects resistance --- wind rising, mud deepening, a wave thrown
across the track. There is none. There is only absence. Ten meters past
the stone, the gravel road stops. It does not end at a cliff edge; it
does not wash out into scree. It simply thins into the texture of the
moor. One meter: road, graded and defined. The next: the vague general
idea of ground. He stops, the toe of his left boot on gravel, the toe of
his right hovering over moss that looks poorly finished, its colour a
flat unshaded green. He steps onto the moss. It holds his weight, but it
feels sponge-like, insufficient.

He walks forward, searching for the bend in the road that skirts the
headland. He remembers it. He has driven it. He has seen it on the map.
He looks up to find the headland, and it is there, a dark mass against
the white sky, but it is wrong --- the profile smooth, lacking the
jagged fractures he documented on the day he arrived, like a shape
sketched in and not yet given its detail. He blinks, and when he opens
his eyes the headland has shifted, its peak now to the left of center.
He turns to check the coastline, and the sea is a flat grey plane, no
waves, no texture, a field of colour with no incident in it. He looks
back at the headland. It has settled into a low ridge.

His heart hammers, but his mind stays coldly analytical. The land is
relying on his attention to keep its shape, he notes; where no one
looks, the geometry lets go. He takes another step and the ground
ripples --- not an earthquake, an uncertainty --- moss becoming slate
and then moss again, cycling through the terrains it might be because no
one has committed the surface to a final form. He scans for the tire
tracks of Brynja's jeep and finds a rut that ends abruptly, the earth
smoothing over as if the vehicle had simply stopped displacing matter.
She did not crash. She drove until the world stopped describing the
road, and she followed the silence out.

He turns in a circle. The landscape is fluid, sickeningly pliable. A
boulder in the middle distance elongates, stretching upward like drawn
taffy, then snaps back to a sphere when his gaze settles on it. The
horizon bends inward, unable to hold the illusion of infinite distance.
This is a place that was held together by being spoken, and out here no
one is speaking it, and so it is coming untied, the way a knot comes
untied when the tension leaves the line.

He looks back. Brynjavík stands in the fog --- small, huddled,
monochrome --- but it is sharp. The rooflines are distinct. The wall
stones are counted and fixed. Smoke rises from Arnar's chimney and obeys
the wind, dispersing in a plume that behaves. The village is rendered in
hard edges because it is still being spoken, sustained by the grammar of
the few throats left in it. Rowan stands at the edge of the unwritten
world, looking at the blur of moor where a sheep appears and dissolves
into white, and then at the village, where the walls are hard and the
consequences are fatal. He does not feel relief. He feels the cold
administrative weight of a decision with no good option. To stay is to
be held inside a dying sentence. To leave is to become a smudge in a
deleted draft, like the granddaughter, unmourned because unnameable.

He adjusts the strap of his bag. He turns his back on the fluid,
uncommitted horizon and walks back toward the gravel, toward the
boundary stone, and when his boot hits the road the sound is crisp --- a
clean, defined crunch, a grammatical sound. Rowan Hale returns to
Brynjavík not because he is brave and not because he is resigned. He
returns because it is the only place left where he still has a name, and
because his mother left the debt unpaid, and because he has finally
understood that the envelope against his chest is not a letter to a
stranger. It is a forwarding address.

\sceneline

The church is on his way, and he does not intend to stop, and he stops.

It is a small white timber building with a red roof, of a design
repeated four hundred times around the coast of Iceland, and it is warm
inside, and it smells of paraffin and old wood and the wet wool of a
congregation that left an hour ago. Séra Hjálmar is at the front,
gathering hymnals off the pews and stacking them in the crook of his
arm, and he does not look up when the door opens, because in nineteen
years no one has come through that door whom he needed to look up for.

``Dr.~Hale.''

``I wanted to ask you something before I go up.''

``You are not going anywhere, though.'' The pastor sets the stack on the
front pew and straightens, and his face is kind and entirely unsurprised.
``Forgive me. Arnar came by at seven. He said you had walked out to the
road and come back. He did not say it as gossip. He said it the way a
man tells you the tide is in.''

Rowan takes the hymnal off the top of the stack and opens it to page
nine and holds it out, and does not say anything, because the question
is on the page and they both know what it is.

Hjálmar looks at it for a moment. Then he sits down on the pew, heavily,
like a man who has been waiting a long time to be asked and has not
enjoyed the wait.

``\emph{Verði að þínum vilja,}'' he says.

``Every book in the village. Including the ones bought in Reykjavík.''

``Including the one on my own shelf that my mother gave me at my
ordination, in Selfoss, in 2001, which read \emph{verði þinn vilji} when
she wrote in the front of it, and which reads what it reads now.'' He
says this without drama, as a fact he checked years ago and has had to
live beside since. ``I did not do it. My predecessor did not do it. It
is not done. It is what happens to the sentence when it is kept here.''

``Why that line?'' Rowan sits down across the aisle. ``There are twenty
lines in the prayer. Why the will?''

``Because it is the only one where something other than a man is
permitted to want something.'' Hjálmar folds his hands. ``Think about
what you have been doing all autumn, Doctor. You have watched this
village spend a thousand years teaching itself one rule, and only one:
that a will in the nominative can act. That is the whole of it. That is
the entire theology of this coast, and it is not a theology, it is a
load rating.''

And it arrives, all at once, the way the true reading of a difficult
line arrives --- not as a discovery but as a thing that has been legible
for weeks and merely unread.

The rule cannot make exceptions. That is what makes it a rule. If the
congregation of Brynjavík stands up on a Sunday morning, eighty voices,
half a kilometre from the water, and says aloud that a will --- any
will, even and especially the highest one --- may stand in the
nominative and be done, then they have not made a statement about God.
They have made a demonstration. They have shown the listener, in chorus,
weekly, for four centuries, that the case is available. That it works.
That a thing which wants may be a subject.

And there is something in the bay that has been listening to every
service since 1732, and it does not know theology, and it does not care
whose will it is. It only learns the grammar.

``You demoted Him,'' Rowan says. ``To keep the rule airtight.''

``We did not demote anyone.'' Hjálmar's voice does not rise, but
something in it hardens, and Rowan understands that he has just heard
the pastor's one point of pride. ``That is what a man from Cambridge
would call it and I would ask him to consider whether he has understood
what he is looking at. We did not make God smaller. We stopped shouting
His name across the water.'' He turns the hymnal on his knee. ``There is
no difference, Doctor, between \emph{thy will be done} and \emph{let it
be done according to thy will,} except one. In the second, the will is
still perfect and still sovereign and still the thing the whole world
bends toward. It simply is not standing up where the sea can see it.''

He hands the book back.

``I say it that way,'' he says. ``I have said it that way for nineteen
years, in front of a font, in a cassock, in a church of the Evangelical
Lutheran Church of Iceland, and I have never once felt that I was lying,
and I have never once written to the bishop about it, and I am not going
to. He would send someone. Someone would correct the books.'' A long
breath. ``And then, one Sunday, eighty people would stand up in this
room and say the other thing.''

Rowan looks at the altar, at the plain wooden cross, at the numbers
still up on the hymn board from the morning.

``Does it trouble you?'' he asks. ``Not knowing which one holds this
place up?''

``It troubled me for eleven years,'' Hjálmar says. ``Then I buried a
child, and I said the burial over her, and I meant every word of it, and
afterward I walked up the hill and stood at the back while Edda said the
other thing over her, and I meant that too.'' He stands, and picks up
his stack of books. ``One of them tells her mother why. The other one
keeps the ground under the churchyard. I have stopped requiring them to
compete. I would recommend it. It is late in the season to be sorting
your convictions into the correct boxes.''

At the door he says, without turning: ``There will be a service on
Sunday. There is always a service on Sunday. You do not have to come and
you would be welcome.'' And then, quieter, in the tone of a man setting
down something he has carried: ``Hold the line well, Dr.~Hale. I will
hold mine. Between us, perhaps, the two of them will hold each other.''

\sceneline

He does not return to the Annex. He walks past the darkened windows of
the guest house, his stride measuring the remaining pavement, and goes
directly to the shoreline, to the basalt shelf that juts into the bay
like a tongue of hardened lava. The wind has died. The silence is not
empty; it is expectant, pressing on the eardrums, waiting for an input.
He climbs onto the platform and stands where Edda stood, where Arnar
failed. He does not set down his bag. He does not clear his throat or
close his eyes to center himself. He looks at the water.

The bay is high, swollen with the grey mass of the Atlantic, lapping
inches from his boots, confident, having already claimed the margin.
Rowan opens his mouth. He speaks.

\emph{``Hér stendur maður.''} \emph{Here stands a man.} The nominative.
He places himself as subject. He does not shout; the volume is
conversational, the tone of a surveyor reading a coordinate aloud. He
shifts to the accusative and addresses the boundary. \emph{``Ég bind
steininn við landið.''} \emph{I bind the stone to the land.} His voice
is different now, not the borrowed accent of a Cambridge academic
imitating a dialect but something produced deep in the chest, resonant
and flat, the consonants striking the air like flint. He moves to the
dative, the giving. \emph{``Ég gef þér salt, en ekki blóð.''} \emph{I
give you salt, but not blood.} The water at his feet ripples --- not a
wave, a shudder, as if its surface tension had changed. He feels the
grammar lock into place, not magic but engineering, a sea-wall built out
of syntax. He reaches for the genitive. \emph{``Hafsins er dýpið, en
landsins er brúnin.''} \emph{The deep belongs to the ocean, but the edge
belongs to the land.} Separation. Ownership defined.

Then the forbidden cases. He does not hesitate. He steps into the
structures Edda died teaching. The Aflending, the unbinding:
\emph{``Aldan skal unda. Straumur skal undi.''} Forward momentum
cancelled, the suffix applied, and the water flinches, the bay pulling
back and sliding smoothly off the basalt, no crash, no foam, a retreat
like a heavy curtain drawn across a stage. The Sæfangi, the sea-holding:
he raises his hand and traces the shape in the air --- the hook, the
loop, the slash. \emph{``Nafnið er fast. Staðurinn er hér.''} \emph{The
name is fixed. The place is here.}

As he speaks the locking clause a shadow moves at the edge of his
vision. Arnar steps onto the rock, quiet, his boots silent on the wet
stone. He does not look at Rowan with awe. He looks with grim
recognition, the expression of a man watching another soldier pick up a
live grenade. He stands two paces behind, to the left, and does not take
the lead. He waits for the gap in the syntax. Rowan initiates the final
case, the Hrimfall, Winterfall, the contract sealed. \emph{``Ísinn
kemur---''} he begins.

\emph{``---og þögnin geymir,''} Arnar finishes. \emph{The ice
comes\ldots{} and the silence keeps.} Arnar's voice is a low rumble,
ballast under Rowan's line, reinforcing the structure. Where Rowan
defines a wall, Arnar mortars the joints. Together they speak the new
coordinates. They do not ask the sea to leave; they concede the outer
fields, abandon the western ridge, surrender the memory of sheep pens
and old drying racks. The circle tightens around the village's core.
``Hér er markið,'' Rowan says. \emph{Here is the mark.} ``Og ekki
lengra,'' Arnar adds. \emph{And no further.}

The sea responds with a massive hydraulic concession. The water level in
the harbour drops two meters, exposing the black slime-slicked wood of
eighteenth-century pilings, revealing the stone foundations of
boathouses lost in the storms of the 1920s. The revealed structures are
not ruins; they are entries. The sea has kept them catalogued in the
dark, and now it lays them bare on the mudflats, wet and black and
smelling of ancient decay, as proof that it remembers everything and
forgives nothing without being asked in the correct case. It is not a
victory. The sea has not fled. It has accepted the edit and agreed to
the margin, for now.

Rowan stops speaking. The silence that follows is instant and concrete.
The vibration in the basalt ceases. The air feels thinner, sharper,
cleared of the static that has hummed in his ears for days. Arnar
exhales and takes the tobacco pouch from his pocket, and his hands
shake, but he rolls a cigarette with practiced muscle memory. He does
not light it. He holds it, staring at the exposed ribs of the old pier.
``The western wall will need pointing,'' he says, his voice raspy.
``Since we are keeping it.''

``The geometry holds,'' Rowan replies. He is drained, heavy as wet rope.

``It holds,'' Arnar agrees. He looks at the village, huddled and small
beneath a vast indifferent sky. ``It is a small room. But it has a
floor.'' They stand together on the wet rock, two men holding a heavy
door shut against a gale, neither celebrating the survival nor mourning
the loss of the view.

And standing there, spent, Rowan lets himself think the thought he has
been holding at arm's length since the cottage, because he owes it to
the three of them not to look away. He is on this rock because a woman
poured out her voice and died, and because another woman before her
poured out her nerve in a schoolhouse and was folded into a cairn under
her own withheld name, and because a third drove out of the world
unnamed carrying a card full of images that curdled into script. He
inherits the stewardship not despite the fact that they were used up but
because they were --- the net is woven from the ones who were spent
thin. He is the man left standing, and the village will call it
strength, and it is not strength; it is only survival, and survival here
has always been paid for by someone else, and always, so far, by the
women. He does not know how to carry that. He decides only that he will
carry it awake, that he will not let it become one more thing held at a
distance and re-labelled, that whatever else the boundary costs, it will
not cost him the honesty of knowing who paid for it before him.

\sceneline

Night does not bring darkness so much as a suspension of light. Rowan
sleeps in the Annex, his body motionless, his breath shallow and
synchronized with the draught whistling through the keyhole, and he is
not in the Annex. He is in the schoolhouse. The dream is high-contrast,
etched in the stark black and white of a woodcut. Sun streams through
clean windows, the glass free of grime and condensation, the dust motes
in the beams not debris but pollen, potential life. Edda stands at the
blackboard --- not the dying woman, nor the corpse on the platform, but
the teacher, her spine straight, her hand steady as the chalk drives
across the slate. She writes the endings. \emph{-und. -ing. -hrafn.} The
classroom is full. The children of Brynjavík sit in rows --- not the
children who left for Reykjavík, but the children who are always here,
the potential, the memory, the future tense, in wool sweaters, with
attentive eyes.

``Again,'' Edda says. \emph{``Sæfangið bindur,''} the children chant,
their voices a single harmonic chord. \emph{The Sea-Holding binds.}
Rowan stands in the doorway --- not seated at a desk, no notebook ---
leaning against the frame, arms crossed, watching the lesson. He
understands the hierarchy. He is not here to learn conjugation; he is
here to keep the classroom standing so the lesson can go on. Edda looks
up from the board and sees him and does not smile. She nods, a single
sharp depression of the chin. The handover of the chalk. ``The margin is
set,'' she says to the class, but her eyes are on him.

Rowan wakes. The transition is gradual, a rising from the dream's depth
to the surface of a grey and cold morning. He lies still, feeling the
dream recede, but the weight of the chalk remains in his hand. He flexes
his fingers. The stiffness is gone. The tremors that plagued him, the
somatic symptoms of a body being asked to hold a grammar it was not
built for, have vanished. He looks at the window. The pane is frosted,
and it is not random fernwork; the ice has crystallized into a rigorous
geometric structure, a perfect lattice of intersecting lines and hooks
and closures. The Sæfangi case. The ending he spoke on the platform.

He throws off the quilt. The cold bites, but it feels bracing rather
than hostile. He walks to the window on floorboards that are solid under
his feet, the house settled, its joints quiet, its foundation gripping
the earth. He raises his hand and presses his index finger to the glass
and traces the frost, following the binding rune, warm skin sliding over
cold ice, and the sea has answered --- its counter-signature written on
the glass. Accepted. The ice melts beneath his touch and a drop of water
runs down the pane, cutting through the pattern, erasing the sign, and
it does not matter, because the sign does not need to persist. The
agreement has been made.

He steps back. He does not say \emph{I am the guardian.} He does not
think \emph{I have saved them.} He adjusts his stance --- feet widening
slightly, center of gravity lowering, shoulders rolling back, settling
the village's weight into the architecture of his spine, heavy but
balanced now, because he knows at last how to carry it. He looks out at
the village as the fog lifts, revealing black clouds stacked over the
ridge and the narrow safe streets between the remaining houses. The
world has snapped into a smaller, harder shape. There is less of it, but
what remains is real. Rowan breathes in. The air is cold and tastes of
salt, and he holds it a moment, feeling the grammar of it in his lungs,
then exhales, ready to begin the long work of maintenance.

\chapter{}

The dawn dilutes the grey dark just enough to distinguish water from
rock. Rowan Hale stands on the basalt shelf, his feet finding the same
two depressions in the stone they found yesterday and the day before,
the leather of his boots long since molded to the geology. He is no
longer a visitor standing on the surface of the place. He is a fitted
part of the platform. He begins the morning maintenance.

He speaks the nominative of position, and the sound that leaves his
throat is not smooth; it is textured, roughened by salt air and
repetition, ground by thousands of utterances into something harder and
more fibrous than the voice he arrived with. \emph{``Hér er
steinninn.''} \emph{Here is the stone.} The water at the shelf's edge,
which has been creeping upward with the stealth of rising damp, halts
and accepts the definition. He shifts his weight, and his left knee
clicks, a dry popping sound that marks the slow erosion of the joint,
and he does not wince. The pain is a kind of information. It confirms
that the cold has settled deep enough to be respected.

He moves to the genitive of boundary. \emph{``Vatnsins er djúpið.''}
\emph{The deep belongs to the water.} He extends his right hand, the
skin of the fingertips thickened, the ridges of the prints worn nearly
smooth by years of stone and wood, and traces a horizontal line in the
air, a cut severing the rising tide from the vertical cliff. The sea
obeys, sliding back to expose a strip of black algae that glistens like
wet oil, the retreat precise, exactly to the line his voice prescribed.
There is no negotiation this morning, only the execution of a standing
agreement. The ocean is not an enemy to be fought. It is a great volume
of water that must be spoken to, in the right case, at the right hour,
forever.

He has stood eleven Great Keepings. He does not think about them the way
he thought about the first, that midwinter, when he walked the whole
perimeter in the dark with sixty-odd people and four hurricane lamps and
braced himself all night for the reading of the names to be the terrible
thing a season in Brynjavík had taught him to expect. It was not
terrible. It took four hours. A woman beside him complained quietly and
without interruption about her hip, and a boy was sent home at the
halfway stone for laughing, and someone had brought a flask of coffee
that had gone cold by the second hour and was drunk anyway.

He understood something that night that he has never been able to put
into a sentence a colleague would have accepted, and has therefore never
written down, which is now his standard test for whether a thing is
true. The names are not said to the water. They are said in front of it.
There is a difference, and the difference is the entire practice, and it
is the reason a village can stand on a cliff in the dark and speak
aloud, for four hours, every name it has, and be safe --- and the reason
that a single man, saying a single name, in the wrong case, at the wrong
hour, is not.

Behind him the village waits. It is smaller than it was in the time of
the buses, but the contraction is not violent; there are no heaps of
rubble where the western farmstead stood, only moor, green and spongy,
quietly reclaiming the ground where a kitchen once was. When a house
loses its speaker, it loses its right to displace the air; the fog moves
in, the definition blurs, and one morning the grass is simply
uninterrupted. Rowan turns from the water, satisfied with the margin,
and walks back along the path toward the cluster of remaining
structures. He passes Bjarni, seated on a plastic crate mending a net
--- Bjarni, who was hauled off the harbour wall by a rope because no one
could find the case that would have let them use his name, and who came
back up the steps with a vowel in his throat he has never quite put
down. He has a daughter now --- small, grave, better at the boundary
walk than her father was at twice her age --- and he does not talk.
Whatever the water taught him on the way under, he has kept it the way
this place keeps everything, which is to say completely and without
reference. The two men nod once
to each other, downward, a shift change acknowledged, the perimeter
held. Further along, Einar sits in the lee of the curing house with the
glass eye turned to the weather and the good one on his hands, and lifts
two fingers off the twine as Rowan passes, and that is the whole of the
morning's conversation, and it is sufficient.

He reaches the cottage that was Edda's and is now the mechanism that is
Rowan. The hinges are greased and silent. Inside, the air is dry, kept
so by the stove he feeds with regulated turf, and he hangs his coat on a
peg, the wool stiff with salt. In the corner his academic notebooks are
stacked, the pages yellowing and curling --- theories of linguistic
drift, diagrams of isolation, publication dates that have stopped
meaning anything, the drawings of a machine made by someone who never
once had to run one. He ignores them and picks up a chisel.

The roof's main beam, a massive length of driftwood pine, is covered in
carvings. Near the ceiling are Edda's marks, sharp angular runes that
handled the storms of the nineties; below them are Rowan's own
additions. He runs his thumb over a fresh gouge. Yesterday the wind
tested the north gable, seeking a weakness in the accusative, and he
felt the draught and knew the verb was too weak to hold it. He sets the
chisel and carves the correction into the infrastructure of the house.
\emph{``Styrkja.''} \emph{Strengthen.} He taps the handle --- \emph{tak}
--- and the wood curls away, and resin fills the room, a necessary
scent, a binding agent, and he cuts the suffix \emph{-ur,} locking the
strength into the noun of the house. His hand cramps, the muscles
seizing into the shape of the tool, and he waits, breathing steadily,
until the spasm passes. The body wears down. The tool wears down. The
grammar holds. He surveys the lexicon on the walls --- not a history but
the wiring of a house that will fall if he stops carving, dissolve if he
stops speaking. He is not preserving a culture. He is keeping the rain
out. He sets the chisel down and pours water from the kettle, and his
reflection in the dark window shows a face eroded by the same weather
that strips the cliffs, the hair gone grey, the eyes flat, become
instruments that no longer look for meaning, only for fault lines. He
drinks. The heat soothes the throat that will speak again at dusk.

\sceneline

The schoolhouse does not hold heat; the radiators are cold iron ribs,
dead since the oil route was struck from the district map. Rowan comes
here because the last fragments of the curriculum remain, and because
Sigrún is here. She sits at a child's desk by the window, wrapped in a
blanket that smells of mothballs and damp wool, her skin gone the yellow
of the parchment maps on the wall. She does not cough. The decay is
silent, a slow unbinding at the level of the cell.

He sits opposite and sets a tin of dried fish on the desk. She ignores
it. ``The wind has changed,'' she says. ``North-northwest. The
declension must shift.''

``I've been using the dative. Giving the wind to the hill.''

She shakes her head. ``Too heavy. Give it all to the hill and the hill
slides. Split the load.'' She lifts a hand, the fingers crooked, the
joints swollen, and breathes on the cold slate of the old blackboard
beside her until a patch of fog forms on the stone. Into the fog she
draws a hook that loops back on itself, the knot used to secure a boat
in a gale. ``The Sæfangi,'' she murmurs. ``Only the first clause. Don't
draw the second.''

``If I don't close it, the loop stays open.''

``Yes.'' Her eyes are sharp, lucid within the ruin of her. ``Leave it
open. Let Einar speak the second clause. Or let the silence speak it.''
She withdraws her hand, and the breath-fog holds the hook a moment on
the slate before it fades. ``Edda tried to hold the whole bay,'' she
says. ``She tried to be the wall. Walls crack. You are not a wall,
Rowan. You are a net. Nets have holes, and the holes are how they
hold.''

``Distributed grammar,'' Rowan says, the academic term surfacing like a
phantom limb.

``Spread the load,'' she corrects, in the plain English she has
scavenged from some old manual. ``Carry everything yourself and you
fail. Carry it together and the failure of any one throat is only a
hole, and a net survives holes.'' She closes her eyes. ``Go now. The
lesson is done. I need to save my breath.'' He leaves her in the grey
light, her finger resting near the fading hook.

The next morning the chair is empty. There is no body, no sign of
struggle, the blanket folded with its edges squared to the grain of the
desk. On the blackboard, written steadily in white chalk: \emph{Gengin.}
Gone, or Walked. The damp air has already erased the lower halves of the
letters.

The village convenes at the edge of the schoolyard --- not to mourn but
to re-sort the register. They stand in a triangle, Rowan at the apex,
Einar to the left, old Halla to the right, and the sea before them is
flat, a perfect mirror, waiting. Rowan speaks first. \emph{``Sigrún var
hér.''} \emph{Sigrún was here.} He steps back. Einar steps forward.
\emph{``Við kveðjum hana.''} \emph{We release her.} He retreats, and
Halla steps up. \emph{``Til moldarinnar, ekki til hafsins.''} \emph{To
the earth, not the sea.} She leaves the final vowel unresolved, and the
silence carries the rest of the sentence, the grammar thinned
deliberately across three throats so that there is no single point that
can fail. The sea scans the phrasing, finds no amplitude to challenge,
and stays dormant. The water laps against the rock. Hush. Hush. The
books balance --- this once, and Rowan lets the old word come, because
Sigrún would have used it, because she died correcting his arithmetic,
because a debt that closes cleanly deserves, for one morning, to be
called by the name she would have given it.

\sceneline

The perimeter tightens again. The outlying sheds lose their standing in
the sentence and the moor absorbs them, and the distance from one end of
the village to the other shortens by subtraction, month by month, until
a man can cross the whole of Brynjavík in the time it once took to walk
to the quay.

Arnar lasts the winter. The friction of it wears him down; his movements
go segmented, each one costing him a visible sum, and when the spring
tide rises he does not come back from the harbour. Rowan finds him at
low water, curled like dried kelp between the pilings, one hand extended
--- not clawing, offering. ``The bollards,'' Arnar whispers. ``The third
one. Loose.''

``I'll fix it.''

``No.~Speak to it. The iron is tired. Use \emph{festar.}'' Metal
fatigue, given its true case. A wheeze, and something that on another
face might have been humor. ``And the moon --- lie to it. Tell it the
boats are stone.'' His hand drops, and he becomes an object, and the sea
does not gloat, and does not hurry. Rowan closes the pale eyes and lays
his own hand on the loose bollard. \emph{``Þú ert berg,''} he says.
\emph{You are rock.} The metal vibrates once and settles, and the
maintenance continues, because the maintenance is the only thing that
does not end.

Summer brings a few visitors, blown off course, arriving by accident
because Brynjavík is no longer a place a person can mean to go. A
doctoral candidate finds her way down the ruined track one bright
afternoon, a recorder blinking red in her hand, her gaze hungry in a way
Rowan recognizes with a sorrow he did not expect. ``A linguistic time
capsule,'' she says, delighted, looking at the carved beams, the
truncated speech.

``It's a dam,'' Rowan says.

She does not understand him, and he does not try again. He gives her
sheep words and fish words, the safe currency, the vocabulary a language
leaves out on the counter for strangers.

She is packing the recorder into its case when she says, conversationally,
that she should thank the woman who pointed her this way in the first
place --- a producer at the council, media side, very organised, drove
her out to the junction herself and drew the last part of the route on
the back of a receipt because the satnav would not take it.

``Þórhallsdóttir,'' the candidate says, hunting the name. ``Brynja
something. Do you know her? She said she'd been up here once, years
ago.''

Rowan holds very still. ``Where is she now?''

``Hvalsker.'' The girl zips the case. ``Small place, further round the
coast. She moved out there after the divorce, I think. Very pretty. They
have a lovely harbour --- well, they're rebuilding it again, apparently.
Third time. She was funny about it, actually, she said the whole village
does nothing but write grant applications for that wall.'' She laughs.
``Anyway. She was lovely. She said to say hello if I found the place.''

``Did she,'' Rowan says.

``She said --- '' and the girl frowns slightly, retrieving it with the
mild puzzlement of someone repeating a remark that had struck her as odd
and then stopped mattering. ``She said, \emph{tell them I did get out,
if there's anyone left who'd know what that means.}'' She shoulders the
bag. ``I assumed it was a joke about the roads.''

Rowan says nothing. He walks her to the car, and shakes her hand, and
watches her go.

She is alive. She is well. She lives in a bright house in Hvalsker,
which is the harbour that put the grammar down two generations ago
because a sensible people decided it was a superstition, and where the
sea has taken the wall three times in forty years, and where they are
still applying, and where they will go on applying until the
applications outlive the applicants. She got out. She drove through the
part of the world that had stopped describing the road and came out the
other side into a country that had never heard of any of this, and she
is fine, and she has a job and a divorce and a view, and once a year or
so she says a thing at a party that makes people look at her strangely,
and she has learned to say it as a joke about the roads.

She escaped into the exact future Brynjavík has spent a thousand years
refusing.

And she is right. That is the part Rowan makes himself hold. She was
right to go, and she is happier than anyone here, and the price of her
being right is only that the wall at Hvalsker will go again in the
spring, and that nobody there will connect the two facts, and that
nobody there will have to.

He watches the dust of the candidate's going settle back onto the track.
The disappointment he sent her away with is a mercy. His relief is the
relief of a man who has kept someone out of a machine that would have
taken her fingers. And under it, unwelcome and exact, is the thing he
has spent years not looking at: that he could have gone too, and that
the difference between him and Brynja Þórhallsdóttir is not courage and
is not virtue. It is that she was never bound and he was, and being kept
is not the same as being spared.

The light shifts. The grey returns, as it always returns. Rowan stands
on the basalt, his hair white now, his hands stained with the grain of
the wood he cuts. He is not the linguist who arrived. He is the operator
of a thing the linguist could only describe. He speaks into the dusk.
\emph{``Hér erum við.''} \emph{Here we are.} Plural --- because he is
not alone now, because Einar holds the second clause and Halla the
third, because Sigrún taught him that a net is stronger than a wall
precisely where it is full of gaps. \emph{``Sjórinn bíður,''} he says.
\emph{The sea waits.} \emph{``Landið stendur.''} \emph{The land stands.}
He traces the Sæfangi in the air, the motion mechanical, functional,
worn to a habit, and the wave hesitates, reads the syntax, and recedes.
Rowan stays, the sentence left open, the verb still active, the loop
unclosed and held by other throats. The boundary holds tonight. Because
it holds tonight, and because the grammar is distributed now across more
than one failing body, it will hold tomorrow.

\clearpage
\thispagestyle{fancy}
\vspace*{2.2\baselineskip}
\begin{center}
{\fontsize{26}{30}\selectfont\codastave}\\[0.7em]
{\fontsize{11}{13}\selectfont\ebgaramond ❦}
\end{center}
\vspace{2.2\baselineskip}
\noindent

The light that bleeds into the sky above the North Atlantic is a
negotiation. It bargains with the dark water for the right to lay itself
across the wet black stone of the coast, and on the basalt shelf that
juts from the broken jaw of the land, Rowan Hale stands as the mediator.
He is not the man who arrived with a recorder and a waterproof notebook;
that man has been edited out of the narrative, his softness abraded by
salt, his academic curiosity replaced by the heavy rhythmic necessity of
maintenance. Rowan is thinner now. The cables of his neck stand against
the collar of a sweater that smells of lanolin and woodsmoke, and his
hands, hanging at his sides, are maps of the terrain --- the knuckles
swollen with cold, the skin thickened into a hide that resists the
stinging spray.

He breathes. The air is frigid and tastes of iodine and ancient decay,
and his exhalation is a white plume that drifts a moment before the wind
takes it and disperses the heat of his body into the vast indifferent
cold. Behind him the village huddles, smaller than it was a year ago,
its geography revised --- the outer farmsteads, the rotting fences of
the sheep runs, the connector road that once wound toward the pass, all
of them not ruined but simply absent, the land healed over them, smooth
and green and amnesiac. The village has contracted to its grammatical
core, a few cottages that still hold warmth, the smoke rising from their
chimneys in vertical lines of defiance, the silence between the walls
doing the work of mortar. Rowan does not look back at what has been
lost. To acknowledge an absence is to invite the void to expand. He
keeps his eyes on the water.

The ocean is restless. It heaves against the stone shelf, a grey
muscular mass that does not answer to the tide tables but to the logic
of memory, slapping the rock, hissing, testing the perimeter, seeking a
lapse in the syntax, a stutter in the sentence that decides where the
land ends and the water begins. Rowan plants his feet in the two smooth
hollows worn by the boots of Edda, and her mother, and the speakers
before them, and his feet fit the mold, and the stone grips him. He
raises his hands, the movement stiff, the cold long since settled into
the joints as a permanent ache that serves to remind him of the stakes:
if he stops, the pain stops, because the body stops.

He waits for the trough of a wave, because the word must land in the
silence between crashes or the meaning will be drowned before it can
set. The water draws back, sucking at the shingle with the sound of a
thousand small stones inhaled, and Rowan speaks.

\emph{``Hafið man.''} \emph{The sea remembers.} The nominative, the
naming of the subject. He does not shout; shouting suggests panic, and
panic is a weakness the grammar cannot support. He projects from the
diaphragm, resonating the vowels in his chest so that they leave him as
hard distinct objects. He names the water not as a body but as a
consciousness, granting it subject status, because to hold a thing you
must first admit that it exists and that it has a will. The Atlantic
answers --- not with a surge but with a subtle, terrifying reordering of
its surface. Ten meters out, the foam on the crest of a dark swell does
not scatter at random; the white arrests its own decay and drifts
together against the wind, forming lines and curves on the black canvas
of the trough, spelling back the echo of his claim, transcribing the
input. It acknowledges the subject.

He holds his stance while the wind tears at hair gone the colour of
surf, feeling the vibration of the acknowledgment travel up through the
soles of his boots, the platform humming, the entity in the water
accepting the call. He does not relax. The acknowledgment is only a
handshake, the connection made and the argument not yet won. The sea
remembers, yes; it remembers everything --- the valley before the lava,
the breath of the men it took last winter, the taste of the wood of the
pier it swallowed. Now he must tell it what to let lie.

\sceneline

He shifts the grammar, moving from subject to object, from naming to
command. \emph{``Landið gleymir hafinu.''} \emph{The land forgets the
sea.} He pushes the direct object away, placing the water in the
position of the thing forgotten, and the logic is brutal in its
simplicity: if the land forgets the sea, the sea cannot touch the land,
because memory is the bridge and forgetting is the chasm. The effect is
hydraulic. The water at the base of the platform does not merely flow
back; it recoils, shoved by the invisible hand of the case, the tide
withdrawing three meters in a single smooth motion, sliding off the rock
as if the stone had turned white-hot. In the newly exposed seabed,
stones that have not seen air since the last high tide shine like wet
eyes, and a tangle of black kelp lies steaming in the cold dawn.

He does not pause; the vacuum must be filled with the next case before
the water understands the trade. He steps forward, his toes over the
lip, and offers the tribute. \emph{``Hafinu skilum við orðum.''}
\emph{To the sea we return words.} The dative, the giving. He feeds the
bay's hunger with syllables so that it does not feed on silt and
foundation stone, and the water ripples in perfect concentric rings that
spread from the point beneath his face, flattening the incoming swell,
and he feels the drain of it --- the dative takes, it takes warmth from
the blood and moisture from the throat, and his skin tightens as if the
water in his own cells were being offered up along with the nouns. He
persists, reaching for the genitive, the root of ownership.
\emph{``Minninga er málfræðin sköpuð.''} \emph{Of memory, all grammar is
made.} His voice loses its human rasp and takes on the timbre of the
basalt itself, grinding and deep, asserting that the structure of the
world belongs to memory and therefore to those who remember. The shadows
on the rock elongate and snap back into hard pools of black; the light
itself seems to struggle to hold the world steady between the fluid
chaos of the ocean and the rigid order of the sentence, and cold sweat
runs down his spine and freezes at his waistband, because the genitive
is heavy, and it requires him to hold all of it in his mind at once ---
every drowned sailor, every lost ship, every word spoken in Brynjavík
for a thousand years --- to hold it and to claim it.

Then the precipice. The Aflending, the unbinding. This is the case that
broke Edda: the impossible one, that unbinds a force by naming its
freedom and, in the same breath, binds it to a limit. His hands tremble
as he raises them high, fingers splayed, cutting the sigil into the grey
air --- the hook, the slash. \emph{``Með aflendingu\ldots{}''} he
begins, and his voice shakes, the syllables like glass in his throat.
\emph{``\ldots und\ldots{} unda\ldots{} við höldum.''} \emph{By
unbinding, we hold.} He speaks the suffixes that cancel motion ---
\emph{-und,} a dull impact, a dead stop; \emph{-unda,} the echo of the
stop --- and the sea answers with a violence held in check. A swell
rises in the middle distance, a wall of water dark and smooth as
obsidian, higher than the platform, rolling forward in silence,
gathering mass, meaning to crush the syntax and wash the speaker off the
rock and edit him into the past tense. Rowan does not flinch. To flinch
is to break the case. He holds the suffix, grinding it through his
clenched jaw, imposing the limit, and the wave hangs --- cresting, never
breaking, suspended ten meters out, a monument to potential energy held
in place by the invisible wire of the grammar. It trembles with the
desire to fall. But the word holds it, because the Aflending has unbound
it from the urge to strike. Edda died because a doubt let the weight of
water into her lungs before the water ever touched her. Rowan has no
doubt left. He burned the last of it as fuel a long time ago.
\emph{Stay,} his whole body says. \emph{You are unbound from the shore.}
The wave groans, a deep tectonic sound of water displacing air, and
then, slowly and grudgingly, it subsides, collapsing not forward but
downward, melting back into the flat plane of the bay. Air rushes into
his lungs. He has survived the cancellation. The boundary is still
fluid, but it obeys the rules, for now.

\sceneline

The ocean is subdued, not tamed. It waits, heavy and sullen, for the
final locks, and Rowan must finish the sequence, because to leave the
sentence incomplete is to leave a door unlatched in a gale. He summons
the Sæfangi, the Sea-Holding --- not a negotiation but a definition, an
anchor. His voice drops an octave; he abandons projection and adopts a
crushing flatness, mimicking the sound of the horizon line itself.
\emph{``Með sæfanga\ldots{}''} He draws the circle in the air with his
right hand, the movement stiff but exact. \emph{``\ldots nöfnin lifa.''}
\emph{Through sea-holding, names remain.} He engages the continuous
present, the state that does not decay, riveting the names of the
village --- house, stone, path, self --- to their coordinates. The water
undergoes a change of phase: the restless chop and the interfering waves
and the chaotic texture all vanish, and in a perfect semicircle fifty
meters out the Atlantic becomes grey glass, unnaturally flat, reflecting
nothing, a zone of stillness made purely out of the density of a word.
The silence deepens as the wind slides over the flattened water without
friction enough to make a sound. Rowan sways; the energy required to
enforce stillness is greater than the energy required to move; his knees
buckle, the cartilage grinding, and he locks them, because he is the
bollard now, the stanchion, the fixed point.

One more. The seal. \emph{``Hrimfall.''} \emph{Winter-Falling.}
\emph{``Í vetrarfalli\ldots{}''} He looks at the grey glass and speaks
to the temperature of the universe. \emph{``\ldots tíminn sefur.''}
\emph{In winter-falling, time sleeps.} The raven ending, the one that
flies and lands and does not move again. The air snaps; the moisture in
it freezes instantly; frost blooms along the water's edge, the grammar
overriding the thermometer, white crystals racing across the wet basalt
and locking the liquid margin into a solid lattice, a rim of ice forming
on the lip of the sea, a fragile white border marking the exact limit of
the ocean's permission. It is done. The sentence is closed. The full
stop is set.

Rowan exhales, and the sound is a ragged tear in the quiet. His chest
heaves against the wool. He lowers his hands, which hang numb and
curled, and he stands in the silence of his own making, the boundary
holding, the village behind him safe for one more rotation of the earth,
the sea agreed to its terms. He turns to go, the movement slow, the turn
of a heavy thing on a worn pivot, meaning to walk back to the cottage
and the kettle and the chair where he will sit and wait for the light to
fail so that he can sleep and do this again.

\sceneline

But this morning, before he goes, he does the thing he has not done in
all the years he has stood on this rock. He takes the envelope from
against his chest, where it has ridden through every storm, softened now
almost to cloth, the ink of \emph{for E.B. only} faded to a watermark.
He is old enough now. He can hold the boundary alone, or nearly --- he
can hold it with a net of other throats, which Sigrún taught him is the
same as holding it well --- and so, by Edda's measure and his own, he is
finally finished enough to read it. He slides his thumb under the flap.
The old glue gives without resistance, as if it had been waiting to be
let go.

Inside is a single sheet, folded once, in his mother's hand --- not the
cramped English hand he knew, but an older hand, rounder, the hand of
the girl on the black coast, a hand he has never seen and recognizes
anyway. It is short. It is not an explanation; he understands, reading
it, that his mother was constitutionally incapable of an explanation,
that she died as she lived, in the locative, placing things near the
truth without ever letting the truth take a subject. He does not read it
aloud. Some of it is only for the woman it was addressed to, who kept it
against her own skin for a season and gave it back unopened, and he
honors that. But two lines he lets his eyes rest on for a long time,
there in the wind, and they are these: \emph{I could not stand where you
are standing. Forgive me the leaving; I have paid for it every day in a
country with no word for this coast.} And below, smaller, the way a
person writes the thing they most mean and least want to be caught
meaning: \emph{If it ever comes back to you through me, hold it more
gently than I could.}

He does not weep, because he has forgotten how, the way one forgets a
case one has stopped declining. But something in his chest, a knot he
has carried so long he took it for the shape of his own spine, eases by
a single degree, and he stands with the letter open in his cracked hands
and understands, at last and completely, that he did not come to
Brynjavík to study a dying grammar. He came to finish a sentence his
mother walked out of forty years before he was born into the middle of
it. His whole life --- the ear for what is withheld, the tools built to
translate a silence, the career he told his department was an accident
of interest --- all of it was the debt taking the only route home it
had. He is not the one who inherited the boundary despite the women who
were spent for it. He is the one his mother made, in advance, out of her
own refusal, to carry back what she could not. The inheritance was never
mechanical. It was the most personal thing that ever happened to him,
and it happened before he could speak.

He folds the letter along its old crease and puts it back against his
chest. He will not carry it much longer; he can feel that. But he will
carry it the rest of the way.

\sceneline

Then he hears it --- a sound on the wind, thin and high and distinct.
Not a gull; gulls scream, and this articulates. He freezes, his hand
finding the stone, and listens with nerves dulled by years of surf.
Again: a whisper, a syllable. He scans the horizon, and the line between
sea and sky is empty, no boat, no wreckage, and the moor to his left is
fog and green nothing.

He looks down at the patch of sand that gathers in the lee of the
platform, a small sheltered drift of shell and grit. There is a
footprint in it. Fresh, the edges sharp, the sand compressed and dark, a
heavy-soled boot, the heel dug deep, and it points toward the water. He
looks at his own boots, planted on the stone. He has not stepped in the
sand today, nor for weeks.

He looks back at the print. And beside it, half a stride off, where he
had not looked before, there is another --- smaller, much smaller, the
tread a shallow waffle of the kind stamped into the sole of a child's
rubber boot, the sort sold in every hardware shop in the north in four
sizes and three colours.

He knows the tread. He has seen it in the mud outside the hall on the
fourth Saturday of every month since its owner was old enough to walk
the boundary, and he saw it most recently at the north gatepost, where
she stood with her grandfather's hand over her hand on the cold stone
and got the dative wrong twice and right once and was told she was
clever, and went pink, and pretended she had not heard.

Sóley Bjarnadóttir. Eleven years old.

The two prints face the same grey water. Neither leads out of the
village and neither leads into the sea. They simply exist, two
impressions of weight arriving from nowhere, side by side, as if two
people had stood here in the dark and considered the work and then gone
back to sleep to be ready for it.

He does not go and ask her. He understands, standing over the sand with
his hands hanging at his sides, that if he went to the house and knocked
and asked her whether she had come down to the rock in the night she
would look at him with her father's flat grey eyes and say no, and would
not be lying.

Rowan looks at the sea. The grey glass holds still, but beneath it, deep
in the dark, he feels attention. The sea is watching him look at the
prints. It is patient. It knows he is finite, that his grammar is made
of meat and bone and that meat and bone fail. But the grammar itself ---
the rules, the cases, the need to speak --- endures, and it will find
throats to live in, and it does not much care what shape those throats
come in, or whose, or how old. When his voice finally cracks, when the
Aflending fails at last on his tongue, the station will not be
abandoned. The sea ensures its own audience.

And the prints in the sand, he understands, are not a threat. They are a
roster.

He told himself on this rock, the morning he took it up, that he would
carry the knowledge of who had paid for the boundary before him awake,
and would not let it become one more thing held at a distance and
re-labelled. He has kept that. Through Sigrún, through Arnar, through
every winter since, and he has been quietly proud of it in the way a man
is proud of a discipline nobody asked him for.

He looks at the small print in the sand and understands that keeping it
awake is not the same as being able to do anything about it. That was
the last consolation and it was never a consolation, only a better grade
of looking. The village will not choose her. Nobody will choose her.
There will simply come a morning, eight or twelve or twenty years from
now, when his voice fails on the third case and hers does not, and the
thing in the bay will note which throat finished the sentence, and the
matter will have been settled without one person having decided
anything.

Her father was twenty when the water got a handle on him. He has spoken
perhaps four hundred words in the years since, and he has never in that
time said his daughter's name where the bay could hear it, and one of
the four hundred --- patiently, twice, at a gatepost, on a Saturday, in
the rain --- was the dative.

``Someone else,'' Rowan says to the cold air, and then, correcting
himself, more quietly, in the plural the language keeps for the things
that must be shared to survive: \emph{``Einhver önnur. Fleiri en
einn.''} \emph{Some other. More than one.} He turns his back on the two
footprints and begins the long walk to the core of the village, leaving
the marks in the sand to be read by the rising tide, and carried, and
passed on.

\clearpage
\thispagestyle{empty}
\vspace*{2.4in}
\begin{center}
{\scshape\small About the Author}
\end{center}
\vspace{1.5em}
\noindent Elian Voigt writes literary and speculative fiction concerned with language, memory, and the labor of stewardship. \textit{Declensions of Dark Water} is the third novel in a sequence of works set along the threshold where grammar meets the physical world.

\vfill
\begin{center}
{\scshape\small Colophon}

\vspace{0.8em}
\ebgaramond{\large ❦\par}
\vspace{0.8em}
{\small This book was set in EB Garamond, a contemporary revival of Claude Garamont's sixteenth-century romans, designed by Georg Duffner and Octavio Pardo. The display ornament is the Aldine leaf, a printer's mark dating to fifteenth-century Venice.\par}
\vspace{1.5em}
{\ebgaramond\addfontfeature{Letters=SmallCaps}\small Form\ae trix\par}
\end{center}
\vspace*{0.6in}

\end{document}
